%!TEX program = lualatex
%%%%%%%%%%%%%%%%%%%%%%%%%%%%%%%%%%%%%%%%%%%%%%%%%%%%%%%%%%%%%%%%%%%%%%%%%%%%%%%
%% A Reference Grammar of the Iridian Language                               %%
%% 2026 Rewrite                                                               %%
%% Copyright (C) 2019--2026 by Roel Christian Yambao                         %%
%% This work may be distributed and/or modified under the conditions of the  %%
%% MIT License.                                                               %%
%%%%%%%%%%%%%%%%%%%%%%%%%%%%%%%%%%%%%%%%%%%%%%%%%%%%%%%%%%%%%%%%%%%%%%%%%%%%%%%

\documentclass[12pt]{amsbook}

\usepackage[
	papersize={6in,9in},
	inner=0.9in,
	outer=0.75in,
	top=0.85in,
	bottom=0.9in
]{geometry}

\usepackage{fontspec}
\setmainfont{EB Garamond}[
	RawFeature={+liga,+clig,+dlig},
	Numbers={OldStyle,Monospaced}
]

\usepackage{microtype}
\usepackage{booktabs}
\usepackage{tabularx}
\usepackage{ragged2e}
\usepackage{multicol}
\usepackage{enumitem}
\usepackage{hyperref}

%% Convenience column types for tabularx
%% L = left-aligned flex column (wrapping)
%% C = centred flex column
\newcolumntype{L}{>{\RaggedRight\arraybackslash}X}
\newcolumntype{C}{>{\Centering\arraybackslash}X}

%% Interlinear glosses (expex)
\usepackage{expex}
\lingset{everygla=\it, everyglb=\small, everyglft=, aboveglftskip=0pt}

%% Leipzig glossing abbreviations
\usepackage[glosses]{leipzig}

%% Custom glossing abbreviations for the 2026 Iridian grammar
%% Core case / alignment (Iridian grammatical tradition)
\newleipzig{Dir}{dir}{direct case}
\newleipzig{Trs}{trs}{transitive case}
\newleipzig{Ind}{ind}{indirect case}
\newleipzig{Voc}{voc}{vocative}

%% Aspect
\newleipzig{Pf}{pf}{perfective}
\newleipzig{Ret}{ret}{retrospective}
\newleipzig{Cont}{cont}{continuous}
\newleipzig{Prog}{prog}{progressive}
\newleipzig{Ctp}{ctp}{contemplative}

%% Valency operations
\newleipzig{Antip}{antip}{antipassive}
\newleipzig{Appl}{appl}{applicative}
\newleipzig{Caus}{caus}{causative}

%% Mood
\newleipzig{Imp}{imp}{imperative}
\newleipzig{Proh}{proh}{prohibitive}
\newleipzig{Hort}{hort}{hortative}
\newleipzig{Irr}{irr}{irrealis}

%% Non-finite / converb
\newleipzig{Seq}{seq}{sequential converb}
\newleipzig{Sim}{sim}{simultaneous converb}
\newleipzig{Purp}{purp}{purposive converb}
\newleipzig{Cncs}{cncs}{concessive converb}
\newleipzig{Cntr}{cntr}{contrastive converb}
\newleipzig{Cond}{cond}{conditional converb}
\newleipzig{Cv}{cv}{converb (generic)}

%% Nominal / discourse
\newleipzig{Neg}{neg}{negation}
\newleipzig{Def}{def}{definiteness marker}
\newleipzig{Spec}{spec}{specificity marker}
\newleipzig{Pl}{pl}{plural}
\newleipzig{Dem}{dem}{demonstrative}
\newleipzig{Prox}{prox}{proximal}
\newleipzig{Med}{med}{medial}
\newleipzig{Dist}{dist}{distal}
\newleipzig{Anim}{anim}{animate}
\newleipzig{Inan}{inan}{inanimate}
\newleipzig{Refl}{refl}{reflexive}
\newleipzig{Quot}{quot}{quotative}
\newleipzig{Exst}{exst}{existential}
\newleipzig{Cop}{cop}{copula}
\newleipzig{NegExst}{neg.exst}{negative existential}
\newleipzig{Nmlz}{nmlz}{nominalizer}

%% Custom commands (matching 2020 archive conventions)
\newcommand{\ird}[1]{\textit{#1}}
\newcommand{\foreign}[1]{\textit{#1}}
\newcommand{\rec}[1]{*\textit{#1}}
\newcommand{\trsl}[1]{`#1'}
\newcommand{\irdp}[2]{\ird{#1}, \trsl{#2}}
\newcommand{\orth}[1]{$\langle$#1$\rangle$}
\newcommand{\phon}[3]{\ird{#1} [#2] \trsl{#3}}
\newcommand{\mk}[1]{\textsc{#1}}
\newcommand{\sx}[1]{\textsuperscript{#1}}
\newcommand{\lingcontext}[1]{{\small\itshape [#1]}}

%% Abbreviations used in morpheme templates
\newcommand{\Stem}{\textsc{stem}}
\newcommand{\Asp}{\textsc{asp.align}}

\title{A Reference Grammar of the Iridian Language}
\author{Roel Christian Yambao}
\date{2026}

\begin{document}

\maketitle
\tableofcontents

%% Chapter 1: Overview
\include{chapters/01-overview}

%% Chapter 2: Phonology
\include{chapters/02-phonology}

%% Chapter 3: The Verb
\chapter{Verbs}\label{ch:verbs}

\section{Introduction}

The verb is the morphological nucleus of the Iridian clause, and upon it falls the principal burden of grammatical expression. While nouns admit of inflection for case alone and are in all other respects morphologically inert, the verb concentrates within itself all information pertaining to the nature of the event described, the relation of its participants, and the speaker's orientation toward what is asserted. Every finite verb bears the aspect-alignment suffix, which is the sole obligatory element of the verbal complex; beyond this the verb may carry markers of voice, mood, and evidentiality, in an order fixed by the morphological template set out in the following section. Tense, gender, number agreement, and pronominal cross-reference find no expression on the Iridian verb.

\section{The morphological template}\label{sec:morph-template}

The structure of a finite verb in Iridian is given by the following template, wherein parentheses enclose optional elements and the sole obligatory position is that of the aspect-alignment suffix:

\medskip
\begin{center}
(\Neg{}) + \Stem{} + (\Antip{} / \Appl{}) + (\Caus{}) + \Asp{}.\textsc{align} + (\textsc{mood}) + (\textsc{evid})
\end{center}
\medskip

The \textsc{asp.align} slot fuses aspect and alignment into a single morpheme and is the one element no finite verb may lack. The mood slot may be analysed as containing a null morpheme in the indicative. The negative prefix \ird{zá-} stands before the entire stem complex. Evidential suffixes are attested in the literature but remain insufficiently described; they are accordingly omitted from the present chapter. Person is not morphologically expressed on the Iridian verb: there is no pronominal clitic series and no agreement layer. Overt person reference is conveyed by full pronoun DPs, demonstratives in the third person, and zero anaphora where discourse recoverability licenses argument omission.


\subsection{The Mirror Principle and functional projections}\label{sec:vm-overview}

\noindent Baker's (1985) \term{Mirror Principle}
predicts that morphological derivations directly reflect syntactic
derivations: each morpheme corresponds to a syntactic head, and the
order of morphemes on the verbal complex mirrors the order in which
those heads are merged. Reading outward from the root, the Iridian
template projects the following hierarchy of functional
projections:

\begin{table}[ht]
  \footnotesize\sffamily
  \caption{Morphological template positions and corresponding
  syntactic projections (innermost to outermost).}\label{tab:vm-projections}
  \medskip
  \begin{tabularx}{\textwidth}{l l L l}
    \toprule
    {\sc position} & {\sc morpheme} & {\sc syntactic projection} &
    {\sc see} \\
    \midrule
    Root         & \Stem{}           & $\surd$P / VP (lexical predicate)
                 & §\,\ref{sec:vm-root} \\
    Voice        & \textsc{trs} / \Antip{} / \Appl{} & VoiceP / ApplP
                 & §\,\ref{sec:vm-antipassive}, §\,\ref{sec:vm-n} \\
    Causative    & \Caus{}           & CausP
                 & §\,\ref{sec:vm-caus} \\
    Asp.Align    & \Asp{}.\textsc{align} & AspP (fused Asp + Case)
                 & §\,\ref{sec:vm-asp} \\
    Neg (prefix) & \Neg{}            & NegP (see §\,\ref{sec:vm-neg})
                 & §\,\ref{sec:vm-neg} \\
    Mood         & \textsc{mood}     & MoodP / IP
                 & §\,\ref{sec:vm-moodp} \\
    Evidential   & \textsc{evid}     & EvidP$_\text{inner}$
                 & §\,\ref{sec:vm-evid} \\
    \bottomrule
  \end{tabularx}
\end{table}

\noindent The prefix position of \Neg{} (§\,\ref{sec:vm-neg})
presents a principled complication: the Mirror Principle predicts
that NegP, being scopally above AspP, should correspond to a suffix
that appears \emph{outside} the Asp.Align morpheme. The Iridian
prefix instead appears before the root. This is a well-known
typological type --- prevalent in verb-final languages --- and its
theoretical resolution is discussed in §\,\ref{sec:vm-tension-neg}.



\subsubsection{The root and VP}
\label{sec:vm-root}

The lexical stem is the exponent of the Root merged with the
functional category V (or, in Distributed Morphology terms, the
root $\surd$ inserted under a categorizing head $v_\text{intr}$ for
intransitives, or $v^*$ for transitives). The Iridian verb allows no
citation form; the stem is bound and always requires at least the
Asp.Align suffix. This confirms that the root must be licensed by a
selecting functional head: a bare root without an Asp projection is
unlicensed in the language.

The VP formed by the root + its arguments constitutes the innermost
layer of the verbal clause. In a transitive clause the internal
argument (patient) is merged as the complement of V:

\begin{center}
\begin{forest}
  syntax tree
  [VP
    [{DP\\{\footnotesize patient}}]
    [V
      [{\textit{root}}]
    ]
  ]
\end{forest}
\end{center}

\noindent Because Iridian is rigorously head-final
(§\,\ref{sec:th-head-finality}), this complement surfaces to the
left of the verbal head at every embedded level. The preverbal
position of the object is therefore not derived by movement but is
the base-generated configuration of the head-final VP.



\section{The verb stem and thematic classes}

Iridian verbs admit no citation form. Lexicographers record the bare stem with a following hyphen to indicate its character as a bound form requiring at least one suffix; the verb rendered in English as \textit{to eat} is thus entered as \ird{piašt-}, where \ird{t} constitutes the \term{thematic consonant}.

The thematic consonant, being the final consonant of the stem, is phonologically active: it undergoes regular alternations when certain suffixes are adjoined. Two mutation systems must be distinguished. Some suffixes trigger \term{sibilant replacement}, a closed morphologized stem mutation whose chief exponent is the antipassive \ird{-aš-}. Others trigger ordinary \term{palatalization}, the productive softening caused by front-vocalic suffixes such as \ird{-ěv}, \ird{-ěc}, and the conjunctive linker \ird{-e}. The five thematic classes in Table~\ref{tab:verbclasses} govern sibilant replacement only.

\begin{table}[ht]
	\footnotesize\sffamily
	\caption{Thematic consonant classes.}\label{tab:verbclasses}
	\medskip
	\begin{tabularx}{\textwidth}{l L L L}
		\toprule
		{\sc class} & {\sc thematic consonants} & {\sc change before \Antip{}} & {\sc example} \\
		\midrule
		I    & -t, -k, -c, -č  & → -č                    & \ird{piašt-} → \ird{piašč-} \\
		II   & -s, -š          & → -š                    & \ird{tras-} → \ird{traš-} \\
		II-A & -l, -p          & add -š                  & \ird{dal-} → \ird{dalš-} \\
		III  & -d, -h, -j, -z, -ž & → -ž               & \ird{vad-} → \ird{váž-} \\
		III-A & all remaining  & add -ž                  & \ird{ščen-} → \ird{ščenž-} \\
		\bottomrule
	\end{tabularx}
\end{table}

Loan stems with surface \orth{g} are assigned analogically to the voiced velar/fricative series and do not require a native phoneme /g/. Palatalizing suffixes do \emph{not} follow Table~\ref{tab:verbclasses}; they invoke the broader front-vocalic mutation system described in Section~\ref{sec:front-palatalization} of Chapter~\ref{ch:phon}.

Outside the closed antipassive classes, verbs use the same general morphophonemic engine as nouns: suffix class selection, weak-\ird{e} deletion or retention, grade formation where licensed, soft-stem selection, and finally cluster and voicing cleanup. The difference is one of density rather than mechanism. Nouns expose that engine most heavily in their case paradigms, while verbs deploy it more selectively in the irrealis, the simultaneous converb, the conjunctive linker, and certain nominalizing formations.

\section{Aspect and alignment}
\label{sec:aspect}

Iridian grammaticalizes \term{aspect} rather than \term{tense} on the finite
verb. There is no dedicated past-, present-, or future-tense suffix; temporal
location is instead inferred from discourse context, temporal adverbials, and
the lexical semantics of the predicate. The finite verb distinguishes five
aspectual categories: the {perfective}, the {retrospective}, the {continuous},
the {progressive}, and the {contemplative}. Aspect and alignment are fused in a
single suffix, and the aspect selected on the verb determines the clause's
alignment frame: the perfective and retrospective take ergative-absolutive
alignment, whereas the continuous, progressive, and contemplative take
nominative-accusative alignment.

Three temporal notions must be kept distinct. Tense locates a situation in time
relative to a reference point, most commonly the moment of speaking. Grammatical
aspect does not locate the situation in this way, but construes its internal
temporal contour: whether it is viewed as bounded, ongoing, prior from a later
vantage point, or not yet realized. \term{Lexical aspect} or \term{Aktionsart},
by contrast, belongs to the predicate itself. A predicate may denote a state,
activity, accomplishment, or achievement independently of the inflection it
bears. Iridian aspectual morphology therefore does not determine the inherent
temporal type of the eventuality; rather, it imposes a viewpoint upon an
eventuality whose lexical contour is supplied by the verb phrase.

A useful descriptive frame is Reichenbach's distinction among speech time
(\textit{S}), reference time (\textit{R}), and event time (\textit{E}):
\textit{S} is the moment of utterance, \textit{E} the time at which the
situation itself takes place, and \textit{R} the temporal vantage point from
which that situation is presented. The usefulness of \textit{R} is that speakers
do not always describe an event directly from the moment of speaking, but may
instead locate it from some later or earlier point already established in the
discourse; English \textit{had left}, for example, places the leaving
(\textit{E}) before a reference point (\textit{R}) that is itself prior to the
speech moment (\textit{S}). In a tense-prominent system, finite morphology often
fixes the relation among all three. In Iridian, by contrast, the finite suffix
primarily construes the relation between \textit{E} and \textit{R}, while the
relation of \textit{R} to \textit{S} is often left to context, temporal
adverbials, and discourse sequencing. This is why the same aspectual category
may receive past-, present-, or future-oriented readings in different
environments without ceasing to be aspectual rather than temporal.

The five categories may be sketched in these terms. The \term{perfective}
presents an eventuality as bounded and viewed as a whole, with \textit{R}
aligned to the event in its entirety. The \term{retrospective} places \textit{E}
before \textit{R}, yielding a later vantage point from which present relevance,
result, prior experience, or other discourse effects of anteriority may be
construed. The \term{continuous} and \term{progressive} both place \textit{R}
within \textit{E}, but they differ in the kind of internal view imposed: the
continuous presents the situation as durative or internally extended, whereas
the progressive more narrowly profiles a dynamic event as actively unfolding at
the reference point. The \term{contemplative} places \textit{E} after
\textit{R}, presenting the event as projected, intended, expected, or otherwise
not yet realized. English translation equivalents therefore vary by context:
none of these categories should be reduced to simple labels such as ``past,''
``present,'' or ``future.'' Moreover, because viewpoint aspect interacts with
lexical aspect, the same inflection does not yield identical interpretive
effects with every predicate class. With states and activities, the continuous
naturally yields persistent or ongoing readings; with accomplishments it
typically selects an internal phase rather than a completed endpoint; with
achievements it may invite iterative, preparatory, or inceptive coercions.
Conversely, the perfective and retrospective readily package telic predicates as
bounded wholes or prior results, but with atelic predicates they instead delimit
a stretch, occasion, or state viewed from an external vantage point. The exact
semantic and distributional consequences of these interactions are discussed in
the subsections that follow.

Typologically, this pattern is an instance of aspect-conditioned split
ergativity. Iridian does not maintain a single alignment pattern across all
finite clauses; rather, the aspectual category selected on the verb determines
whether the clause is built on an ergative-absolutive or a nominative-accusative
frame. Aspect in Iridian therefore has consequences at two levels at once: it
construes the temporal profile of the eventuality and it determines the
morphosyntactic alignment through which the clause is organized.

\begin{figure}[t]
\centering
\begin{tikzpicture}[
  x=1.05cm,y=1cm,
  evt/.style={line width=1.1pt},
  refpt/.style={circle,fill,inner sep=1.3pt},
  speech/.style={diamond,draw,inner sep=1.2pt},
  rowlab/.style={anchor=east,font=\small},
  every node/.style={font=\small}
]

% All rows hold S at x = 4.5 so the deictic baseline is visually constant.

% perfective (y = 0)
\node[rowlab] at (0,0) {perfective};
\draw[->] (0.5,0) -- (8.5,0);
\draw[evt] (3.0,0) -- (6.0,0);
\draw (3.0,0.15) -- (3.0,-0.15);
\draw (6.0,0.15) -- (6.0,-0.15);
\node[refpt] at (3.0,0) {};
\node[refpt] at (6.0,0) {};
% R bracket above the timeline, E label between bracket and line
\node at (4.5,0.30) {$E$};
\draw[decorate,decoration={brace,amplitude=3pt}]
  (3.0,0.55) -- (6.0,0.55)
  node[midway,above=4pt] {$R$};
% S below, matching spacing of other rows
\draw[dashed] (4.5,-0.08) -- (4.5,-0.42);
\node[speech] (s1) at (4.5,-0.55) {};
\node[below=2pt of s1] {$S$};

% retrospective (y = -2.0)
\node[rowlab] at (0,-2.0) {retrospective};
\draw[->] (0.5,-2.0) -- (8.5,-2.0);
\draw[evt] (1.8,-2.0) -- (3.4,-2.0);
\draw (1.8,-1.85) -- (1.8,-2.15);
\draw (3.4,-1.85) -- (3.4,-2.15);
\node at (2.6,-1.65) {$E$};
\node[refpt,label=above:{$R$}] at (4.5,-2.0) {};
\node[speech] (s2) at (4.5,-2.55) {};
\node[below=2pt of s2] {$S$};
\draw[dashed] (4.5,-2.08) -- (4.5,-2.42);

% continuous (y = -4.0)
\node[rowlab] at (0,-4.0) {continuous};
\draw[->] (0.5,-4.0) -- (8.5,-4.0);
\draw[dotted] (1.6,-4.0) -- (2.2,-4.0);
\draw[evt] (2.2,-4.0) -- (6.8,-4.0);
\draw[dotted] (6.8,-4.0) -- (7.4,-4.0);
\node at (3.0,-3.65) {$E$};
\node[refpt,label=above:{$R$}] at (4.5,-4.0) {};
\node[speech] (s3) at (4.5,-4.55) {};
\node[below=2pt of s3] {$S$};
\draw[dashed] (4.5,-4.08) -- (4.5,-4.42);

% progressive (y = -6.2)
\node[rowlab] at (0,-6.2) {progressive};
\draw[->] (0.5,-6.2) -- (8.5,-6.2);
\draw[evt] (2.4,-6.2) -- (7.0,-6.2);
\draw (2.4,-6.05) -- (2.4,-6.35);
\draw (7.0,-6.05) -- (7.0,-6.35);
\node at (6.5,-5.85) {$E$};
\draw[->] (2.8,-5.55) -- (4.2,-5.55)
  node[midway,above,font=\scriptsize] {dynamic phase};
\node[refpt,label=above:{$R$}] at (4.5,-6.2) {};
\node[speech] (s4) at (4.5,-6.75) {};
\node[below=2pt of s4] {$S$};
\draw[dashed] (4.5,-6.28) -- (4.5,-6.62);

% contemplative (y = -8.2)
\node[rowlab] at (0,-8.2) {contemplative};
\draw[->] (0.5,-8.2) -- (8.5,-8.2);
\node[refpt,label=above:{$R$}] at (4.5,-8.2) {};
\draw[evt] (5.8,-8.2) -- (7.5,-8.2);
\draw (5.8,-8.05) -- (5.8,-8.35);
\draw (7.5,-8.05) -- (7.5,-8.35);
\node at (6.65,-7.85) {$E$};
\node[speech] (s5) at (4.5,-8.75) {};
\node[below=2pt of s5] {$S$};
\draw[dashed] (4.5,-8.28) -- (4.5,-8.62);

\end{tikzpicture}
\caption{Reichenbachian schematic for Iridian aspect. $E$ = event time, $R$ =
reference time, $S$ = speech time. In the perfective, $R$ spans the bounded
event (marked by dots at the endpoints and a bracket); elsewhere, $R$ is a
point-like vantage. In this figure, dashed connectors mark the unmarked deictic baseline in which
$R$ coincides with $S$. $S$~is held at a constant horizontal position across all
rows; in discourse, adverbials or sequencing may shift $R$ independently of $S$.}
\label{fig:iridian-aspect-reichenbach}
\end{figure}

\subsection{Alignment in nominative aspects (continuous, progressive, contemplative)}

In the continuous, progressive, and contemplative aspects the main argument of
the verb stands in the direct case regardless of transitivity, while the patient
of a transitive verb appears in the indirect. This is a nominative-accusative
pattern, the direct case serving as the absolutive and the indirect as the
oblique.

\pex
\a
\begingl
\gla Marek zmacime.//
\glb Marek.\Dir{} run-\Prog{}//
\glft \trsl{Marek is running.} \lingcontext{intransitive}//
\endgl
\a
\begingl
\gla Marek travě piastime.//
\glb Marek.\Dir{} bread.\Ind{} eat-\Prog{}//
\glft \trsl{Marek is eating bread.} \lingcontext{transitive; P = \Ind{}}//
\endgl
\a
\begingl
\gla Vytam shoda.//
\glb shirt.\Dir{} wear-\Cont{}//
\glft \trsl{(I) am wearing a shirt.} \lingcontext{stative continuous}//
\endgl
\a
\begingl
\gla Vytam shodime.//
\glb shirt.\Dir{} wear-\Prog{}//
\glft \trsl{(I) am putting on a shirt.} \lingcontext{dynamic progressive}//
\endgl
\xe

\subsection{Alignment in ergative aspects (perfective and retrospective)}

In the perfective and retrospective aspects the patient of a transitive verb is placed in the direct case, and the agent takes the transitive case. This is an ergative-absolutive pattern; what Iridian grammar calls the \emph{transitive case} is, in typological terms, the ergative, and the \emph{direct case} serves as the absolutive. For a transitive verb to trigger ergative alignment, it must carry the Voice head \ird{-n-} immediately before the aspect suffix.

\begin{table}[ht]
	\footnotesize\sffamily
	\caption{Aspect-alignment suffixes.}\label{tab:aspect}
	\medskip
	\begin{tabularx}{0.85\textwidth}{L l L}
		\toprule
		{\sc aspect} & {\sc suffix} & {\sc notes} \\
		\midrule
		Perfective              & \ird{-ek} / \ird{-ik}  & \ird{-ik} when stem vowel is \ird{ě} or \ird{e} (dissimilation; see §\ref{sec:aspect-pf}); \ird{-nik} after transitive \ird{-n-} \\
		Retrospective           & \ird{-ac} & \\
		Continuous              & \ird{-a}   & \\
		Progressive             & \ird{-ime} & \\
		Contemplative           & \ird{-ach} & \\
		\bottomrule
	\end{tabularx}
\end{table}

Three configurations are available in the ergative aspects. When the Voice head \ird{-n-} is present, the clause takes full ergative alignment: the patient is licensed as a core argument and appears in the direct case, and the agent takes the transitive. When \ird{-n-} is absent and the verb bears the antipassive \ird{-aš-}, the agent appears in the direct case and the patient is removed from core argument structure. When neither is present, the bare perfective suffix is used, a form grammatical only with inherently intransitive verbs, for which the question of alignment does not arise. For a transitive verb the bare perfective is inadmissible, because the patient lacks the overt patient-licensing Voice head required for core realization.

\pex
\a
\begingl
\gla Travat Marku piaštnik.//
\glb bread.\Dir{}.\Def{} Marek.\Trs{} eat-\textsc{trs}-\Pf{}//
\glft \trsl{Marek ate the bread.} \lingcontext{ergative; P = \Dir{}}//
\endgl
\a
\begingl
\gla Marek piašček.//
\glb Marek.\Dir{} eat-\Antip{}-\Pf{}//
\glft \trsl{Marek ate.} \lingcontext{antipassive; A = \Dir{}}//
\endgl
\xe

The choice between these constructions is governed primarily by the discourse status of the patient. The ergative construction foregrounds the patient: it is the construction of a specific, discourse-trackable participant who is either newly introduced or reactivated, and the event is asserted as bringing about a change of state in that participant. It is accordingly characteristic of narrative foreground, where the consequences of events for particular persons or things constitute the substance of what is related. Patient definiteness correlates strongly with the ergative, and the definite marker \ird{-t} most commonly appears on the patient noun in this construction, though it is not formally required.

The antipassive construction, by contrast, foregrounds the agent's activity without reference to any outcome in a patient. It is the construction of choice when the patient is non-referential, semantically incorporated, or of no discursive moment: in habitual and generic statements, in activity-focus contexts wherein the event itself and not its effect upon some participant is what is asserted, and wherever the clause contains no patient at all. The guiding principle is that the patient, where it exists, must either be present through the ergative construction or wholly absent; the antipassive admits of no intermediate state in which a patient is demoted to an oblique position.

\subsection{The perfective}\label{sec:aspect-pf}

The perfective encodes a completed action viewed as a whole at a specific point in time, and is the primary narrative aspect.

The bare intransitive perfective suffix presents two allomorphs: \ird{-ek} is the default form, and \ird{-ik} appears where the stem's final syllable has \ird{ě} or \ird{e} as its nucleus. The raising is dissimilatory in character: adjoining \ird{-ek} directly to such a stem would produce a CěCe or CeCe sequence across a single intervening consonant, which the language avoids. The allomorphy is thus conditioned by the stem vowel and not by the thematic consonant class. Where a valency morpheme --- the antipassive \ird{-aš-}, the applicative \ird{-éb-}, or the causative \ird{-na-} --- intervenes between stem and suffix, it contributes its own vowel and breaks the sequence; raising accordingly does not apply in derived forms. The ergative perfective \ird{-nik}, which always follows the Voice head \ird{-n-}, is a fixed form and does not derive from this rule. % Cross-reference: this dissimilation should be described in the phonology chapter alongside the other vowel alternations.

When negated, the perfective carries a specific implicature: the negated perfective describes something that was expected to occur but did not. For the neutral negation of a past event, free from this implication, the retrospective is preferred.

\pex
\a
\begingl
\gla Marek zapiaščik.//
\glb Marek.\Dir{} \Neg{}-eat-\Antip{}-\Pf{}//
\glft \trsl{Marek didn't eat (but he should have).}//
\endgl
\a
\begingl
\gla Marek zapiaščac.//
\glb Marek.\Dir{} \Neg{}-eat-\Antip{}-\Ret{}//
\glft \trsl{Marek didn't eat (neutral past).}//
\endgl
\xe

\subsection{The retrospective}

The retrospective (\ird{-ac}) describes a past action with continuing present relevance, or an action that was non-volitional or accidental. It corresponds in certain uses to the English perfect, in others to a non-volitional past.

\pex
\a
\begingl
\gla Pražě de stožac.//
\glb Prague.\Ind{} to go-\Ret{}//
\glft \trsl{(I) have (been to) Prague before.}//
\endgl
\a Compare:
\begingl
\gla Bych Pražě de stožek.//
\glb yesterday Prague.\Ind{} to go-\Pf{}//
\glft \trsl{Yesterday I went to Prague.} \lingcontext{specific past event; perfective}//
\endgl
\xe

The ergative and antipassive constructions encode distinct kinds of continuing relevance in the retrospective. The transitive retrospective (\ird{-n-ac}) situates the relevance in the outcome state of the patient: something has been affected, and that effect persists or bears upon the present situation. The antipassive retrospective (\ird{-aš-ac}) locates the relevance in the experience of the agent: the event is part of the agent's history, not a change in some external participant. This distinction parallels the English opposition between \textit{the bread has been eaten} --- outcome-relevant, a fact about the bread --- and \textit{I have eaten} --- experiential, a fact about the speaker.

\subsection{Ambitransitive verbs}\label{sec:ambitransitive}

A small class of verbs possesses two readings of different argument structure: an intransitive reading, in which no patient is structurally available, and a transitive reading, in which a patient is licensed. These are lexical ambitransitives; neither reading is derived from the other through the antipassive, applicative, or causative, but both are basic senses of the verb. Two examples are \ird{lov-} (\trsl{hunt, forage [intrans.]; catch, take [trans.]}) and \ird{kěl-} (\trsl{sing, perform [intrans.]; sing a song, chant a text [trans.]}).

The diagnostic for membership in this class is whether the Voice head \ird{-n-} is admissible in one reading but impossible in the other. For \ird{lov-}: \ird{lovnik} is grammatical in the transitive reading, as something was caught; no \ird{-n-} form is available in the intransitive reading, wherein the event is construed as the activity of hunting itself, with no quarry in the event frame. The bare perfective accesses only the intransitive reading; the antipassive and ergative access only the transitive.

\pex
\a
\begingl
\gla Marek bych lovek.//
\glb Marek.\Dir{} yesterday hunt-\Pf{}//
\glft \trsl{Marek went hunting yesterday.} \lingcontext{intransitive reading; bare \ird{-ek}}//
\endgl
\a
\begingl
\gla Zvěrat Marku lovnik.//
\glb game.\Dir{}.\Def{} Marek.\Trs{} hunt-\textsc{trs}-\Pf{}//
\glft \trsl{Marek caught the game.} \lingcontext{transitive reading; ergative}//
\endgl
\a
\begingl
\gla Marek bych lovžek.//
\glb Marek.\Dir{} yesterday hunt-\Antip{}-\Pf{}//
\glft \trsl{Marek was hunting yesterday.} \lingcontext{transitive reading; antipassive}//
\endgl
\xe

The two non-ergative forms, though they may coincide in surface expression in many contexts, are not interchangeable. \ird{Lovek} belongs to the intransitive reading: hunting is construed as a pure activity, with no quarry present even in the conceptual frame of the event. \ird{Lovžek} belongs to the transitive reading: the hunt is directed toward prey, and that prey --- though it may not be expressed in any case form, the antipassive having removed it from the clause's argument structure --- is none the less understood as part of the conceptual event. The distinction corresponds, in practice, to the difference between going out to hunt and being engaged in the act of hunting something.

For \ird{kěl-}, the stem vowel is \ird{ě}, so the bare intransitive perfective takes the \ird{-ik} allomorph by dissimilation: \ird{kělek} would create a CěCe sequence across the single intervening \ird{-l-}, which is avoided. The antipassive \ird{kělšek} is unaffected, the \ird{-š-} of the antipassive introducing its own vowel and breaking the sequence. The ergative \ird{kělnik} is the fixed ergative form; it and the bare \ird{kělik} are distinguished only by the presence or absence of \ird{-n-}:

\pex
\a \ird{kělik} \trsl{sang, performed} \lingcontext{intransitive; bare \ird{-ik}: no song in event frame}
\a \ird{kělnik} \trsl{sang it (a specific song)} \lingcontext{ergative; \ird{-n-ik}: song = \Dir{}}
\a \ird{kělšek} \trsl{was singing} \lingcontext{antipassive; \ird{-š-ek}: Class~II-A augment}
\xe

The minimal pair \ird{kělik} / \ird{kělnik} provides the clearest illustration in the grammar of the bare perfective opposed to the ergative, the two forms differing only by the inserted \ird{-n-}.

\subsection{The continuous}

The continuous (\ird{-a}) encodes an ongoing state or a habitual, generic action and is stative in character.

\ex
\begingl
\gla Praže to mulša.//
\glb Prague.\Ind{} \textsc{loc} live-\Cont{}//
\glft \trsl{(I) live in Prague.}//
\endgl
\xe

\subsection{The progressive}

The progressive (\ird{-ime}) encodes a dynamic activity in progress, distinguished from the continuous by its dynamic rather than stative character. Compare the stative \ird{Vytam shoda}, \textit{I am wearing a shirt} --- a state of affairs obtaining at the time of speaking --- with the progressive \ird{Vytam shodime}, \textit{I am putting on a shirt}, which describes an action actively in course.

\subsection{The contemplative}

The contemplative (\ird{-ach}) encodes a mental state directed toward a potential future action: intention, consideration, or disposition toward an act not yet performed. It may be rendered in English as \textit{about to}, \textit{considering}, or \textit{intending to}, according to context.


\subsection{Formal analysis: the Asp.Align morpheme}
\label{sec:vm-asp}

The aspect-alignment suffix is the sole obligatory morpheme in every
finite Iridian verbal form. Its double burden --- encoding both the
aspect of the event and the alignment pattern of case assignment ---
makes it the central theoretical challenge of the Iridian verbal
system.

\subsubsection{Aspect-driven case assignment}
\label{sec:vm-asp-case}

Iridian alignment is organized by the Asp.Align head, whose aspectual
values determine the clause's default absolutive orientation. The
alignment defaults are:

\begin{description}[leftmargin=2em, labelwidth=!, labelindent=0em]
\item[\upshape ERG-aspect domain (perfective, retrospective).]
  Asp.Align assigns \textsc{abs} (direct case) to the internal
  argument (patient) and \textsc{erg} (transitive case) to the
  external argument (agent). The ergative is the marked, dependent
  case: it is assigned to the agent only when a patient has already
  received absolutive (Marantz 1991). The default absolutive is
  therefore the \emph{patient}.
\item[\upshape NOM-aspect domain (continuous, progressive,
  contemplative).] Asp.Align assigns \textsc{abs} (direct case) to
  the external argument (agent) and \textsc{obl} (indirect case) to
  the internal argument (patient). The default absolutive is
  therefore the \emph{agent}.
\end{description}

\noindent These defaults constitute the unmarked shape of the
Iridian alignment system. They may be altered by overt Voice
morphology. When a transitive stem is merged with an overt Voice
head carrying a specification for absolutive orientation, that
specification takes precedence over the aspectual default and is
respected by Asp.Align in the realization of the clause. The
aspectual head therefore determines the default alignment pattern,
but an overt Voice head may impose a different absolutive selection
in marked constructions.

The morpheme \ird{-n-} is such a head: it licenses the internal
argument and specifies \textsc{[pat-abs]}, yielding ordinary
ergative transitives in the ERG domain and passive constructions in
the NOM domain (§\,\ref{sec:vm-n}, §\,\ref{sec:vm-passive}).
The antipassive \ird{-aš-} likewise constitutes an overt Voice head,
but it is subject to a structural input condition: it may apply only
where the patient already stands as the absolutive participant, and
it therefore surfaces exclusively in ERG-aspect clauses
(§\,\ref{sec:vm-antipassive}).

In the Minimalist framework, the Asp.Align head's default case
assignment is mediated by Agree: the head probes its c-command domain
for DPs bearing unvalued Case features and values them according to
its featural content (Chomsky 2001). When an overt Voice head is
present, its absolutive specification is satisfied before Asp.Align
probes, and the remaining non-absolutive argument receives the case
predicted by the aspectual domain. The fusion of aspect and alignment
into a single suffix is not an accident: it reflects the close
coupling in Iridian between the temporal framing of an event and the
information-structural decision of which participant is the clause's
absolutive pivot.

\subsubsection{Split ergativity and the information-structure interface}
\label{sec:vm-asp-info}

The grammar's description of aspect selection as a
\emph{pragmatic} decision (the speaker selects the aspect that
places the discourse-salient argument in absolutive) requires
formal reconciliation with the syntactic Agree mechanism described
above.

Two strategies are available within Minimalism:

\begin{description}[leftmargin=2em, labelwidth=!, labelindent=0em]
\item[\upshape Top-down feature percolation.] Topic assignment
  operates in the left periphery (Spec,TopP), where an argument
  already bears topic features. These features are passed down
  (through some mechanism of syntactic binding or coindexation) to
  the Asp head, constraining which Vocabulary Item is inserted at
  Asp. A patient-topic clause selects an ergative Asp morpheme; an
  agent-topic clause selects a nominative Asp morpheme. This
  requires either late lexical insertion (consistent with
  Distributed Morphology) or a bidirectional Agree relation between
  TopP and AspP.
\item[\upshape Late Vocabulary Insertion.] In a DM-style account
  (Halle \& Marantz 1993; Embick \& Noyer 2007), the Asp node in narrow syntax is
  underspecified for the choice of morpheme. Vocabulary Items are
  inserted post-syntactically by matching feature bundles. The
  speaker's information-structural decision (which argument is
  discourse-topic) is registered as a feature [\textsc{top}] on one
  of the DP arguments; this feature is present when Vocabulary
  Insertion takes place and conditions which Asp morpheme is
  inserted.
\end{description}

\noindent The DM approach has the advantage of keeping
narrow-syntax Agree simple (the Asp head probes for DPs in its
c-command domain without knowing in advance which DP will be the
absolutive) while assigning information-structural conditioning to
the PF component. This is the preferred analysis here, though
it requires elaboration of how [\textsc{top}] features are
represented in the DM morphological component.

The key theoretical desideratum: the grammar's informal claim that
``aspect choice is partly pragmatic'' must be formalized as a
specific interface mechanism. Until this is done, the claim risks
being merely redescriptive. A concrete proposal: the
[\textsc{top}] feature on a DP argument is set at the information
structure level (in the C-domain or in a dedicated information
structure component co-indexed with TopP), and this feature
determines which Asp Vocabulary Item is inserted. The split is
then not a quirk of case morphology but a reflex of the language's
deep commitment to topic-comment organization at every level of
the grammar.

\subsubsection{Mood as a competing head}
\label{sec:vm-mood-asp}

When a mood suffix is present (imperative \ird{-ím}, prohibitive
\ird{-é}, hortative \ird{-ká/-ku}, irrealis \ird{-ěv}), the
Asp.Align slot is suppressed entirely: mood replaces aspect. In
syntactic terms, this is most naturally analyzed as the MoodP head
blocking the Asp head: when MoodP is merged above AspP, the Asp
head is defective and receives no Vocabulary Item at the AspP node.
The morphological output is therefore Mood without Aspect.

This has a consequence for case assignment. In the directive moods,
the case-assigning features of the Asp.Align head are absent, and
case is instead determined by the principles of the directive domain
described in §\,\ref{sec:mood-directive}: the direct case is assigned
to the theme or the clause's sole participant, the transitive to an
overt speech-act participant subject where competition would otherwise
arise, and the indirect to applicative beneficiaries and causative
causees. The directive case geometry is therefore not a carry-over
of either the ergative or the nominative-accusative indicative
alignment, but the output of its own case-assignment principles once
the Asp.Align head is suppressed.


\section{Valency operations}

Three derivational operations may be applied to the verb stem before the aspect-alignment suffix: the antipassive, the applicative, and the causative. Where more than one is present, the order is fixed: antipassive or applicative precedes the causative.

\subsection{The antipassive (\ird{-aš-})}
\label{sec:vm-antipassive-desc}

The antipassive is an overt Voice head that applies only where the
patient is structurally the absolutive participant. It displaces that
patient from absolutive status, promotes the agent to absolutive, and
removes the patient from core argument structure. Because the
structural target it requires --- a patient standing as the
absolutive --- is available only in ERG-aspect clauses, the
antipassive is exclusively an ergative-domain construction. It is
therefore the obligatory non-ergative construction for any verb
capable of bearing the transitive/passive Voice head \ird{-n-} in
the perfective and retrospective aspects: such a verb must appear
with either \ird{-n-} (ergative transitive) or \ird{-aš-}
(antipassive). The bare perfective serves no antipassive function and
is reserved exclusively for inherently intransitive verbs, as
described in Section~\ref{sec:ambitransitive}.

The antipassive suffix is \ird{-aš-}, though its phonological realisation depends upon the thematic consonant class of the stem, following the alternations set out in Table~\ref{tab:verbclasses}. The stem-final consonant is replaced or augmented as shown; the underlying \ird{-aš-} string surfaces as the modified stem consonant followed by the relevant sibilant. Thus \ird{piašt-} yields \ird{piašč-} as its antipassive stem:

\pex
\a
\begingl
\gla Marek piašček.//
\glb Marek.\Dir{} eat-\Antip{}-\Pf{}//
\glft \trsl{Marek ate.}//
\endgl
\a
\begingl
\gla Marek houlam se piašček.//
\glb Marek.\Dir{} spoon.\Ind{} \textsc{inst} eat-\Antip{}-\Pf{}//
\glft \trsl{Marek ate with a spoon.} \lingcontext{instrument adjunct in \Ind{}}//
\endgl
\xe

The patient, having been removed from core argument structure by the antipassive, may not be expressed in any case form; it is accordingly inadmissible as an indirect-case argument, as may be seen from the following:

\ex
\begingl
\gla *Marek travě piašček.//
\glb Marek.\Dir{} bread.\Ind{} eat-\Antip{}-\Pf{}//
\glft \trsl{(intended: Marek ate bread.)}//
\endgl
\xe

\noindent Where a patient is to be expressed in any form, the ergative construction is required. Instrument and adjunct arguments, which bear no patient relation to the predicate, remain fully admissible in the indirect case, as in the spoon example above. The absence of an expressible patient is not merely a discourse choice to background a participant that is otherwise retained in syntax: under the antipassive, the patient is removed from the clause's core argument structure and therefore cannot surface in any case form. The construction does not merely withhold the patient from prominence; it eliminates the patient from the overt argument structure of the clause. For an inherently transitive verb the patient is typically still understood as part of the conceptual event, but what is asserted is the agent's activity; the patient has no realization as a syntactic argument.

Equally, for a transitive verb, the bare perfective without either \ird{-n-} or \ird{-aš-} is ungrammatical:

\ex
\begingl
\gla *Marek piaštek.//
\glb Marek.\Dir{} eat-\Pf{}//
\glft \trsl{(intended: Marek ate.)} \lingcontext{bare perfective illicit for transitive verb}//
\endgl
\xe

\noindent\textit{Note:} The antipassive and the applicative are in complementary distribution; a verb cannot take both simultaneously.


\subsubsection{Minimalist analysis: the antipassive as VoiceP}
\label{sec:vm-antipassive}

The antipassive \ird{-aš-} is an overt Voice head (Kratzer 1996;
Alexiadou et~al. 2006) subject to a structural input condition: it
may be merged only in a clause whose input to VoiceP contains a
patient that is eligible for absolutive realization. It then removes
that patient from core argument structure and promotes the agent to
the absolutive position vacated by that removal.

\paragraph{Input condition.}
The antipassive requires, as the structural context for its
application, a clause in which the patient stands in the absolutive
position made available by the clause's alignment architecture. In
the ERG-aspect domain, Asp.Align's default assigns absolutive to the
patient (§\,\ref{sec:vm-asp-case}); the input condition is therefore
satisfied, and the antipassive is licensed. In the NOM-aspect domain,
the absolutive is assigned by default to the agent; no patient stands
in the absolutive position, and the structural target required by the
antipassive is absent. The antipassive is consequently unlicensed in
NOM-aspect clauses. This is a hard grammatical restriction, not a
preference or economy effect.

\paragraph{Argument-structure removal.}
Under the antipassive, the patient is not demoted to an oblique slot
within the clause. It is removed from the clausal core entirely.
Oblique expression of the patient is therefore not merely
dispreferred but impossible: the construction does not retain a
suppressed patient that has been reassigned to a peripheral position.
The patient ceases to be part of the overt argument structure of the
clause. This is why an indirect-case patient is ungrammatical with
the antipassive (see §\,\ref{sec:vm-antipassive-desc} for examples),
while instrument and adjunct arguments in the indirect case remain
fully admissible.

The template order --- antipassive between root and Asp.Align ---
confirms that VoiceP immediately dominates VP:

\begin{center}
\begin{forest}
for tree={align=center}
[AspP
  [{DP\\{\footnotesize agent}\\{\footnotesize\textsc{[abs]}}}]
  [{Asp$'$}
    [{Asp\\{\footnotesize\textsc{[erg]}}}]
    [VoiceP
      [{DP\\{\footnotesize agent}}]
      [{Voice$'$}
        [{Voice\\\ird{-aš-}\\{\footnotesize\textsc{[agt-abs]}}\\{\scriptsize requires p-abs}}]
        [VP
          [{\textit{root}}]
        ]
      ]
    ]
  ]
]
\end{forest}
\end{center}

\noindent The label \textsc{[p-abs required]} on the antipassive Voice
head represents the input condition: the head selects for a VP whose
internal argument would stand as the absolutive participant in the
relevant alignment environment. In ERG aspects this condition is met;
in NOM aspects it is not. The complementary distribution of the
antipassive and the applicative (\ird{-aš-} and \ird{-éb-} cannot
co-occur on the same stem) is a direct consequence of both occupying
the same VoiceP projection: only one Voice head may be merged at this
tier.

In ERG-aspect transitives, the patient requires overt Voice licensing
to be realized as a core argument (§\,\ref{sec:vm-n}). A verb bearing
neither \ird{-aš-} nor \ird{-éb-} nor \ird{-n-} in an ergative
aspect therefore fails for a different reason: the patient lacks the
overt patient-licensing head required for core realization. That
exclusion is distinct from the structural input condition that blocks
the antipassive in NOM aspects; the two empty cells in the alignment
system are each excluded by a separate principle (§\,\ref{sec:vm-asp-case}).


\subsection{The applicative (\ird{-éb-})}

The applicative promotes a peripheral argument --- most commonly a beneficiary or recipient, but also a goal, location, or instrument --- to core argument status, placing it in the absolutive. The suffix \ird{-éb-} derives from an older benefactive morpheme and still defaults to a beneficiary or recipient reading when such an interpretation is available. Where the lexical semantics of the verb or the wider clause exclude a benefactive reading, the applicative is interpreted instead as promoting the most discourse-salient non-core participant, typically a goal, location, or instrument. Only one applicative argument may be promoted in a single clause.

The contrast is most readily observed in nominative-alignment contexts, where the applied participant becomes the direct-case object that would otherwise have remained oblique:

\pex
\a
\begingl
\gla Marek bylam piastime.//
\glb Marek.\Dir{} child.\Ind{} eat-\Prog{}//
\glft \trsl{Marek is eating for the child.} \lingcontext{beneficiary remains oblique}//
\endgl
\a
\begingl
\gla Marek bylam piaštébime.//
\glb Marek.\Dir{} child.\Dir{} eat-\Appl{}-\Prog{}//
\glft \trsl{Marek is feeding the child.} \lingcontext{beneficiary promoted to core status}//
\endgl
\xe

\noindent With verbs of motion and placement the same suffix promotes a goal or location rather than a beneficiary:

\pex
\a
\begingl
\gla Marek nužam stachime.//
\glb Marek.\Dir{} town.\Ind{} go-\Prog{}//
\glft \trsl{Marek is going toward the town.}//
\endgl
\a
\begingl
\gla Marek hrad stéběme.//
\glb Marek.\Dir{} town.\Dir{} go-\Appl{}-\Prog{}//
\glft \trsl{Marek is entering the town / making the town his destination.}//
\endgl
\xe

\noindent Instrumental promotion is likewise possible, though it is pragmatically more marked, being preferred where the instrument itself is the discourse-relevant participant:

\pex
\a
\begingl
\gla Marek houlam se piašček.//
\glb Marek.\Dir{} spoon.\Ind{} \textsc{inst} eat-\Antip{}-\Pf{}//
\glft \trsl{Marek ate with a spoon.} \lingcontext{instrument remains adjunct}//
\endgl
\a
\begingl
\gla Marek houlam piaštébime.//
\glb Marek.\Dir{} spoon.\Dir{} eat-\Appl{}-\Prog{}//
\glft \trsl{Marek is using the spoon for eating / is eating with the spoon as the relevant implement.}//
\endgl
\xe

\noindent In the imperative and hortative the same promotion yields benefactive readings by default. Directive clauses require the beneficiary to be overtly realized; the applicative is ungrammatical without one (see §\,\ref{sec:mood-directive-valency}):

\pex
\a
\begingl
\gla Trava bylam piaštébím.//
\glb bread.\Dir{} child.\Ind{} eat-\Appl{}-\Imp{}//
\glft \trsl{Feed the child bread.}//
\endgl
\a
\begingl
\gla Trava bylam piaštébká.//
\glb bread.\Dir{} child.\Ind{} eat-\Appl{}-\Hort{}//
\glft \trsl{Let's feed the child bread.}//
\endgl
\xe

\noindent The antipassive and the applicative are in complementary distribution. This is not an accidental gap: the antipassive removes the patient from core argument structure, whereas the applicative creates a new core object out of a peripheral participant. The two operations compete for the same VoiceP slot and cannot be combined on a single finite stem.


\subsubsection{Minimalist analysis: the applicative as ApplP}
\label{sec:vm-appl}

The applicative (\ird{-éb-}) promotes a peripheral (non-core)
participant to absolutive status. Following Pylkkänen
(2008), this can be analyzed as a
\term{high applicative} head (ApplP) that introduces the promoted
argument into the event structure between VoiceP and CausP:

\begin{center}
\begin{forest}
for tree={align=center}
[ApplP
  [{DP\\{\footnotesize applied argument}}]
  [{Appl$'$}
    [{Appl\\\ird{-éb-}}]
    [VoiceP
      [{\ldots}]
    ]
  ]
]
\end{forest}
\end{center}

\noindent The beneficiary, goal, or instrument promoted by the
applicative becomes the absolutive argument of the resulting vP,
effectively taking the slot that would otherwise be occupied by the
patient. The interaction with discourse is mediated at the AspP
level: which of the available absolutive candidates (patient or
applied argument) becomes the topic absolutive is determined by the
Asp.Align morpheme selected (§\,\ref{sec:vm-asp}).

%% STUB: The question of whether Iridian uses a high or low
%% applicative (Pylkkänen's distinction) should be tested
%% diagnostically. High applicatives introduce an argument that
%% relates to the event; low applicatives introduce an argument
%% related directly to the patient. In Iridian, beneficiary
%% promotion appears to be high (no entailment of patient
%% affectedness), while instrument promotion may be low.
%% This distinction bears on whether there are two distinct ApplP
%% projections, or a single one with type-sensitive semantics.


\subsection{The causative (\ird{-na-})}

Cross-linguistically, causation may be expressed lexically, analytically, or morphologically. Some languages employ distinct lexical verbs comparable to English \textit{kill}, while others use periphrastic constructions equivalent to \textit{make} or \textit{cause} plus a second predicate. Iridian belongs to the third type: causation is expressed by a productive morphological operation, realized synchronically as the suffix \ird{-na-}.

The causative \ird{-na-} adds an external causer argument to the event. In nominative-alignment contexts the causer stands in the direct case, the original agent or sole participant is demoted to the indirect, and any patient licensed by the base verb remains indirect. In ergative-alignment contexts the causer takes the transitive case. With transitive bases, the original patient remains the direct argument and the causee falls into the indirect; with intransitive bases, the causee itself occupies the direct as the sole affected participant. The causative thus increases valency by one while preserving the aspect-driven alignment pattern of the clause as a whole. Before a following vowel-initial aspect or mood suffix, \ird{-na-} inserts an epenthetic glide and surfaces as \ird{-naj-}; before consonant-initial material, including the Voice head \ird{-n-}, it remains \ird{-na-}.

\pex
\a
\begingl
\gla Marek bylam zmacnajime.//
\glb Marek.\Dir{} child.\Ind{} run-\Caus{}-\Prog{}//
\glft \trsl{Marek is making the child run.} \lingcontext{causative of intransitive base; nominative alignment}//
\endgl
\a
\begingl
\gla Travat Marku bylam piašt-na-nik.//
\glb bread.\Dir{}.\Def{} Marek.\Trs{} child.\Ind{} eat-\Caus{}-\textsc{trs}-\Pf{}//
\glft \trsl{Marek made the child eat the bread.} \lingcontext{causative of transitive base; ergative alignment}//
\endgl
\xe

\noindent The semantic character of the causation is expressed primarily through aspect. In the perfective the causative presents the causer as exercising effective control over the caused event and therefore tends toward direct, interventionist, or coercive readings. In the retrospective the same morphology presents a looser relation between causer and event, so that the clause more readily receives indirect, delegated, accidental, or permissive interpretations. For the discourse consequences of this contrast, see Chapter~\ref{ch:discourse}.

\pex
\a
\begingl
\gla Marek bylam piaštnajime.//
\glb Marek.\Dir{} child.\Ind{} eat-\Caus{}-\Prog{}//
\glft \trsl{Marek is having the child eat / is feeding the child.}//
\endgl
\a
\begingl
\gla Travat Marku bylam piašt-na-n-ac.//
\glb bread.\Dir{}.\Def{} Marek.\Trs{} child.\Ind{} eat-\Caus{}-\textsc{trs}-\Ret{}//
\glft \trsl{Marek has let the child eat the bread / has ended up getting the child to eat it.}//
\endgl
\xe

\noindent The causative is fully compatible with the applicative, provided the applicative is formed first in accordance with the morphological template:

\pex
\a
\begingl
\gla Marek bylam piaštébnajime.//
\glb Marek.\Dir{} child.\Ind{} eat-\Appl{}-\Caus{}-\Prog{}//
\glft \trsl{Marek is making someone feed the child.}//
\endgl
\a
\begingl
\gla Trava bylam piaštnajím.//
\glb bread.\Dir{} child.\Ind{} eat-\Caus{}-\Imp{}//
\glft \trsl{Make the child eat the bread.} \lingcontext{directive; causee in \Ind{} required}//
\endgl
\a
\begingl
\gla Trava bylam piašt-na-ká.//
\glb bread.\Dir{} child.\Ind{} eat-\Caus{}-\Hort{}//
\glft \trsl{Let's make the child eat the bread.}//
\endgl
\xe

\noindent Synchronically, \ird{-na-} is distinct from the transitive/passive Voice head \ird{-n-}. The former is a valency-changing derivation occupying the causative slot of the verbal template (CausP); the latter is the Voice head (VoiceP) that licenses the internal argument and specifies patient absolutive, inserted immediately before the aspect suffix. Earlier notes used \ird{-ne-} for the causative, but the present grammar regularises \ird{-na-} in order to keep the causative head formally distinct from \ird{-n-}.


\subsubsection{Minimalist analysis: the causative as CausP}
\label{sec:vm-caus}

The causative (\ird{-na-}) introduces an external causer above the
base event structure, corresponding to a \term{CausP} projection
that dominates VoiceP/ApplP. The causer argument is merged in
Spec,CausP:

\begin{center}
\begin{forest}
for tree={align=center}
[CausP
  [{DP\\{\footnotesize causer}}]
  [{Caus$'$}
    [{Caus\\\ird{-na-}}]
    [ApplP
      [{\ldots}]
    ]
  ]
]
\end{forest}
\end{center}

\noindent The morphological order (applicative before causative
before Asp.Align) is predicted: CausP dominates VoiceP/ApplP, and
the Mirror Principle assigns them to the innermost positions relative
to AspP. The interaction with aspect is semantically significant
(§\,\ref{sec:aspect-caus}): perfective causation presents direct
causer control (CausP features trigger tight Agree with the event
head), while retrospective causation presents indirect or
permissive causation (looser Agree relation, admitting delegated or
accidental readings).

\subsubsection{The transitive/passive Voice head \ird{-n-}}
\label{sec:vm-n}

The morpheme \ird{-n-} is an overt Voice head occupying the same
VoiceP slot as the antipassive \ird{-aš-} and the applicative
\ird{-éb-}. It carries two properties simultaneously: it licenses
the internal argument (patient) as a core participant in the clause's
argument structure, and it specifies that the patient is to stand as
the absolutive (\textsc{[pat-abs]}). These two properties are not
environment-specific; they are always present. What varies across
aspectual domains is their structural consequence within the
Asp.Align system.

\paragraph{ERG-aspect behaviour.}
In the perfective and retrospective aspects the default absolutive
orientation of Asp.Align already favours the patient (§\,\ref{sec:vm-asp-case}).
The \textsc{[pat-abs]} specification contributed by \ird{-n-} is
therefore harmonious with the aspectual default. The indispensable
contribution of \ird{-n-} in this environment is accordingly not to
override the default but to license the patient as a core argument:
without an overt Voice head in VoiceP, a transitive predicate cannot
realize its patient as a properly licensed participant in the clause's
core argument structure. The bare perfective or retrospective is
consequently illicit for any transitive reading, as described in
§\,\ref{sec:ambitransitive}. With \ird{-n-} present, the patient is
licensed as the absolutive and the agent takes the transitive case;
the result is the ordinary ergative transitive construction.

\begin{center}
\begin{forest}
for tree={align=center}
[AspP
  [{DP\\{\footnotesize patient}\\{\footnotesize\textsc{[abs]}}}]
  [{Asp$'$}
    [{Asp\\{\footnotesize\textsc{[erg]}}}]
    [VoiceP
      [{DP\\{\footnotesize agent}\\{\footnotesize\textsc{[trs]}}}]
      [{Voice$'$}
        [{Voice\\\ird{-n-}\\{\footnotesize\textsc{[pat-abs]}}}]
        [VP
          [{DP\\{\footnotesize patient}}]
          [{\textit{root}}]
        ]
      ]
    ]
  ]
]
\end{forest}
\end{center}

\paragraph{NOM-aspect behaviour.}
In the continuous, progressive, and contemplative aspects the same
overt Voice head overrides the default agent-absolutive orientation
of Asp.Align. Because \ird{-n-} specifies \textsc{[pat-abs]}, the
patient rather than the agent becomes the direct-case pivot, and the
agent is demoted to the oblique domain. The result is a
\term{passive} construction. Asp.Align remains in the nominative
domain; it is the Voice head that imposes a patient-absolutive
realization upon the clause. The full analysis and examples are in
§\,\ref{sec:vm-passive}.

\ird{-n-} is thus not a pair of context-sensitive morphemes whose
function changes by aspect. It is a single Voice head whose constant
properties --- patient licensing and patient absolutive --- produce
different surface constructions depending on whether the aspectual
domain reinforces or is overridden by its absolutive specification.
The phonological resemblance to the causative \ird{-na-} is
addressed in §\,\ref{sec:vm-tension-n}.


\subsection{The passive}\label{sec:vm-passive}

The passive is the construction that arises when the transitive/passive
Voice head \ird{-n-} is merged in a NOM-aspect clause. Since
\ird{-n-} specifies patient absolutive (\textsc{[pat-abs]}) and
the NOM-aspect domain would by default assign absolutive to the agent,
the Voice head overrides the aspectual default: the patient surfaces
in the direct case as the absolutive pivot, and the agent is demoted
to the indirect (oblique) case. Asp.Align remains in the
nominative domain throughout; only the absolutive orientation of
the clause is altered by the overt Voice head.

The passive is thus the marked construction by which a
nominative-domain clause is reorganized around the patient, allowing
the patient rather than the agent to function as the discourse pivot.
Its counterpart in the ergative domain is the ordinary ergative
transitive, which arises from the same Voice head under ERG-aspect
alignment. The two constructions are formally identical at the Voice
level; their different surface alignment is a consequence of the
aspectual domain they are embedded in.

\subsubsection*{Morphological form}

The passive consists of a NOM-aspect suffix on a transitive stem
bearing \ird{-n-}. The NOM-aspect suffixes available for passive
formation are the continuous \ird{-a}, the progressive \ird{-ime},
and the contemplative \ird{-ach}; the patient-absolutive
specification of \ird{-n-} is realized in each case.

\subsubsection*{Case frame}

The patient surfaces in the direct case (\textsc{abs}/\textsc{dir});
the agent, displaced from the absolutive, surfaces in the indirect
case (\textsc{obl}/\textsc{ind}). The agent may be omitted when it
is recoverable or irrelevant.
\begin{center}
\begin{forest}
for tree={align=center}
[AspP
  [{DP\\{\footnotesize patient}\\{\footnotesize\textsc{[abs]}}}]
  [{Asp$'$}
    [{Asp\\{\footnotesize\textsc{[nom]}}}]
    [VoiceP
      [{DP\\{\footnotesize agent}\\{\footnotesize\textsc{[obl]}}}]
      [{Voice$'$}
        [{Voice\\\ird{-n-}\\{\footnotesize\textsc{[pat-abs]}}}]
        [VP
          [{DP\\{\footnotesize patient}}]
          [{\textit{root}}]
        ]
      ]
    ]
  ]
]
\end{forest}
\end{center}

\noindent Asp.Align remains nominative-domain (\textsc{Asp[nom]}),
but the overt Voice head requires patient absolutive; the agent
therefore surfaces in the oblique domain.

\subsubsection*{Examples}

With agent expressed in the oblique:

\pex
\a
\begingl
\gla Travat Marku piastnime.//
\glb bread.\Dir{}.\Def{} Marek.\Ind{} eat-\textsc{pass}-\Prog{}//
\glft \trsl{The bread is being eaten by Marek.} \lingcontext{passive; P = \Dir{}, A = \Ind{}}//
\endgl
\a
\begingl
\gla Vytam Marku shodnime.//
\glb shirt.\Dir{}.\Def{} Marek.\Ind{} wear-\textsc{pass}-\Prog{}//
\glft \trsl{The shirt is being put on by Marek.}//
\endgl
\xe

With agent omitted:

\pex
\a
\begingl
\gla Travat piastnime.//
\glb bread.\Dir{}.\Def{} eat-\textsc{pass}-\Prog{}//
\glft \trsl{The bread is being eaten.} \lingcontext{passive; agent omitted}//
\endgl
\a
\begingl
\gla Doma piasčna.//
\glb house.\Dir{} build-\textsc{pass}-\Cont{}//
\glft \trsl{The house is being built / is under construction.}//
\endgl
\xe

\noindent The discourse value of the passive is to constitute the
patient as the clause's absolutive topic, making it the argument
that subsequent clauses may pick up pronominally or by topic drop.
The agent, when expressed, appears in the indirect case as a
non-core oblique participant.


\subsection{Pragmatic use of valency operations}

The choice among ergative, antipassive, applicative, and causative is not purely
mechanical. Each construction offers the speaker a distinct way of packaging the
event for discourse. The antipassive is used when the patient is generic,
non-referential, semantically incorporated into the activity, recoverable from
context, or intentionally backgrounded. Its effect is to keep the agent in
absolutive position and to present the clause as being about what the agent did
rather than about what happened to some patient.

The applicative serves the opposite discourse function: it pulls a formerly
peripheral participant into the centre of the clause. A beneficiary,
destination, location, or instrument is promoted when that participant is the
one the discourse is tracking, when it will remain topical in subsequent
clauses, or when the speaker wishes to construe the event from its perspective
rather than from that of an unexpressed or backgrounded patient.

The causative likewise has a pragmatic life beyond its valency-changing effect.
By choosing a causative rather than a simple predicate, the speaker frames the
event in terms of inducement, delegated agency, or control. In direct contexts
this may foreground authority or coercion; in weaker retrospective contexts it
may soften the causer's responsibility, presenting the caused event as something
merely permitted, arranged, or indirectly brought about rather than personally
carried out.

The applicative and causative are compatible with the imperative and hortative
and occupy their regular positions in the morphological template before the mood
suffix.


\section{Negation}\label{sec:negation}

The negative prefix \ird{zá-} --- reduced to \ird{za-} in unstressed environments --- stands before the entire verbal complex. In the indicative it accordingly precedes the stem itself. In the imperative, hortative, and irrealis, negation is expressed by distinct mood-specific negative suffixes rather than by \ird{zá-}; see Section~\ref{sec:mood} and §\,\ref{sec:vm-moodp}.

This section describes verbal negation --- the use of \ird{zá-} with finite indicative predicates --- exclusively. Negation in equative clauses (zero-copula constructions) and existential clauses (the \ird{jec-}/\ird{zad-} alternation) is treated in §\,\ref{sec:clausal-negation} of Chapter~\ref{ch:simple}.

One framing point is worth flagging before the formal analysis. \ird{Zá-} is the most morphologically prominent of the three NegP exponents in Iridian, but it is also the structurally anomalous one: it surfaces as a prefix rather than in the head-final position that NegP occupies in the clause spine. In equative clauses, by contrast, Neg\textsuperscript{0} is realized as the free morpheme \ird{staly}, which surfaces exactly where the clause spine predicts --- clause-finally, after its two nominal arguments. \ird{Staly} thus provides the clearest direct evidence for the head-final status of NegP; the prefix position of \ird{zá-} is a PF-level cliticization effect, not the canonical position of the negation head. The full argument is in §\,\ref{sec:vm-neg-exponents}.

The absence of \ird{zá-} in the prohibitive (\ird{-é}), the negative hortative (\ird{-ku}), and the negative irrealis (\ird{-ir}) is not arbitrary. Each of these suffixes carries [+\textsc{neg}] as a lexically inherent feature of the MoodP Vocabulary Item; when such a head is selected, NegP is simply absent from the derivation and \ird{zá-} has no Neg\textsuperscript{0} node to spell out. There is no separate blocking principle at work: the absence of NegP is the full explanation. See §\,\ref{sec:vm-moodp} for the MoodP paradigm and analysis.


\subsection{Formal analysis: NegP and the prefix paradox}
\label{sec:vm-neg}
\label{sec:vm-neg-exponents}

The negative prefix \ird{zá-} stands before the entire verbal stem.
Within the clause spine (Table~\ref{tab:th-clause-spine}), NegP
occupies projection 4, above AspP but below IP. If the Mirror
Principle held perfectly, \ird{zá-} should appear \emph{after} the
Asp.Align suffix, since NegP dominates AspP and therefore corresponds
to a more peripheral morpheme. Instead it appears as a prefix, inside
the verbal complex to the left of the root.

Three analyses are possible. Under V-to-Neg head movement, the verbal
complex undergoes further movement to the Neg head in narrow syntax,
making \ird{zá-} the Neg head with the raised complex to its right;
but in a head-final language the Neg head should follow its complement,
predicting \ird{*piaštek zá-}, so the account requires either
surface-right adjunction or a PF reordering rule. A scope-based prefix
convention (Bybee 1985) treats the prefix position as a cross-linguistic
tendency for high-scope operators to surface as prefixes in verb-final
languages (cf.\ Turkish, Mongolian, Basque \textit{ez}), making
the Mirror Principle simply inapplicable to negation. The preferred
account is DM morphological incorporation at PF (Embick \& Noyer 2007):
NegP is a syntactic head above AspP, but \ird{zá-} is specified as a
prefix and at PF phonologically cliticizes to the left edge of the
verbal complex, leaving the syntactic hierarchy intact. The
procliticization itself is direct evidence: \ird{zá-} does not
form an independent prosodic word.

That the DM account is correct emerges most clearly from a unified
treatment of Neg\textsuperscript{0} across all clause types. Treating
\ird{zá-} as the sole reflex of NegP is incomplete: Neg\textsuperscript{0}
has three PF exponents conditioned by predicate type. With a finite
verbal predicate, Neg\textsuperscript{0} is realized as the procliticized
prefix \ird{zá-} attached to the left edge of the verbal complex; because
that complex is itself clause-final by head-finality, the prefix surfaces
immediately before the root, yielding propositional negation of the event
description. With a zero-copula equative or class-membership predicate,
there is no verbal host for \ird{zá-} to cliticize onto, and
Neg\textsuperscript{0} is instead spelled out as the free morpheme
\ird{staly}, which surfaces in the canonical head-final position following
its two nominal arguments --- the position expected of any clause-final
head, requiring no additional movement rule. When the predicative
complement is a plain DP the result is \term{predicational negation}
(\ird{Marek doktor staly}, \trsl{Marek is not a doctor}); when it is a
participant nominal (§\,\ref{sec:nmlz-predicative}) the result is
\term{identificational negation} (§\,\ref{sec:vm-ident-neg} below).
With the existential verb \ird{jec-}, Neg\textsuperscript{0} triggers a
DM readjustment: the existential root is realized as the suppletive
allomorph \ird{zad-} and Neg\textsuperscript{0} itself has no independent
PF exponent, negation being absorbed into the root via Vocabulary
Insertion --- the standard behavior of suppletive negative existentials
cross-linguistically (cf.\ Turkish \textit{yok}/\textit{var};
Japanese \textit{nai}/\textit{aru}).

The complete distribution, including mood-internal negation, is
summarized in Table~\ref{tab:neg-exponents} of Chapter~\ref{ch:simple}.
The three environments are a conditioned allomorphy of a single
Neg\textsuperscript{0} head, not distinct lexical negation strategies.
NegP occupies the same structural position (above AspP, below MoodP)
in all clause types; what varies is the PF exponent of
Neg\textsuperscript{0} and the semantic scope of negation, which is
determined by the semantics of the predication NegP takes in its scope.

\begin{figure}[ht]
  \footnotesize\sffamily
  \caption{Comparative exponence of Neg\textsuperscript{0} across clause types. In all three structures NegP occupies the same position above AspP or the zero-copula predication; what varies is the PF realization of Neg\textsuperscript{0}. Only the verbal clause shows the prefix-paradox effect, with \ird{zá-} cliticized to the left edge of the verbal complex rather than surfacing in canonical head-final position.}\label{fig:neg0-comparative}
  \medskip
  \centering
  \scriptsize
  \begin{minipage}[t]{0.30\textwidth}
    \centering
    \textsc{Verbal clause}\par
    \medskip
    \begin{forest}
      syntax tree, for tree={l sep=8pt,s sep=6pt,inner sep=0.5pt}
      [NegP
        [{DP\\{\footnotesize topic}}]
        [{Neg$'$}
          [{Neg$^{0}$}\\\ird{zá-}\\{\small(PF prefix)}]
          [AspP
            [{Asp$'$}
              [{Asp\\{\footnotesize\textit{asp.align}}}]
              [VP
                [{V\\{\footnotesize\textit{root}}}]
              ]
            ]
          ]
        ]
      ]
    \end{forest}
  \end{minipage}
  \hfill
  \begin{minipage}[t]{0.30\textwidth}
    \centering
    \textsc{Equative clause}\par
    \medskip
    \begin{forest}
      syntax tree, for tree={l sep=8pt,s sep=6pt,inner sep=0.5pt}
      [NegP
        [{DP\\{\footnotesize topic}}]
        [{Neg$'$}
          [{Neg$^{0}$\\\ird{staly}}]
          [EquP
            [{DP\\{\footnotesize pred.}}]
            [{Eq$'$}
              [{Eq$^{0}$\\{$\varnothing$}}]
            ]
          ]
        ]
      ]
    \end{forest}
  \end{minipage}
  \hfill
  \begin{minipage}[t]{0.30\textwidth}
    \centering
    \textsc{Existential clause}\par
    \medskip
    \begin{forest}
      syntax tree, for tree={l sep=8pt,s sep=6pt,inner sep=0.5pt}
      [NegP
        [{DP\\{\footnotesize pivot}}]
        [{Neg$'$}
          [{Neg$^{0}$\\{\small($\varnothing$: absorbed)}}]
          [AspP
            [{Asp$'$}
              [{Asp\\{\footnotesize\textit{asp}}}]
              [ExistP
                [{Exist$'$}
                  [{Exist$^{0}$\\\ird{zad-}\\{\small(suppletive)}}]
                ]
              ]
            ]
          ]
        ]
      ]
    \end{forest}
  \end{minipage}
\end{figure}

An important consequence follows for the prefix paradox. The paradox
arises specifically for verbal predicates, where \ird{zá-} must attach
to a phonological host; Figure~\ref{fig:neg0-comparative} makes clear
that this is the outlier among the three exponents, not the structural
default. In equative clauses, where no such host exists, \ird{staly}
surfaces in the expected head-final position --- after its complement
(the nominal predication) --- and therefore gives the clearest direct
evidence for the head-final status of NegP. In existential clauses,
\ird{zad-} shows the same Neg\textsuperscript{0} head being discharged
through suppletive Vocabulary Insertion at Exist\textsuperscript{0}
rather than by an independent overt exponent. The prefixhood of
\ird{zá-} is therefore a PF-level cliticization effect, not the
canonical behavior of Neg\textsuperscript{0}.

For equative clauses specifically, the word order reflects the
head-final NegP structure transparently:

\begin{center}
\begin{forest}
  syntax tree, for tree={l sep=8pt,s sep=6pt,inner sep=0.5pt}
  [NegP
    [{DP\\{\footnotesize topic}}]
    [{Neg$'$}
      [{DP\\{\footnotesize pred.}}]
      [{Neg$^{0}$\\\ird{staly}}]
    ]
  ]
\end{forest}
\end{center}

\noindent The two nominal constituents before \ird{staly} may permute
in accordance with the topic-comment architecture (§\,\ref{sec:topic-comment}),
while \ird{staly} is invariably clause-final.

The incompatibility of \ird{zad-} with \ird{zá-} also follows without additional
stipulation. If Neg\textsuperscript{0} triggers the suppletive Vocabulary
Insertion of \ird{zad-} at the root node, then Neg\textsuperscript{0} has
already been discharged structurally; stacking a second Neg\textsuperscript{0}
on top would create an uninterpretable double-negation configuration. The
prohibition is a consequence of the analysis rather than a separate blocking
rule.

The relationship between Neg\textsuperscript{0} and the inherently negative mood
forms (prohibitive \ird{-é}, negative hortative \ird{-ku}, negative irrealis
\ird{-ir}) parallels the \ird{zad-} case in its effect on NegP projection: when
MoodP carries [+\textsc{neg}] as a VI-internal feature, NegP is absent from the
derivation and \ird{zá-} accordingly does not surface. The full analysis of
MoodP [+\textsc{neg}] --- including the semantic distinction between
directive-scoped and propositional negation, and the restricted palatalizing
behavior of \ird{-ir} --- is in §\,\ref{sec:vm-moodp}.

The DM morphological-incorporation account makes the following empirical
commitments. (a)~\ird{Zá-} does not constitute an independent prosodic word: it
is a bound clitic on the verbal complex and cannot receive independent stress or
be separated from the stem by an intervening focus particle. (b)~Nothing may
intervene between \ird{zá-} and the stem; evidence of such separation would
require the cliticization analysis to be revised. (c)~In ellipsis and fragment
contexts, \ird{zá-} cannot appear stranded without the full verbal complex,
unlike a free-standing negative head that could survive ellipsis of its
complement. Two further commitments bear on the three-exponent analysis:
(d)~\ird{staly} is fully invariant, bearing no aspect, person, mood, or
iterative marking; paradigmatic variation in future data would favour a
defective-verbal-root analysis over the Neg\textsuperscript{0} account. (e)~The
three exponents (\ird{zá-}, \ird{staly}, \ird{zad-}) are in strict complementary
distribution, conditioned by the presence or absence of a finite verbal host;
the equative/stative boundary is treated in §\,\ref{sec:vm-tension-stative}.
Cross-linguistic typology is consistent with this picture: prefix negation is
well attested in SOV languages (Dryer 2013), suggesting the prefix position of
\ird{zá-} reflects a systematic PF convention rather than a Mirror Principle
violation.

\subsubsection{Identificational negation and the scope of \ird{staly}}
\label{sec:vm-ident-neg}

The availability of participant nominals as equative predicates
(§\,\ref{sec:nmlz-predicative} of Chapter~\ref{ch:verbs}) introduces a
semantically significant variation within the \ird{staly} environment. The
contrast can be stated precisely in terms of what NegP takes in its scope.

In a plain equative (\ird{Marek doktor}), the predication is
\term{predicational}: it ascribes a property or class to the topic. The negated
version (\ird{Marek doktor staly}) denies property-membership. NegP takes a
property-ascription in its scope and yields propositional negation of that
ascription.

In an identificational equative (\ird{Lobek Jancí piaštnikou}, \trsl{The apple
is what Janek ate}), the predication is \term{specificational} or
\term{identificational}: it asserts that the topic DP is the unique entity
satisfying the nominalized description. The semantic content of the predicate
(the participant nominal) is inherently exhaustive --- it picks out exactly one
referent. NegP over a specificational predication yields \term{identificational
negation}:

\begin{center}
$\lnot$[\textsc{topic} = \textsc{unique-satisfier}(\textit{nom.\ description})]
\end{center}

The contrastive implication --- that some other entity is the actual satisfier
--- follows from the exhaustivity of the original predication, not from any
special semantic feature of \ird{staly}. This is an automatic consequence of
negating an exhaustive identification, parallel to the contrastive implication
of cleft negation in English (\textit{It was not the apple that Janek ate}).

Crucially, the semantic difference between propositional \ird{zá-} and
identificational \ird{staly}+nominal is not ambiguity in the NegP head but a
difference in what NegP scopes over. \ird{Zá-} scopes over a finite event
description; \ird{staly} over a (zero-copula) equative predication, which may be
either property-ascribing or identificational depending on the type of its
predicate DP.

\begin{tabularx}{\linewidth}{l L L}
\toprule
& \textsc{propositional neg.\ (\ird{zá-})} & \textsc{identificational neg.\ (\ird{staly}+nom.)} \\
\midrule
Event presupposed?       & no                        & yes \\
Contrastive implication? & no                        & yes \\
NegP scope               & event-VP                  & equative predication \\
PF of Neg\textsuperscript{0} & prefix to verbal complex & free morpheme, clause-final \\
\bottomrule
\end{tabularx}

\medskip

he distinction between these rows is what the user should consult when deciding
which construction to employ: if the question is whether an event occurred,
\ird{zá-} is used; if the question is which entity was involved in a presupposed
event, \ird{staly} with a participant nominal is used. Both are structurally
expressions of a single NegP.

Note also that the nominal-domain negative operators (\ird{zide} \trsl{no one},
\ird{niho} \trsl{nothing}, etc.; see §\,\ref{sec:nm-pron-operators} of
Chapter~\ref{ch:nouns}) are not part of the clausal NegP system described here.
Their interaction with the NegP projection --- in particular, whether they
project their own NegP, co-occur with \ird{zá-}, or block NegP projection --- is
an open descriptive question discussed at the relevant section of
Chapter~\ref{ch:nouns}.

\section{Mood}\label{sec:mood}

The mood system comprises the indicative, the imperative, the prohibitive, the
hortative, and the irrealis. Mood is expressed by a suffix occupying the
position immediately following the aspect-alignment morpheme; in the indicative
this suffix is null. Where a mood suffix is present, the aspect-alignment slot
is suppressed, the mood suffix replacing the full aspect complex. The
imperative, prohibitive, and hortative together form the directive domain, a
unified subsystem with its own case properties; it is treated as a group
beginning in §\,\ref{sec:mood-directive}.

\subsection{Indicative}

The indicative is the unmarked mood. It alone hosts the full five-way aspect
system and the split-ergative alignment described in Section~\ref{sec:aspect};
it is accordingly the only mood in which the antipassive-transitive distinction
is operative.

\subsection{The directive domain}\label{sec:mood-directive}

The imperative, the prohibitive, and the hortative form a unified directive
domain: a compact finite subsystem used for commands, negative commands,
suggestions, proposals, soft requests, permissive uses, and wishes. This
subsystem operates on principles distinct from those of the indicative. Where
the indicative couples mood with the full five-way aspect system and the
split-ergative alignment described in Section~\ref{sec:aspect}, directive moods
exclude the ordinary aspect suffixes entirely. The directive mood suffix is
itself clause-defining and therefore occupies the aspect-alignment slot in the
verbal template; no aspect morpheme is present alongside it.

The directive domain is further governed by a case principle not operative in
the indicative: directive clauses ordinarily permit only one core nominal in the
direct case. The understood subject of imperative and prohibitive clauses is
second person and is normally left unexpressed, being recoverable from the
speech situation. Where an overt speech-act participant subject would otherwise
compete with an overt direct object for the clause's single direct slot, that
subject is expressed in the transitive case instead. The patient or theme of the
clause is thus the clause's sole core direct nominal. Beneficiaries licensed by
the applicative and causees licensed by the causative are obligatorily overt and
appear in the indirect. Antipassive morphology is excluded from all directive
moods.

\subsection{Imperative and prohibitive}\label{sec:mood-imp}

The imperative (\ird{-ím}) is the ordinary mood for direct commands,
instructions, warnings, and immediate discourse-management directives addressed
to the hearer. The prohibitive (\ird{-é}) is its dedicated negative counterpart,
used for negative commands and preventative directives. The two are treated
together because they share the same null-subject pattern and the same case
geometry.

The subject of both the imperative and the prohibitive is second person and is
syntactically understood but ordinarily unexpressed; no overt pronoun is
required in fluent discourse. The theme or patient of the clause occupies the direct case as the
clause's single core direct nominal:

\pex
\a
\begingl
\gla Trava piaštím.//
\glb bread.\Dir{} eat-\Imp{}//
\glft \trsl{Eat the bread.}//
\endgl
\a
\begingl
\gla Tak hrabním.//
\glb \Dem{}.\Prox{}.\Dir{} hear-\Imp{}//
\glft \trsl{Hear this.}//
\endgl
\xe

The prohibitive carries the negative force within the suffix itself and does not
co-occur with the ordinary negative prefix \ird{zá-}. This is a structural
consequence of the morphology: the prohibitive suffix bears the inherent
negative value of the MoodP Vocabulary Item, rendering NegP absent from the
derivation; see §\,\ref{sec:vm-moodp} for the formal account. A prohibition
accordingly takes the following form:

\pex
\a
\begingl
\gla Trava piašté.//
\glb bread.\Dir{} eat-\Proh{}//
\glft \trsl{Do not eat the bread.}//
\endgl
\xe

Overt second-person subjects are permitted in both the imperative and the
prohibitive, but they are pragmatically marked, expressing foregrounding,
contrast, insistence, or rebuke toward the addressee. The second-person singular
pronoun is \ird{já} in the direct and \ird{jí} in the transitive; the
second-person plural is \ird{tová} in the direct and \ird{tvy} in the
transitive. In a transitive imperative or prohibitive clause where an overt
direct object is already present, an overt speech-act participant subject
appears in the transitive case rather than the direct, so that the affected
participant may remain the clause's sole direct nominal:

\pex
\a
\begingl
\gla Jí trava piaštím.//
\glb 2\textsc{sg}.\Trs{} bread.\Dir{} eat-\Imp{}//
\glft \trsl{You eat the bread.} \lingcontext{marked; overt 2\textsc{sg} subject in \Trs{}}//
\endgl
\a
\begingl
\gla Tvy trava piašté.//
\glb 2\textsc{pl}.\Trs{} bread.\Dir{} eat-\Proh{}//
\glft \trsl{You lot do not eat the bread.} \lingcontext{marked; overt 2\textsc{pl} subject in \Trs{}}//
\endgl
\xe

\noindent The pressure for the transitive case arises specifically from
competition: an overt subject would otherwise displace the direct object from
the clause's single direct slot. In intransitive directive clauses lacking a
competing direct object, no such competition arises, and the rule does not force
the transitive. The rule is therefore competition-sensitive rather than
absolute.

\subsection{Hortative}\label{sec:mood-hort}

The hortative covers a broader range of directive uses than the imperative:
inclusive proposals of the type \textit{let us}, softened direct commands,
permissive uses of the type \textit{let the child come}, and optative
expressions of wish or hope. It serves as the default polite register in
ordinary speech, typically preferred over the imperative in everyday
conversation. The hortative takes two forms: the positive \ird{-ká} and the
negative \ird{-ku}. Like the imperative and prohibitive, it excludes the
ordinary aspect suffixes and operates within the single-direct principle of the
directive domain.

The hortative defaults to a null subject whose interpretation is supplied by
clause type and context. In proposal uses the null subject is most readily read
as first person plural, inclusive of the addressee; in softened directive uses
the same null subject may receive a second-person reading. Null-subject
hortatives are the unmarked form:

\pex
\a
\begingl
\gla Piaštká.//
\glb eat-\Hort{}//
\glft \trsl{Let's eat.} \lingcontext{null subject; inclusive proposal}//
\endgl
\a
\begingl
\gla Tak hrabnká.//
\glb \Dem{}.\Prox{}.\Dir{} hear-\Hort{}//
\glft \trsl{Please hear this.} \lingcontext{null subject; softened directive with direct object}//
\endgl
\xe

The negative hortative \ird{-ku} is used for proposals to refrain, cautious
collective restraint, and mild negative requests. Like the prohibitive, it is an
inherently negative Vocabulary Item and accordingly excludes \ird{zá-}:

\pex
\a
\begingl
\gla Tak piaštku.//
\glb \Dem{}.\Prox{}.\Dir{} eat-\Neg{}\Hort{}//
\glft \trsl{Let's not eat this.} \lingcontext{negative hortative; collective restraint}//
\endgl
\xe

Overt subjects in the hortative fall into two structurally distinct classes. The
first class comprises overt speech-act participant subjects --- first person
plural, second person singular, second person plural --- appearing in proposal
or soft-command uses. These behave exactly like marked subjects in the
imperative and prohibitive: where they would otherwise compete with an overt
direct object for the clause's single direct slot, they appear in the transitive
case:

\pex
\a
\begingl
\gla Mně trava piaštká.//
\glb 1\textsc{pl}.\Trs{} bread.\Dir{} eat-\Hort{}//
\glft \trsl{Let us eat the bread.} \lingcontext{overt 1\textsc{pl} subject in \Trs{}}//
\endgl
\a
\begingl
\gla Jí tak hrabnká.//
\glb 2\textsc{sg}.\Trs{} \Dem{}.\Prox{}.\Dir{} hear-\Hort{}//
\glft \trsl{Please, you hear this.} \lingcontext{marked; overt 2\textsc{sg} subject in \Trs{}}//
\endgl
\xe

The second class comprises overt lexical subjects in optative or permissive
hortatives, where the clause expresses a wish or permission directed toward a
third participant rather than a command to the addressee. Such clauses are
typically intransitive. Their lexical subject occupies the direct case as the
clause's single core direct nominal --- the same position that in transitive
directive clauses would be held by the theme or patient. Because there is no
competing direct object, no competition for the direct slot arises and the
transitive case is not required:

\pex
\a
\begingl
\gla Byl ščenká.//
\glb child.\Dir{} arrive-\Hort{}//
\glft \trsl{May the child come.} \lingcontext{intransitive optative; lexical subject in \Dir{}}//
\endgl
\a
\begingl
\gla Je ščenká.//
\glb \Dem{}.\Dist{}.\Dir{} arrive-\Hort{}//
\glft \trsl{May that come to pass.} \lingcontext{intransitive optative; distal demonstrative in \Dir{}}//
\endgl
\xe

\noindent The contrast between these two subject classes captures the essential
geometry of the directive domain. An overt speech-act participant subject
appears in the transitive when a direct object is present because it would
otherwise compete with that object for the single direct slot. A lexical subject
of an intransitive optative or permissive hortative faces no such competition
and therefore remains in the direct. The reason for the difference is not
grammatical person but structural: in the first class a competing direct nominal
is present; in the second it is not.

\subsection{Case and subject expression in directive clauses}\label{sec:mood-directive-case}

The case properties of the directive domain may be stated compactly as a set of
rules that govern all three moods.

Directive clauses ordinarily permit only one core nominal in the direct case. In
transitive uses, that nominal is the theme or patient of the predicate. In
intransitive uses it is the clause's sole participant. The understood
second-person subject of imperative and prohibitive clauses, and the understood
subject of hortative clauses, are syntactically present but normally left
unexpressed.

Overt speech-act participant subjects are pragmatically marked throughout the
directive domain. Where their expression would otherwise compete with an overt
direct object, they appear in the transitive case rather than the direct,
preserving the affected participant as the clause's sole direct nominal. The
following minimal pair makes the contrast plain:

\pex
\a
\begingl
\gla Trava piaštím.//
\glb bread.\Dir{} eat-\Imp{}//
\glft \trsl{Eat the bread.} \lingcontext{null subject; theme in \Dir{}}//
\endgl
\a
\begingl
\gla Jí trava piaštím.//
\glb 2\textsc{sg}.\Trs{} bread.\Dir{} eat-\Imp{}//
\glft \trsl{You eat the bread.} \lingcontext{overt subject in \Trs{}; theme in \Dir{}}//
\endgl
\xe

\noindent By contrast, an overt lexical subject of an intransitive optative or
permissive hortative may remain in the direct when no other direct nominal is
present, since the single-direct principle is not under pressure. The
competition-sensitive nature of the rule means it neither applies to
intransitive clauses nor generalizes beyond the contexts in which direct-slot
competition actually arises.

\subsection{Valency-changing morphology in directive clauses}\label{sec:mood-directive-valency}

Applicative and causative morphology are compatible with the directive moods.
Both occupy their regular positions in the morphological template immediately
before the directive mood suffix, which is accordingly the outermost suffix in
the verbal complex within the mood domain. The directive suffix thus remains
outer to all valency morphology.

The applicative (\ird{-éb-}) is licensed in directive clauses only if the
beneficiary is overtly realized. The beneficiary appears in the indirect case,
typically with the postposition \ird{my}, while the theme retains the direct.
This fits the directive case geometry without disturbance: the direct holds the
theme, the indirect holds the beneficiary, and an overt addressee subject, if
present, occupies the transitive:

\pex
\a
\begingl
\gla Trava bylam piaštébím.//
\glb bread.\Dir{} child.\Ind{} eat-\Appl{}-\Imp{}//
\glft \trsl{Feed the child bread.} \lingcontext{applicative; beneficiary in \Ind{}}//
\endgl
\a
\begingl
\gla Jí trava bylam piaštébím.//
\glb 2\textsc{sg}.\Trs{} bread.\Dir{} child.\Ind{} eat-\Appl{}-\Imp{}//
\glft \trsl{You feed the child bread.} \lingcontext{overt subject in \Trs{}}//
\endgl
\xe

\noindent Without an overt beneficiary the applicative is ungrammatical in
directive clauses:

\pex
\rec{trava piaštébím} \trsl{intended: `Feed someone bread.'} \lingcontext{applicative without beneficiary}
\xe

The causative (\ird{-na-}) likewise requires an overtly realized causee in
directive clauses. The causee appears in the indirect case, keeping the direct
available for the theme. The directive case pattern for causatives therefore
parallels that for applicatives exactly: theme in the direct, causee in the
indirect, overt addressee subject in the transitive if present. Before the
vowel-initial mood suffixes (\ird{-ím}, \ird{-é}), the causative takes its
epenthetic form \ird{-naj-}; before the consonant-initial hortative \ird{-ká} it
remains \ird{-na-}:

\pex
\a
\begingl
\gla Trava bylam piaštnajím.//
\glb bread.\Dir{} child.\Ind{} eat-\Caus{}-\Imp{}//
\glft \trsl{Make the child eat the bread.} \lingcontext{causative; causee in \Ind{}}//
\endgl
\a
\begingl
\gla Jí trava bylam piaštnajím.//
\glb 2\textsc{sg}.\Trs{} bread.\Dir{} child.\Ind{} eat-\Caus{}-\Imp{}//
\glft \trsl{You make the child eat the bread.} \lingcontext{overt subject in \Trs{}}//
\endgl
\xe

\noindent Without an overt causee the causative is ungrammatical in directive clauses:

\pex
\rec{trava piaštnajím} \trsl{intended: `Make someone eat the bread.'} \lingcontext{causative without causee}
\xe

The requirement for overt beneficiaries and causees is categorical. The
directive subsystem resolves case pressure not by multiplying direct nominals
but by reserving the direct for the theme and assigning added non-theme
participants to the indirect. An applicative without a beneficiary or a
causative without a causee would leave an indirect role structurally implied but
unresolved. Since the language tolerates multiple indirect nominals, a directive
clause bearing both applicative and causative morphology is structurally
possible, even if stylistically heavy; the principle throughout is that added
participants enter the indirect while the theme retains the direct.

Antipassive morphology is excluded from all directive moods. The discourse
function of the antipassive --- demoting a patient from core status in order to
foreground the agent's activity --- is incompatible with the directive domain.
In the directive subsystem the null-subject orientation already removes the
agent-versus-patient tension that motivates antipassive use in the indicative,
and the single-direct principle already resolves any structural pressure by
assigning the direct to the theme and the transitive to the overt subject. The
antipassive-transitive distinction is therefore operative only in the
indicative.

The directive subsystem as a whole is characterized by null-subject orientation,
the absence of ordinary aspect morphology, a preference for only one core direct
nominal, marked overt speech-act participant subjects in the transitive case
where competition for the direct slot would arise, and the well-structured
integration of applicative and causative morphology in which beneficiaries and
causees enter the indirect. This pattern does not simply reproduce either half
of the indicative split. Rather, it regularizes its own case geometry in a way
that keeps the direct nominal central and prevents competition within the clause
core.


\subsection{Irrealis}

The irrealis (\ird{-ěv}) is used for counterfactual, hypothetical, or non-actual
situations. It belongs to the palatalizing, not the sibilant-replacing, domain:

\ex
\ird{piašt-} + \ird{-ěv} → \ird{piaščev}
\xe

In stems with the appropriate grade alternations, the front vowels /ɛ/ and /eː/
further raise to /i/ and /iː/ before this palatalizing suffix:

\ex
\ird{zděk-} `to bring' → \ird{zdíčev} \lingcontext{-e- raised to -í-; -k- palatalized to -č-}
\xe

The irrealis is compatible with all three converb moods and serves as the
main-clause mood in concessive constructions; both uses are treated in
Chapter~\ref{ch:complex}.

The negative counterpart of the irrealis is \ird{-ir}, a competing negative
Vocabulary Item at MoodP rather than a suppletive form in the strict sense. Like
the prohibitive \ird{-é} (negative imperative) and the negative hortative
\ird{-ku}, it is inherently negative and therefore does not co-occur with the
negative prefix \ird{zá-}. Phonologically, \ird{-ir} does not trigger the full
palatalizing pattern of \ird{-ěv}: stems in \ird{-k} show the restricted
alternation \ird{k $\rightarrow$ c}, but there is no broader \ird{-ěv}-type
palatalization and no associated vowel raising. Thus the same stem that yields
\ird{piaščev} in the affirmative irrealis yields \ird{piaštir} in the negative
irrealis, while \ird{zděk-} forms \ird{zděcir} rather than \ird{zdíčev}. The
formal analysis is in §\,\ref{sec:vm-moodp}.

\pex
\a \ird{piaštir} \trsl{would/might not eat; irrealis non-eating}
\a \ird{zděcir} \trsl{would not bring} \lingcontext{restricted \ird{k $\rightarrow$ c}; no vowel raising; cf.\ affirmative \ird{zdíčev}}
\xe


\subsection{Formal analysis: the IP domain}\label{sec:vm-ip}

\subsubsection{MoodP}
\label{sec:vm-moodp}

The mood suffixes occupy the projection immediately above AspP. In the
indicative, this projection is present but headed by a null morpheme (the mood
suffix has a null exponent). In the non-indicative moods (imperative, hortative,
irrealis), a lexical morpheme appears and simultaneously suppresses AspP (see
§\,\ref{sec:vm-mood-asp}).

The ordering of Mood above Asp is consistent with the cross-linguistic tendency
for tense and mood to dominate aspect (Cinque's (1999) functional hierarchy
places Mood$_\text{irrealis}$ above Asp$_\text{perfective/imperfective}$). The
mood paradigm consists of competing Vocabulary Items inserted at the MoodP head:

\begin{center}
\begin{tabular}{l l l}
\toprule
\textsc{mood} & \textsc{affirmative vi} & \textsc{negative vi} \\
\midrule
Indicative   & $\emptyset$      & ---\quad (NegP projected; \ird{zá-}) \\
Imperative   & \ird{-ím}        & \ird{-é} \hspace{3pt}(prohibitive) \\
Hortative    & \ird{-ká}        & \ird{-ku} \\
Irrealis     & \ird{-ěv}        & \ird{-ir} \\
\bottomrule
\end{tabular}
\end{center}

\medskip
\noindent Three of the eight cells above carry a lexically specified
[+\textsc{neg}] feature as part of their Vocabulary Item: the prohibitive
\ird{-é}, the negative hortative \ird{-ku}, and the negative irrealis \ird{-ir}.
This is a property of the individual VI, not a productive morphological
operation; there is no synchronic rule that derives the negative cell from the
affirmative. The affirmative/negative pairing is an artifact of their
paradigmatic organization at MoodP, not of compositional negation.

When a MoodP head bearing [+\textsc{neg}] is selected, NegP is absent
from the derivation entirely: the clause is assembled with the negative MoodP
merging directly above AspP, and no Neg\textsuperscript{0} head enters the
numeration. \ird{Zá-}, as the PF exponent of Neg\textsuperscript{0}, therefore
has no node at which to be inserted and does not surface. No independent
blocking principle is required; the absence of \ird{zá-} is a simple consequence
of NegP not being projected.

An important asymmetry follows. The indicative is the only mood that projects
NegP and uses \ird{zá-}: it is the only cell in the MoodP paradigm whose
negative value is expressed by a separate functional projection rather than by
an inherently negative VI. The three non-indicative moods each have their own
negative VI, rendering NegP superfluous in those environments.

The semantic character of MoodP [+\textsc{neg}] differs from that of NegP
[+\textsc{neg}]. NegP negation is propositional: it operates over a
truth-evaluable event description and inverts its truth conditions. MoodP
[+\textsc{neg}] is \term{illocutionary/modal-scoped}: it negates the force or
modality encoded by the MoodP head, not a proposition. Within this class, two
sub-types are distinguished. The prohibitive \ird{-é} and the negative hortative
\ird{-ku} are \term{directive-scoped}: the negation qualifies the illocutionary
force of a command or suggestion. A prohibitive such as \ird{piaštéb-é}
\trsl{Don't feed anyone!} does not have truth conditions in the ordinary sense;
its negation qualifies the directive illocution, not a predicate. The negative
irrealis \ird{-ir}, by contrast, is \term{modal/hypothetical-scoped}: the
negation denies the relevance or applicability of a hypothetical
event-description --- that is, it operates over the non-actual world introduced
by the irrealis, not over a directive force. \ird{Piaštir} \trsl{would not eat;
irrealis non-eating} asserts that the hypothetical eating does not obtain,
rather than prohibiting it. The two sub-types of MoodP [+\textsc{neg}] are
formally parallel in their blocking of NegP projection; they are semantically
distinct in scope (directive illocution vs.\ hypothetical modality) and
therefore appear in separate rows of the negation summary table
(Table~\ref{tab:neg-exponents}).

The negative irrealis \ird{-ir} shows only a restricted front-vocalic effect. In
stems ending in \ird{-k}, it triggers the alternation \ird{k $\rightarrow$ c},
but it does not otherwise trigger the broader verbal palatalization pattern
associated with \ird{-ěv}: there is no \ird{č}-output and no accompanying vowel
raising. This partial behavior sets it apart both from fully palatalizing
\ird{-ěv} and from wholly non-palatalizing \ird{-ic}. Accordingly, \ird{zděc-ir}
is the correct negative irrealis of \ird{zděk-} \trsl{bring}, not \ird{*zděkir}
and not \ird{*zdíčir}; stems such as \ird{piašt-} still surface simply as
\ird{piaštir}.


\subsubsection{Inner evidential (EvidP$_\text{inner}$)}
\label{sec:vm-evid}

An evidential suffix occupies the outermost position in the finite inflectional
spine, immediately inside the complement linker \ird{-e}. This
\term{inner evidential} (EvidP$_\text{inner}$) is distinct from the clausal
evidential and speech-act operators in the left periphery
(§\,\ref{sec:th-particle-framework}), which are Zone~I particles occupying
EvidP$_\text{outer}$ and SpeechActP in the clause spine.

The morphological evidence for the inner evidential is the position of the
indirect evidential suffix (tentatively \ird{-ejte}) in the verbal template: it
precedes the complement linker \ird{-e} and follows MoodP, giving the order:

\begin{center}
\textsc{stem} + \textsc{asp.align} + (\textsc{mood}) +
(\textsc{ind.evid}) + \ird{-e} + \ird{po}
\end{center}

\noindent The corresponding syntactic projection is therefore above MoodP and
constitutes the outermost layer of the finite inflectional spine. The boundary
with the clausal EvidP$_\text{outer}$ (a left-peripheral projection) must be
kept clear: the inner evidential marks the \emph{source of knowledge} as it
bears on the proposition expressed by the verbal complex; the outer EvidP
particles (Zone~I) operate at the clause-level and take scope over the full
proposition including its discourse-relational frame.

%% STUB: The inner evidential suffix is insufficiently described
%% (as noted in the main verb chapter). A full treatment must:
%% (a) establish the form of the suffix with certainty;
%% (b) determine whether it occupies a projection above or below
%%     MoodP (the current template places it above, i.e. outside);
%% (c) investigate its interaction with the outer EvidP particles
%%     (can they co-occur? What is their combined interpretation?);
%% (d) test whether the indirect evidential is restricted to
%%     matrix clauses or also appears freely in embedded clauses.



\section{Non-finite forms: converb morphology}\label{sec:converb-morphology}

% ============================================================
% STATUS: MORPHOLOGY-ONLY SECTION
% This section documents what the converb suffixes look like and how
% they attach. The syntax and semantics of converb constructions are
% treated in Chapter~\ref{ch:complex}.
% ============================================================

Converb suffixes replace the aspect-alignment slot entirely. A verb carrying a converb suffix accordingly bears no aspect, no mood, and no alignment. The morphological structure of a converb is:

\medskip
\begin{center}
(\textsc{lex-prefix}) + \Stem{} + (\Appl{}) + (\Caus{}) + \textsc{converb-suffix}
\end{center}
\medskip

The converb suffixes and the stem mutations they occasion are set out in the following table:

\begin{table}[ht]
	\footnotesize\sffamily
	\caption{Converb suffixes and their phonological effects.}\label{tab:converb-morph}
	\medskip
	\begin{tabularx}{\textwidth}{L L L L}
		\toprule
		{\sc converb} & {\sc suffix(es)} & {\sc stem mutation} & {\sc notes} \\
		\midrule
		Sequential    & \ird{-u}, \ird{-ení} & none & \ird{-ní} after applicative/causative \\
		Simultaneous/manner & \ird{-ěc} & palatalization & front-vocalic trigger; may coincide with older sibilantized outputs in some stems \\
		Temporal (6)  & \ird{-ezak}, \ird{-eškady}, \ird{-eeby}, \ird{-ešhoume}, \ird{-eebymy}, \ird{-esny} & none & all add \ird{-e-} linker \\
		Concessive    & \ird{-ouc} & none & \\
		Conditional   & \ird{-ic}, \ird{-emy}, \ird{-ezmy}, \ird{-ena} & none for \ird{-ic}; see note & \ird{-ic} does not trigger palatalization \\
		Purposive     & \ird{-it} & irregular; mainly Class I & \\
		Contrastive   & \ird{-ema} & none & \\
		\bottomrule
	\end{tabularx}
\end{table}

The simultaneous/manner converb \ird{-ěc} and the irrealis mood \ird{-ěv} belong to the palatalizing domain, while the antipassive \ird{-aš-} is the clearest instance of true sibilant replacement. The conditional \ird{-ic} does not palatalize despite its initial front vowel; the historical reasons for this exception are discussed in Section~\ref{sec:cond-nopala} of Chapter~\ref{ch:complex}. Stronger irregularity is concentrated instead in a few high-frequency linker and quotative forms, where analogical leveling has obscured the older outputs.


\subsection{Formal analysis: the syntactic status of converbs}
\label{sec:vm-converbs}

\subsubsection{Converbs as aspect-deficient vP/VoiceP projections}
\label{sec:vm-converb-aspless}

Converbs replace the Asp.Align slot entirely. A verb bearing a
converb suffix therefore projects a vP/VoiceP but no AspP. The
converb suffix itself can be analyzed as the Vocabulary Item
inserted at a \term{defective AspP} node: a position in the
functional hierarchy that would normally host an Asp.Align morpheme
but instead hosts a converb-class morpheme that lacks aspect and
alignment features.

Alternatively, converbs project no AspP at all: the converb suffix
is the head of a Voice or vP projection that merges directly with
the higher adverbial or complementizer position without passing
through AspP. The absence of the ergative case marker \ird{-n-}
inside converb clauses (§\,\ref{sec:converb-chains}) is predicted
by either account: if AspP is absent, the Asp head that would have
triggered ergative-absolutive alignment is not present to assign
the transitive case to the agent.

\subsubsection{Case alignment within converb clauses}
\label{sec:vm-converb-case}

The nominative-accusative alignment that is invariant within
converb clauses (absolutive agent, indirect patient) is consistent
with the absence of an Asp head: without the perfective/
retrospective Asp head to assign ergative case to the agent, the
agent defaults to absolutive (nominative), and the patient receives
oblique/indirect case. This is in effect the same pattern as the
nominative-aspect main clause, with the difference that in the main
clause it is the Asp head (continuous/progressive/contemplative)
that assigns this pattern, whereas in the converb clause it is the
absence of any Asp head.

\subsubsection{Converb clauses as adjuncts and embedded topics}
\label{sec:vm-converb-adjunct}

In the clause-combining hierarchy, converb clauses function as
adverbial adjuncts to the matrix clause. They are merged as
adjuncts to the matrix IP or vP, not as arguments in Spec or
Complement positions. This explains their preverbal (and pre-matrix)
position: as adjuncts to IP, they precede all material within IP
and surface before the matrix arguments in the linear string,
consistent with the head-final adjunction pattern.

Switch-reference in converb chains (§\,\ref{sec:switch-reference}) is a
discourse-pragmatic phenomenon: different-subject readings are licensed by
the discursive contrast of two distinct referents in topic position, not by
any morphological clitic on the converb. When both clauses bear full nominal
subjects, those DPs suffice for disambiguation; when a subject is absent,
discourse recoverability determines the reading.


\subsection{The nominalized verb forms}\label{sec:nominalized}

% ============================================================
% STATUS: STUB — to be developed
% The 2020 archive describes the following nominalized forms:
%   - Generic nominal in -ou (from inflected stem)
%   - Participant nominals (agent nominal -kou, patient nominal)
%   - Event nominals
% These interact with definiteness, argument structure, and focus
% (the "nominalization for indefiniteness" construction described in Ch. 4).
% ============================================================

Iridian possesses several forms that nominalize the verb. One is the lexical or resultant nominal in \ird{-ou}, often semi-productive and lexeme-specific, which yields nouns such as \ird{piaštou} \trsl{food}. Another is the inherited gerundal strategy, which remains the productive event-nominal formation. Participant nominals likewise employ \ird{-ou}, but they are built from stems whose morphology already fixes which participant occupies the absolutive pivot.

Antipassive stems yield agent-oriented participant nominals, unmarked transitive stems patient-oriented ones, and applicative stems nominals referring to the promoted participant.

All five indicative aspects are used productively in these participant nominals: perfective, retrospective, continuous, progressive, and contemplative. The same voice-sensitive stem that determines the nominal's orientation also carries the aspectual morphology, so the nominal may refer to a completed participant, an experientially or result-relevant past participant, an ongoing or habitual participant, a participant in an event now in progress, or a participant associated with an intended or contemplated event. In this formation \ird{-ek + -ou} contracts to \ird{-ikou}, retrospective \ird{-ac} simply yields \ird{-acou}, \ird{-a + -ou} to \ird{-ova}, and \ird{-ime + -ou} to \ird{-imou}; the contemplative stem takes \ird{-ou} without further special contraction. Thus \ird{piaščikou} \trsl{the one who ate}, \ird{piaščacou} \trsl{the one who has eaten}, \ird{piaščova} \trsl{the one who eats / is eating}, \ird{piaščimou} \trsl{the one who is in the act of eating}, \ird{piaščachou} \trsl{the one who is about to eat}, \ird{piaštnikou} \trsl{that which was eaten}, \ird{piaštnacou} \trsl{that which has been eaten}, and \ird{piaštébikou} \trsl{the one who was fed}. Their syntax, definiteness properties, and argument structure are treated in Section~\ref{sec:event-nominals} of Chapter~\ref{ch:nouns}.

\subsubsection{Predicative use: participant nominals in identificational equatives}
\label{sec:nmlz-predicative}

Participant nominals may serve as the predicative complement of a zero-copula equative clause, with the topic DP identified as the entity satisfying the nominalized description. This construction has the schematic form:

\begin{center}
[\textit{topic-DP}] \ \ [\textit{(internal-arg.\Ind{})} \ \ \textit{participant-nominal}]
\end{center}

\noindent Compare the base transitive clause with its patient-nominal equative counterpart:

\pex
\a
\begingl
\gla Lobek Janku piaštnik.//
\glb apple.\Dir{} Janek.\Trs{} eat-\textsc{trs}-\Pf{}//
\glft \trsl{Janek ate the apple.} \lingcontext{standard ergative clause}//
\endgl
\a
\begingl
\gla Lobek Jancí piaštnikou.//
\glb apple.\Dir{} Janek.\Ind{} eat-\textsc{trs}-\Pf{}-\textsc{nmlz}//
\glft \trsl{The apple is what Janek ate.} \lingcontext{patient nominal as equative predicate}//
\endgl
\xe

\noindent And with an agent-nominal equative:

\ex
\begingl
\gla Janek lubcí piaščikou.//
\glb Janek.\Dir{} apple.\Ind{} eat-\Antip{}-\Pf{}-\textsc{nmlz}//
\glft \trsl{Janek is the one who ate the apple.} \lingcontext{agent nominal as equative predicate}//
\endgl
\xe

Two features of the case marking are notable. First, the agent of the base
transitive clause, which bears the transitive case in the finite verb
construction (\ird{Janku}, \Trs{}), shifts to the indirect case
(\ird{Jancí}, \Ind{}) when it appears as an internal argument of the
patient nominal. Second, the patient argument removed from core argument structure by the antipassive
in a finite clause may reappear in the indirect case as an internal argument
of the agent nominal (\ird{lubcí} in the second example above). In both
cases all internal arguments of the nominalized event appear in the indirect
case, since the absolutive/direct slot of the equative clause is occupied
by the external topic DP and no verbal case-assigning machinery operates
within the nominalized predicate.

The semantic effect of predicating a participant nominal is
\term{identificational}: the topic DP is being identified as the unique
referent satisfying the description encoded in the nominal. This contrasts
with simple property-ascription equatives (\ird{Marek doktor}, \trsl{Marek
is a doctor}) and with the plain verbal assertion (\ird{Lobek Janku
piaštnik}, \trsl{Janek ate the apple}). Negation of identificational
equatives with \ird{staly} and its semantic contrast with the verbal
\ird{zá-} are treated in §\,\ref{sec:ident-equative} and
§\,\ref{sec:clausal-negation} of Chapter~\ref{ch:simple}.


\section{The conjunctive / quotative linker}\label{sec:conjunctive-linker}

Iridian has a verbal linking form in \ird{-e}, used in clause linking and quotative constructions. Synchronically this \ird{-e} is a front-vocalic palatalizing suffix added to the fully inflected verb. It therefore belongs with the palatalization system described in Chapter~\ref{ch:phon}, not with sibilant replacement.

\pex
\a \ird{piašček} $\rightarrow$ \ird{piaščice}
\a \ird{piaščac} $\rightarrow$ \ird{piaščace}
\xe

The exact output depends on the segment immediately preceding the linker. Velars typically yield \orth{c/ce}; other endings show the regular soft reflexes of the final consonant or syllable in standard spelling, without introducing separate haček letters for soft \orth{t}, \orth{d}, or \orth{l}. Some older conjunctive formations are analogically levelled and remain morphologized, but the synchronic generalization is clear: \ird{-e} selects a palatalized form of the preceding material. The formal syntactic analysis of the linker is developed in the following subsection.


\subsection{Formal analysis: the conjunctive linker \ird{-e}}
\label{sec:vm-tension-linker}

The linker \ird{-e} appears at the right edge of embedded finite
clauses before \ird{po}/\ird{poby} and also in the conjunctive
clause-chaining function described above.
Its dual function raises the question of whether it is a single
morpheme with a unified syntactic status or two homophonous
morphemes. A full resolution requires integrating the phonological
behavior of \ird{-e} with its morphosyntactic distribution and its
position in the verbal template.

\subsubsection{Phonological characterization}

\paragraph{Trigger class and palatalization.}
Synchronically, \ird{-e} is a bound suffix that belongs to the
productive front-vocalic palatalization system described in
§\,\ref{sec:front-palatalization}. Its trigger class is the same as
that of \ird{-ě}, \ird{-í}, and the irrealis suffix \ird{-ěv}: all
are front-vocalic and all soften the segment to which they attach.
The conditional converb \ird{-ic}, though it begins with a front
vowel, is the systematic exception: it does not trigger
palatalization, because it was grammaticalized during a period when
its initial vowel lacked palatalizing force
(§\,\ref{sec:cond-nopala}). This exception is historically
motivated and synchronically opaque; it does not disturb the
generalization that \ird{-e} is a palatalizing trigger in the
productive system.

\paragraph{Segment-specific outputs.}
The palatalizing effect of \ird{-e} targets the final segment of the
inflected form to which it attaches, not the thematic consonant of
the bare stem. The outcomes depend on the place and manner of the
final consonant:

\begin{itemize}
  \item \textit{Velars} (/k/, /g/): the velar yields the palatalized
    affricate series. The most common output is \orth{c} for /k/,
    as in \ird{piašček} $\rightarrow$ \ird{piaščice} (the final
    velar of the perfective suffix \ird{-ek/-ik} palatalizes to
    \orth{c} before \ird{-e}).

  \item \textit{Coronals} (/t/, /d/, /n/, /l/, /r/): the regular
    soft reflexes apply. Standard spelling does not introduce haček
    diacritics for softened \orth{t}, \orth{d}, or \orth{l} in this
    context; the softness is indicated by the following vowel graph
    or is left orthographically unmarked for the consonant itself.
    Thus \ird{piaščac} (the retrospective, whose suffix now ends in
    \ird{-c}) yields \ird{piaščace} under \ird{-e}: the final
    affricate simply surfaces with the linker vowel as \ird{-ce}.

  \item \textit{Labials} (/p/, /b/, /m/, /v/) and liquids (/r/):
    secondary palatalization (a palatal off-glide) is added, giving
    the same soft-series output as before nominal \ird{-í}
    (§\,\ref{sec:palatalization}). These outcomes are less
    phonologically salient than velar or coronal softening and are
    typically not marked distinctively in orthography.
\end{itemize}

\noindent The key generalization is that \ird{-e} selects the soft
allomorph of whichever ending it attaches to, following the same
mechanism as other palatalizing suffixes. The output is therefore
predictable from the palatalization system already operative in
nominal and verbal morphology, without requiring a separate
phonological rule.

\paragraph{Interaction with vowel raising.}
When the final syllable of the inflected form contains a short
front vowel (\ird{e} or \ird{ě}), that vowel is subject to raising
to \ird{i}/\ird{í} in the syllable immediately preceding the
palatalizing trigger (§\,\ref{sec:vowel-raising}). In the derivation
\ird{piašček} + \ird{-e}: the thematic \ird{-e-} of the perfective
suffix raises to \ird{-i-} before the palatalized output, yielding
\ird{piaščice} rather than *\ird{piaščece}. This raising is the same
rule that operates in \ird{zděk-} + \ird{-ěv} $\rightarrow$
\ird{zdíčev} (the \ird{-e-} of the root raises before the
front-vocalic verbal suffix). The interaction of raising with
palatalization is thus not specific to \ird{-e}: it is a general
consequence of attaching any front-vocalic suffix to a stem whose
final syllable contains a short front vowel.

\paragraph{Analogically levelled forms.}
Some high-frequency conjunctive forms are morphologized: they
preserve older palatalizations or otherwise deviate from the
synchronic rule. These are most concentrated in the finite forms
of verbs with irregular thematic consonants or with stems that
underwent earlier lenition or vowel contraction. The synchronic
system tolerates these as exceptions without compromising the
productivity of the rule. Their existence does, however,
complicate any derivational analysis that requires strict
compositionality: the linker form of a high-frequency verb may
need to be stored as a listed item alongside its productive
derivation.

\paragraph{Prosodic status.}
The linker \ird{-e} is a bound morpheme: it does not constitute an
independent prosodic word and does not shift the stress of the
verbal complex. Stress in Iridian is fixed to the lexical stem
and is not influenced by suffixation in the outer inflectional
layers. The prosodic invariance of \ird{-e} is consistent with its
status as a suffix rather than a clitic or particle. If it were a
prosodic word or particle, the stress pattern of \ird{piaščice}
might be expected to differ from that of \ird{piaščik}; there is no
evidence that it does.

\subsubsection{Morphosyntactic distribution}

\paragraph{Template position.}
Within the verbal complex, \ird{-e} attaches outside all finite
inflectional morphology. The full template is:

\begin{center}
\Stem{} + \Asp{}.\textsc{align} + (\textsc{mood}) +
(\textsc{evid}) + \ird{-e}
\end{center}

\noindent The linker is therefore exterior to EvidP$_\text{inner}$, the
outermost projection in the finite inflectional spine
(§\,\ref{sec:vm-evid}). This position is significant: if \ird{-e} were an
agreement or mood morpheme, it would be expected inside the inflectional spine.
Its consistent extra-clausal position supports an analysis in which it marks the
edge of the finite complex as it interfaces with a higher projection --- a
CP-boundary marker, not an inflectional head.

\paragraph{Attested environments.}
The linker appears in exactly three environments:

\begin{enumerate}
  \item \textit{The finite \ird{po}-complement}
    (§\,\ref{sec:vm-po-comp}): the linker marks the right edge of
    the embedded finite TP before the complement C head \ird{po} or
    \ird{poby}. It is obligatory here; an embedded finite clause
    without the linker is ungrammatical before \ird{po}/\ird{poby}.

  \item \textit{The conjunctive chain} (§\,\ref{sec:conjunctive-linker}):
    the linker marks non-final finite clauses in a clause chain
    where the chained clauses are interpreted as coordinated or
    sequentially related propositions. This use resembles the
    clause-linking \textit{-ko} of Korean or the converbal
    linker of other verb-final languages, but unlike those, the
    Iridian form attaches to a fully finite, fully inflected verb.

  \item \textit{The matrixless \ird{po}-construction}
    (§\,\ref{sec:complementation}): even when the matrix predicate
    is absent and \ird{po} functions as a discourse-level
    attribution marker, the linker \ird{-e} is retained on the
    finite verb. This shows that \ird{-e} is not licensed by
    the presence of a lexical matrix predicate but by the presence
    of the C-domain element \ird{po} itself.
\end{enumerate}

\paragraph{Non-environments.}
The linker is absent in two structurally revealing contexts. First,
independent finite clauses lack \ird{-e}: when no C head embeds the
TP, the linker is not licensed. Second, converb forms --- the
non-finite clause-linking morphology of simultaneous, sequential,
conditional, and other converbs --- do not take \ird{-e}. This
follows naturally: converbs are already non-finite (they lack
Asp.Align and Mood morphology), so they do not participate in the
domain of the finite-clause boundary that \ird{-e} marks.

\subsubsection{Competing analyses}

Three main analyses are available; they diverge principally on
whether \ird{-e} has a unified syntactic identity across its two
main environments (complement and conjunctive chain).

\begin{description}[leftmargin=2em, labelwidth=!, labelindent=0em]
  \item[\upshape Unified CP-edge analysis.] Under this analysis,
    \ird{-e} is a single morpheme throughout: the phonological
    exponent of a \textsc{[+fin, +sub]} feature complex that is
    valued on the embedded TP when it enters the scope of a C
    head. In the complement construction, C is overt (\ird{po} or
    \ird{poby}). In the conjunctive chain, C is a null clause-juncture
    head (following the general principle that coordination in
    head-final languages may involve a null C$^0$ that takes a
    finite TP as its specifier or complement). The linker is the
    spell-out of the embedded T's contact with this C head, and its
    front-vocalic palatalizing nature is the PF signature of
    \textsc{[+fin, +sub]} in Iridian. The absence of \ird{-e} in
    independent clauses follows without stipulation: no C head is
    present to value the feature on T.

    The central challenge for the unified analysis is the
    distributional asymmetry: in the complement construction,
    \ird{-e} co-occurs with an overt C head (\ird{po}/\ird{poby});
    in the conjunctive construction, no overt C appears. The unified
    analysis must postulate a null conjunctive C$^0$ whose presence
    is evidenced only by the linker itself, which risks circularity
    unless independent evidence for this head can be adduced.

  \item[\upshape PF boundary-marker analysis.] Under this analysis,
    \ird{-e} is not a syntactic morpheme but a phonological
    convention: a boundary signal inserted at any right edge of a
    finite clause when that clause is followed by further clausal
    material. Its front-vocalic palatalizing character is then a
    phonological property of the boundary diacritic rather than a
    reflex of any syntactic feature. This account handles the
    conjunctive environment straightforwardly (every non-final
    finite clause is followed by another clause) and is empirically
    efficient.

    The drawback is that a PF boundary marker need not itself be a
    palatalizing morpheme: phonological boundary tones and
    juncture markers in other languages do not normally trigger
    segmental mutation. The palatalization caused by \ird{-e} is
    identical in type to the palatalization triggered by inflectional
    suffixes (\ird{-ě}, \ird{-í}, \ird{-ěv}), suggesting that
    \ird{-e} is a genuine suffix, not a boundary diacritic. The PF
    account therefore purchases distributional coverage at the cost
    of explaining a morphologically specific segmental process by
    appeal to an otherwise unmotivated phonological rule.

  \item[\upshape Two-homophone analysis.] Under this analysis,
    \ird{-e}$_1$ (the complement linker) and \ird{-e}$_2$ (the
    conjunctive linker) are distinct morphemes that happen to be
    phonologically identical and share the palatalizing trigger
    class. \ird{-e}$_1$ is the exponent of the TP-CP boundary
    in complement clauses; \ird{-e}$_2$ is a clause-chaining
    morpheme lexically specified for the conjunctive domain. Their
    coincidence in form is not synchronically active but is a
    historical convergence (or retention from a common source).

    This analysis avoids the null-C problem but is the most
    stipulative: it requires two separate lexical entries sharing
    phonological form and the same segmental-mutation property,
    with their distribution jointly covering the environments
    of both complement embedding and clause chaining. The burden
    of proof falls on demonstrating that the two environments
    are genuinely syntactically disjoint in a way that the
    unified analysis cannot accommodate.
\end{description}

\subsubsection{Preferred account and outstanding questions}

The unified CP-edge analysis is preferred here on parsimony grounds.
Its core claim is that \ird{-e} spells out the feature
\textsc{[+fin, +sub]} --- roughly, ``finite clause under
subordination'' --- on the embedded TP's outermost layer, and that
this feature is valued whenever a C head (overt or null) selects the
finite TP. The palatalizing nature of \ird{-e} is then the PF
realization of this morphosyntactic feature, fully parallel to the
palatalizing effect of other front-vocalic suffixes in the
productive system.

The position of \ird{-e} outside EvidP$_\text{inner}$ in the template
supports this analysis directly: the linker is the last element
added to the verbal complex before the clause exits the finite
inflectional domain, which is exactly where a TP-edge or C-interface
marker would be expected.

Two empirical questions remain open:

\begin{enumerate}
  \item \textit{The null conjunctive C$^0$.} The postulated null C
    head in the conjunctive construction should in principle be
    detectable through its interaction with the argument structure of
    the chained clauses: if the conjunctive C$^0$ takes the finite
    TP as its specifier and itself projects a CP, then the chained
    clause is a CP adjunct. This predicts that binding relations and
    scope interactions across conjunctive-chain clauses should pattern
    with embedding rather than flat coordination. Corpus investigation
    of binding and scope in conjunctive chains would provide a
    diagnostic.

  \item \textit{Interaction with the evidential.} The indirect
    evidential \ird{-ejte} appears inside \ird{-e} in the template
    (\ird{-ejte-e}), confirming that the evidential is a finite
    inflectional morpheme whereas the linker is not. But it is not
    yet established whether the palatalizing context of \ird{-e}
    applies to the final segment of \ird{-ejte} or directly to the
    verb stem. If \ird{-ejte} itself ends in a front vowel, the
    palatalizing context would overlap with it, potentially yielding
    an opaque surface form. Documenting the attested surface outputs
    of \ird{-ejte + -e} sequences across all five aspect forms would
    clarify whether the palatalization is invariant (expected under
    the unified analysis) or subject to additional sandhi conditions.
\end{enumerate}

\noindent The front-vocalic palatalizing nature of \ird{-e} is the
strongest evidence against treating it as a mere PF convention, and
the strongest argument for its status as a genuine morphosyntactic
suffix --- one that happens to appear at the clause's structural
edge, but whose phonological behavior is fully integrated with the
inner inflectional system of Iridian.


\section{Stative verbs}\label{sec:stative}

A class of stative verbs describes states rather than events. These verbs are
the language's primary means of expressing properties predicatively: where
English often uses an adjective, Iridian normally uses a stative verb. Stative
verbs are true verbs and accordingly require no copula to function as
predicates.

\subsection{Property roots and the verbalising suffix \ird{-n-}}

A small closed class of property roots each denotes an abstract property as a
nominal. The verbalising suffix \ird{-n-} is added to the root to yield a
stative verb stem, which then conjugates as a regular intransitive verb:

\begin{table}[ht]
	\footnotesize\sffamily
	\caption{Property roots and derived stative verbs.}\label{tab:statroots}
	\medskip
	\begin{tabularx}{\textwidth}{l L l L}
		\toprule
		{\sc root} & {\sc gloss} & {\sc stative stem} & {\sc gloss} \\
		\midrule
		\ird{hrad} & goodness  & \ird{hradn-}  & to be good \\
		\ird{mroc} & beauty    & \ird{mrocn-}  & to be beautiful \\
		\ird{trpa} & badness   & \ird{trpan-}  & to be bad \\
		\ird{svor} & bigness   & \ird{svorn-}  & to be big \\
		\bottomrule
	\end{tabularx}
\end{table}

The derived stative verb conjugates across the full five-way aspect system.
The continuous is closest to an equative property predication; the other
aspects carry characteristically dynamic or inchoative readings. The full
paradigm of \ird{mrocn-} \trsl{be beautiful} illustrates the system:

\pex
\a
\begingl
\gla Nužot mrocna.//
\glb city.\Dir{}.\Def{} be.beautiful-\Cont{}//
\glft \trsl{The city is beautiful.} \lingcontext{continuous: current state}//
\endgl
\a
\begingl
\gla Nužot mrocnime.//
\glb city.\Dir{}.\Def{} be.beautiful-\Prog{}//
\glft \trsl{The city is becoming beautiful.} \lingcontext{progressive: entering the state}//
\endgl
\a
\begingl
\gla Nužot mrocnik.//
\glb city.\Dir{}.\Def{} be.beautiful-\Pf{}//
\glft \trsl{The city became beautiful.} \lingcontext{perfective: ingressive}//
\endgl
\a
\begingl
\gla Nužot mrocnac.//
\glb city.\Dir{}.\Def{} be.beautiful-\Ret{}//
\glft \trsl{The city was beautiful.} \lingcontext{retrospective: past state}//
\endgl
\a
\begingl
\gla Nužot mrocnach.//
\glb city.\Dir{}.\Def{} be.beautiful-\textsc{contempl}//
\glft \trsl{The city will be beautiful.} \lingcontext{contemplative: expected state}//
\endgl
\xe

In bounded or developmental aspects --- the progressive, perfective, and
retrospective --- stative verbs are normally interpreted ingressively unless
context strongly resists it. The perfective thus yields \trsl{became
beautiful}, not simply \trsl{was beautiful}; that reading belongs to the
retrospective, which describes a past state conceived as a complete episode, a
state that no longer obtains, or one that is no longer at issue in the current
discourse. The continuous, by contrast, gives plain property predication
without any implication of change, and is the unmarked form for simple
description.

Stative verbs in the perfective take the raised allomorph \ird{-ik} rather
than \ird{-ek}; the verbalising suffix \ird{-n-} conditions this allomorphy in
the same manner as a stem-final \ird{e} or \ird{ě} (§\ref{sec:aspect-pf}).

\subsection{Attributive and adverbial derivatives}\label{sec:stative-attr}

Unlike regular verbs, stative verbs derived from property roots possess a
dedicated attributive form built by adding \ird{-á} directly to the verbalised
root. This form is tenseless, carries no aspectual or alignment marking, and
is used exclusively as a direct modifier within the noun phrase:

\pex
\a \ird{mrocná nuž} \trsl{a beautiful city}
\a \ird{hradná byl} \trsl{a good child}
\a \ird{svorná dum} \trsl{a big house}
\xe

The \ird{-á} form cannot be used predicatively. For predication, the full
stative verb form is required:

\pex
\a \ird{Nužot mrocna.} \trsl{The city is beautiful.} \lingcontext{predicative: stative verb required}
\a \ird{*Nužot mrocná.} \lingcontext{ungrammatical: attributive form in predicative position}
\xe

\noindent The \ird{-á} attributive functions as the adnominal counterpart of
the continuous stative, expressing the property directly without tense or
aspect. Its position within the noun phrase relative to determiners and other
modifiers is treated in Chapter~\ref{ch:nouns}.

Alongside the attributive in \ird{-á}, property-root-based stative stems
productively form a dedicated class of derived manner adverbs in \ird{-ě}
\label{sec:stative-adv}. This formation is derivational rather than
inflectional: it yields a lexical adverbial adjunct, not a further verbal
form, and it is distinct both from the adnominal attributive and from the
converb system discussed in Chapter~\ref{ch:complex}. Schematically, the
pattern is as follows:

\begin{center}
\term{property root $\rightarrow$ stative stem in \ird{-n-} $\rightarrow$ attributive in \ird{-á} / manner adverb in \ird{-ě}}
\end{center}

\noindent A representative paradigm is the following:

\pex
\a \ird{skor} \trsl{quickness}
\a \ird{skorn-} \trsl{be quick}
\a \ird{skorná} \trsl{quick, fast} \lingcontext{attributive}
\a \ird{skorně} \trsl{quickly}
\xe

The resulting system preserves a clear three-way distinction: the finite
stative predicate expresses property predication, the \ird{-á} form modifies a
noun within the DP, and the \ird{-ě} form modifies the event as a manner
adverb. Thus \ird{skorn-} means \trsl{be quick}, \ird{skorná} means
\trsl{quick} as a noun modifier, and \ird{skorně} means \trsl{quickly} as a
clausal adjunct. The syntactic position of these derived manner adverbs in the
preverbal field is discussed in Chapter~\ref{ch:simple}; compare also the
simultaneous / manner converb in \ird{-ěc} in Chapter~\ref{ch:complex}, which
remains a non-finite verbal form rather than a lexical adverb.

\paragraph{Degree prefixes.}
Scalar degree on the stative class is expressed by a closed set of prefixes
that attach to the stative stem before any derivational material, including
the verbalising \ird{-n-} and the attributive \ird{-á}. The five prefixes are
mutually exclusive and occupy a single slot:

\begin{table}[ht]
  \footnotesize\sffamily
  \caption{Degree prefixes on stative verbs.}\label{tab:degpfx}
  \medskip
  \begin{tabularx}{\textwidth}{l l L L}
    \toprule
    {\sc prefix} & {\sc value} & {\sc finite} & {\sc attributive} \\
    \midrule
    \ird{po-}  & slight; somewhat      & \ird{pomrocna}  & \ird{pomrocná}  \\
    \ird{do-}  & moderate; quite       & \ird{domrocna}  & \ird{domrocná}  \\
    \ird{pre-} & strong; very          & \ird{premrocna} & \ird{premrocná} \\
    \ird{od-}  & elevated; exceedingly & \ird{odmrocna}  & \ird{odmrocná}  \\
    \ird{nad-} & excessive; too        & \ird{nadmrocna} & \ird{nadmrocná} \\
    \bottomrule
  \end{tabularx}
\end{table}

\noindent The prefixes are fully productive across the stative class. Because
they precede the stem complex, they appear before \ird{-n-} in finite forms
and before the root in attributives.

\paragraph{The emphatic suffix \ird{-ěsk-}.}
The suffix \ird{-ěsk-} is distinct from the degree prefix system and does not
occupy the same grammatical slot. It is restricted to the \ird{-á} attributive
and carries a marked, heightened, or literary flavour rather than a simple
scalar value --- closer in function to an emphatic particle than to a degree
adverb. The contrast with a prefixed form is one of register and framing
rather than degree alone: \ird{premrocná} states a degree as a plain fact;
\ird{mrocněská} colours the attribution, implying that the beauty is
remarkable, conspicuous, or worth dwelling on. In ordinary speech
\ird{-ěsk-} attributives may sound affected; in verse, formal oratory, and
deliberate rhetorical prose they are appropriate and expected.

The suffix is not fully productive. It is most natural with roots that carry
inherent aesthetic or evaluative content:

\pex
\a \ird{mrocněská} \trsl{strikingly beautiful; beautiful indeed} \lingcontext{natural in literary register}
\a \ird{hradněská} \trsl{truly good; good in the full sense} \lingcontext{natural in literary register}
\xe

\noindent With dimensional or manner-denoting roots the formation is possible
but tends to sound stilted or unidiomatic outside verse:

\pex
\a \ird{svorněská} \trsl{(strikingly big)} \lingcontext{odd in prose; possible in elevated style}
\a \ird{skorněská} \trsl{(remarkably quick)} \lingcontext{marginal; feels forced outside verse}
\xe

\noindent \ird{trpěská} occupies a middle ground: the negative valence of
\ird{trpa} (\trsl{badness}) makes its \ird{-ěsk-} form feel ironic or
rhetorical rather than sincerely emphatic, and its use is consequently rare
and stylistically loaded.

\ird{-ěsk-} may co-occur with a degree prefix. The two are additive rather
than competing: the prefix supplies the scalar value and \ird{-ěsk-} supplies
the rhetorical register.

\pex
\a \ird{premrocněská} \trsl{very beautiful} \lingcontext{emphatic; degree stated and dwelt upon}
\a \ird{odmrocněská} \trsl{exceedingly beautiful} \lingcontext{elevated register; characteristic of panegyric}
\a \ird{nadmrocněská} \trsl{excessively, ruinously beautiful} \lingcontext{ironic or elegiac; beauty to a fault}
\xe

\noindent Such combinations are licensed but marked; they are characteristic
of formal verse and deliberate oratorical prose. In ordinary speech, a
prefixed form without \ird{-ěsk-} is the normal choice.

\subsection{Nominalisation of stative verbs}

Finite stative verbs are nominalised by the same \ird{-ou} mechanism that
applies to all verbs (§\ref{sec:nominalized}). The nominalised form yields a
substantive or classifying reading rather than a simple description: the
speaker is not merely ascribing a property but identifying a referent by that
property or setting it in contrast to others.

\pex
\a \ird{Nužot mrocna.} \trsl{The city is beautiful.} \lingcontext{simple property predication}
\a \ird{Nužot mrocnova.} \trsl{The city is the beautiful one.} \lingcontext{substantive: identifying or contrastive}
\xe

\pex
\a \ird{Bylot svornik.} \trsl{The child became big (i.e.\ grew up).} \lingcontext{ingressive perfective}
\a \ird{Bylot svornikou.} \trsl{The child is the one who grew up.} \lingcontext{nominalised perfective}
\xe

When a nominalised stative verb is used adnominally, it takes the standard
\ird{-ovní} form:

\ex
\begingl
\gla Svornikovní bylot Marek.//
\glb be.big-\Pf{}-\textsc{nmlz}-\textsc{attr} child.\Dir{}.\Def{} Marek.\Dir{}//
\glft \trsl{Marek is the child who grew up.}//
\endgl
\xe

The continuous is an exception: the continuous stative nominalised in
\ird{-ova} cannot be used adnominally with \ird{-ovní}. In adnominal position
the \ird{-á} attributive must be used instead:

\pex
\a \ird{mrocná nuž} \trsl{a beautiful city} \lingcontext{correct: \ird{-á} attributive}
\a \ird{*mrocnavní nuž} \lingcontext{ungrammatical: \ird{-ovní} built on continuous stem}
\xe

\subsection{State types and aspectual interpretation}\label{sec:stative-stability}

Not all states predicated by stative verbs are equally permanent. Some roots
denote enduring or characteristic properties --- canonical evaluative qualities
such as \ird{hrad} \trsl{goodness} or \ird{mroc} \trsl{beauty}, and
dimensional properties such as \ird{svor} \trsl{bigness} --- while others
denote contingent or temporary conditions: \ird{ptakn-} \trsl{be tired},
\ird{smrn-} \trsl{be sick}, \ird{hvodn-} \trsl{be open}, \ird{žbirn-}
\trsl{be wet}. This distinction is not encoded grammatically, but it shapes
the naturalness and interpretation of aspectual forms.

For a stable property, the continuous yields plain predication and the
progressive is somewhat marked, carrying a clearly developmental reading that
implies a process of change (\trsl{getting better}, \trsl{getting bigger}).
For a temporary condition, the continuous describes a contingent present state
(\trsl{is tired right now}) and the progressive more readily receives an
ingressive reading (\trsl{is starting to feel tired}, \trsl{is getting sick}).

Three semantic classes within the stative system may accordingly be
distinguished: \term{canonical quality statives}, which denote evaluative
properties of things (\trsl{be good}, \trsl{be beautiful}, \trsl{be bad});
\term{dimensional statives}, which denote extent or scale properties
(\trsl{be big}); and \term{experiential and bodily-condition statives}, which
denote states of persons (\trsl{be ill}, \trsl{be hungry}, \trsl{be ashamed}).
The last class most readily admits progressive and retrospective readings, since
temporary and personally felt states support aspectual dynamics more naturally
than stable properties do.

\subsection{Valency operations with stative verbs}\label{sec:stative-valency}

Stative verbs admit the causative \ird{-na-}, following the same alignment
conditions as other verbs (§\ref{sec:aspect}): nominative alignment in the
progressive and contemplative, ergative alignment in the perfective and
retrospective. Because stative stems end in \ird{-n-}, the causative suffix
surfaces as \ird{-an-} rather than \ird{-na-}, the order being reversed to
avoid a geminate cluster at the morpheme boundary:

\pex
\a
\begingl
\gla Marek problemam svornanime.//
\glb Marek.\Dir{} problem.\Ind{} be.big-\Caus{}-\Prog{}//
\glft \trsl{Marek is making the problem bigger.} \lingcontext{causative progressive; nominative alignment}//
\endgl
\a
\begingl
\gla Problem Marku svornanik.//
\glb problem.\Dir{} Marek.\Trs{} be.big-\Caus{}-\textsc{trs}-\Pf{}//
\glft \trsl{Marek made the problem bigger.} \lingcontext{causative perfective; ergative alignment}//
\endgl
\xe

The applicative \ird{-éb-} is admitted with stative verbs but is semantically
more restricted than with active verbs. It yields a circumstantial reading in
which a causally implicated participant --- typically the entity or situation
responsible for bringing about the state --- is introduced. This use is close
to an experiencer-causative in English (\trsl{gladden}, \trsl{please}). The
causing entity stands in the direct case; the affected participant appears in
the indirect:

\ex
\begingl
\gla Senica Marcí všihnébik.//
\glb news.\Dir{} Marek.\Ind{} be.happy-\Appl{}-\Pf{}//
\glft \trsl{The news gladdened Marek.} \lingcontext{\ird{všihn-}: be happy, be glad}//
\endgl
\xe

Stative verbs do not possess a transitive form and accordingly do not admit the
antipassive. The antipassive is licensed only where an ergative transitive
reading would otherwise be available; since stative verbs are inherently
intransitive, no such reading exists.

Stative verbs may further function as attributive modifiers of nouns and remain
the productive attributive strategy in Iridian. Alongside them stands the small
closed set of paired adnominal property-root forms just mentioned; this use is
described in Chapter~\ref{ch:nouns}.

Property-root-based stative stems also feed the productive derived manner
adverbs in \ird{-ě} (§\,\ref{sec:stative-adv}), yielding lexical adverbial
adjuncts such as \ird{skorně} \trsl{quickly}. These belong to the clause's
preverbal adjunct structure, not to the attributive system proper.

\subsection{Syntactic analysis of stative verbs}
\label{sec:vm-tension-stative}

Under the analysis adopted here, stative verbs are regular finite verbs rather
than a partially exceptional predicate class. A lexical property root combines
with the verbalising suffix \ird{-n-}, analysed syntactically as a stativising
verbal head, to yield an intransitive predicate denoting a state. That
predicate is then selected by the same Asp head that selects other finite
verbs in Iridian, so the full five-way Asp.Align system applies to stative
bases exactly as it does elsewhere in the verbal system.

This account captures the central semantic contrast of the paradigm without
requiring a separate copular projection for ordinary property predication. In
the continuous, Asp selects the stative predicate as a state holding at the
reference time, yielding plain predications such as \ird{Nužot mrocna}
\trsl{the city is beautiful}. In the progressive, perfective, retrospective,
and contemplative, the same Asp system contributes developmental, bounded, or
temporally displaced viewpoints over that stative base; the resulting readings
are therefore naturally inchoative or ingressive (\trsl{be becoming
beautiful}, \trsl{became beautiful}) rather than simple present-state
descriptions. Asp does not cease to apply to stative predicates; rather, its
semantic contribution is interpreted over a predicate whose lexical content is
stative.

A consequence of this analysis for negation deserves explicit statement,
since the surface form of stative predications can resemble zero-copula
equatives. Consider the contrast:

\begin{center}
\begin{tabular}{l l l}
\toprule
\textsc{predication type} & \textsc{affirmative} & \textsc{negative} \\
\midrule
Zero-copula equative & \ird{Marek doktor} & \ird{Marek doktor staly} \\
Stative verbal       & \ird{Nužot mrocna} & \ird{Nužot zámrocna} \\
\bottomrule
\end{tabular}
\end{center}

\medskip
\noindent The two clause types are superficially similar: both surface as [DP
Pred V?] with no overt copula and no aspectual suffix. Yet their negation
strategies differ sharply. The zero-copula equative takes \ird{staly} (the
free-morpheme Neg\textsuperscript{0} exponent), because there is no verbal host
for \ird{zá-} to cliticize onto. The stative verbal predication takes \ird{zá-},
because \ird{mrocna} is a finite verb (the continuous form of \ird{mrocn-}
\trsl{be beautiful}), which provides the verbal host that cliticization requires.

The structural distinction is therefore the presence or absence of a finite
verbal head. In zero-copula equatives, the predicate is a bare DP (or
participant nominal); there is no V head in the clause and AspP is not
projected. In stative verbal predications, the property root combines with
the stativizing suffix \ird{-n-} to yield a proper verbal stem, which is
inflected for Asp.Align in the usual way. Neg\textsuperscript{0} therefore
takes a verbal complement and its exponent is \ird{zá-}. The two predicate
types are formally distinguished by the presence of \ird{-n-} (and the
resulting finite Asp slot), not by anything special about the NegP layer.

The boundary becomes potentially unclear only with underived property adjuncts
that lack \ird{-n-}, if any such items exist; the present analysis treats all
property predications as going through the \ird{-n-} stativizer. See
§\,\ref{sec:vm-tension-neg} for the related open question about
complementarity of the three Neg\textsuperscript{0} exponents.

The syntactic infrastructure of the stative domain therefore distinguishes
three outputs from the same lexical base. First, the finite stative verb heads
the clause's verbal predicate and requires no copula in ordinary property
predication. Second, the \ird{-á} form is a non-finite attributive derivative
inside the DP, corresponding to the adnominal strategy discussed in
Chapter~\ref{ch:nouns}; it is not a finite predicate and cannot bear aspect or
alignment. Third, the \ird{-ě} form is a lexical adverbial derivative that
functions as a clausal adjunct in the preverbal field, not as an inflectional
member of the verbal paradigm. The same lexical base thus feeds clause-level
predication, DP-internal attribution, and adverbial modification through three
distinct syntactic realisations.

Because underived stative predicates are intransitive, the alignment component
of Asp.Align is trivial in their simple finite use: the sole argument is the
direct-case absolutive/nominative argument in every aspect. The ordinary
alignment alternations re-enter the system only when valency-changing
morphology is added. Causativised statives behave like other derived
transitives, showing nominative alignment in the progressive and contemplative
and ergative alignment in the perfective and retrospective. The stative system
thus does not challenge the broader verbal theory; it instantiates it with a
lexically stative base whose aspectual interpretations follow from the
interaction of a regular Asp head with a regular intransitive predicate.


\section{Verbal complements and head finality: preverbal position}
\label{sec:vm-comp-headfinal}

This section develops the formal structural analysis of verbal complements within the head-final clause: how the head-final parameter determines the preverbal position of nominal arguments and embedded clauses, and what syntactic projection each complement type occupies.

\subsection{The base position of the object}
\label{sec:vm-obj-base}

In a head-final language, the complement of a head $H$ is to the
left of $H$ in the linear string. For the lexical predicate V, its
internal argument (patient DP) is the complement of V and therefore
surfaces to V's left:

\begin{center}
\term{[VP [DP patient] V]}
\end{center}

\noindent This is a \emph{base-generated} preverbal position, not
the result of A-bar movement or scrambling. In the Minimalist
framework, head-finality is parameterized as the linearization
convention that assigns $<$Complement, Head$>$ ordering (Kayne's
1994 antisymmetric alternative --- deriving all
SOV orders from movement --- is not adopted here; the parametric
head-final approach is more consistent with the thoroughgoing
head-finality of Iridian across all phrase types).

The patient DP, while base-generated as the complement of V, may
also undergo A-movement to a higher specifier position for topic
purposes. In the ergative aspects, the patient is the absolutive
(topic) argument and moves from [Spec-VP] or [Compl-VP] to
Spec,TopP:

\medskip
\begin{center}
[TopP [patient$_i$] [Top$'$ Top [AspP ... [VP $t_i$ V]]]]
\end{center}
\medskip

\noindent In this derivation, the patient originates as the
complement of V (preverbal by head-finality), then moves to
Spec,TopP (clause-initially) to receive [absolutive] case and
topic status. The surface position of the patient (clause-initial,
absolutive) is therefore derived by A-movement, not base-generation
in TopP.

In the nominative aspects, the patient remains in its base position
as the complement of V (preverbal in the SOV order), while the
agent moves to Spec,TopP:

\medskip
\begin{center}
[TopP [agent$_j$] [Top$'$ Top [AspP ... [vP [VP [patient] V] $t_j$]]]]
\end{center}
\medskip

\noindent The surface word order in both cases is SOV (or more
precisely, absolutive-oblique-V), with the absolutive topic
first and the verb final.

\subsection{The po-clause as a nominal CP complement}
\label{sec:vm-po-comp}

The finite propositional complement (\ird{po}-clause) is a CP that
occupies the complement position of the matrix predicate. The
structure of the complement CP is:

\begin{center}
\term{[CP [TP embedded-clause-\ird{e}] \ird{po}]}
\end{center}

\noindent \ird{Po} is the head of the complement CP (a
complementizer), and the embedded finite clause is the TP whose
head is marked by the linker \ird{-e}. In head-final structure,
the C head follows its complement TP:

\begin{center}
\begin{forest}
  syntax tree
  [CP
    [{TP\\\textit{embedded clause}\\\textit{(verb + -e)}}]
    [{C\\\ird{po}}]
  ]
\end{forest}
\end{center}

\noindent This CP is then the complement of the matrix verb:

\begin{center}
\begin{forest}
  syntax tree
  [VP
    [{CP\\{[\textit{emb-clause}-\ird{e}] [\ird{po}]}}]
    [V]
  ]
\end{forest}
\end{center}

\noindent Head-finality places the CP complement to the left of the
matrix V. The surface word order of a propositional complement
construction is therefore:

\begin{center}
\term{[embedded-TP + \ird{-e}] [\ird{po}] [matrix arguments] [matrix V]}
\end{center}

\noindent This is precisely the attested order (§\,\ref{sec:th-po}).
The analysis requires no movement: the complement CP is
base-generated in [Complement, matrix-VP], and head-finality places
it before the matrix V.

A further consequence: the matrix absolutive argument (the topic)
in a propositional complement clause is most often null (topic
drop). This follows from the heaviness of the CP complement: in
a head-final clause, the CP already occupies the preverbal position
that would be available to an overt absolutive DP, and the topic
is preferentially dropped when the complement clause is highly
informative and heavy.

\paragraph{Po as C$^0$: evidence.}

The analysis of \ird{po} as a C head (and not, for example, as an
adposition or particle) is supported by the following diagnostics:
\begin{enumerate}
\item The embedded clause is fully finite (retains aspect,
  alignment, mood, and person morphology); this is what is
  expected of a CP complement, not a reduced VP.
\item \ird{Po} appears after the embedded clause-final verb
  (the right edge of TP in a head-final CP), in exactly the
  position predicted for a head-final C.
\item \ird{Poby} (the non-assertive counterpart) parallels \ird{po}
  structurally and has the semantics of a complementizer: it
  types the embedded clause as a polar-interrogative content
  (\textit{whether}) rather than a declarative. Complementizer
  alternation is a CP-level phenomenon.
\item The linker \ird{-e} on the embedded verb is the edge-marking
  morpheme at the TP-CP boundary; this is consistent with its role
  as a phonological reflex of \mbox{[Spec,CP] → C} movement or of
  the TP entering the scope of C (§\,\ref{sec:vm-tension-linker}).
\end{enumerate}

\subsection{The purposive converb complement}
\label{sec:vm-purp-comp}

The purposive converb (\ird{-it}) serves as the complement of
desiderative, volitional, and modal predicates. Unlike the
\ird{po}-clause, it is non-finite: it lacks aspect, alignment, and
mood. Structurally it is an aspectless vP or VoiceP projection ---
a verbal phrase headed by the root + any applicable valency
morpheme, terminated by the purposive suffix \ird{-it} rather
than an Asp.Align morpheme.

The purposive converb complement is therefore not a CP but a
smaller projection, most plausibly a bare VoiceP or AspP$_\varnothing$
(an AspP whose head is phonologically null, licensed only when
embedded under a controlling predicate). This is the standard
\term{event complement} (Stowell 1982; Wurmbrand 2001): the predicate denotes an
event rather than a proposition, and it is controlled by the matrix
subject in a typical control configuration.

In head-final structure, the purposive converb complement precedes
the matrix predicate:

\begin{center}
\term{[matrix agent-\textsc{abs}] [purposive-\ird{-it}] [matrix V]}
\end{center}

\noindent The shared-subject (control) reading is the default, with
an implicit PRO subject in the converb clause coreferent with the
matrix absolutive (topic).

%% STUB: The formal structure of the purposive converb complement
%% (VoiceP vs. AspP_null vs. TP with reduced features) should be
%% examined through the following diagnostics:
%% (a) Can the purposive converb take its own topic arguments
%%     distinct from the matrix clause? (→ if yes, it has a
%%     full TP or CP projection)
%% (b) Can the subject of the purposive converb be overtly
%%     expressed, or must it be PRO? (→ tests control vs. raising)
%% (c) Does the purposive converb license aspect-sensitive
%%     alternations in the matrix clause? (→ aspect interaction)
%% These tests bear directly on how high a projection the
%% purposive converb clause actually projects.

\subsection{The preverbal zone: a structural hierarchy}
\label{sec:vm-preverbal-zone}

The preceding analysis allows a unified characterization of the
preverbal zone in the Iridian clause. Reading from left to right,
the preverbal constituents occupy the following structural positions:

\begin{center}
\footnotesize
\begin{tabular}{l l l}
  \toprule
  {\sc position} & {\sc constituent type} & {\sc structural slot} \\
  \midrule
  1 (leftmost) & Topic DP (absolutive)     & Spec,TopP \\
  2            & Zone I peripheral particles & Spec/Head of EvidP, SpeechActP \\
  3            & Clause-typing particle \ird{by} & Head of PolP \\
  4            & Wh-phrase (in questions)  & Spec,FocP \\
  5            & Non-topic arguments       & Spec,IP or IP-adjuncts \\
  6            & Comment-domain operators   & Spec/Head of AspP (zone II) \\
  7            & Phasal particles           & Spec/Head of PhasP (zone III) \\
  8            & CP complement (po-clause) & Compl,VP (innermost) \\
  8 (alt.)     & Purposive converb         & Compl,VP or Spec,vP \\
  9 (rightmost) & Finite verb               & Head of VP / vP \\
  \bottomrule
\end{tabular}
\end{center}

\noindent The verbal complement (whether NP, CP, or converb) is
structurally the innermost preverbal constituent, closest to the
verb. This reflects its base position as the complement of V in the
VP (positions 8 in the table). Higher functional projections
(AspP, IP, TopP) are realized by progressively more leftward
material in the string. The verb itself appears at the absolute
right edge, consistent with the head-final parameter operating
recursively at every projection.

\section{Tension points and areas requiring development}
\label{sec:vm-tensions}

The following issues represent the principal theoretical challenges
in the current analysis. Each will require targeted empirical
investigation, corpus diagnostics, or theoretical elaboration before
the verbal morphosyntax chapter can be considered fully settled.

\subsection{The Asp.Align portmanteau and the Mirror Principle}
\label{sec:vm-tension-asp}

The most fundamental tension in the morphological template is the
\term{portmanteau} character of the Asp.Align morpheme: a single
morpheme does double duty, encoding both the aspectual value of the
predicate and the case-alignment pattern of the clause. The Mirror
Principle predicts that each morpheme corresponds to a single
syntactic head. If Asp and Align are two distinct syntactic
features (housed in AspP and CaseP respectively), the Mirror
Principle predicts separate morphemes for each.

Three responses are available:
\begin{description}[leftmargin=2em, labelwidth=!, labelindent=0em]
\item[\upshape Portmanteau morpheme by fusion (DM).] In DM,
  morphological Fusion (Halle \& Marantz 1993)
  allows two syntactic heads to be realized by a single Vocabulary
  Item when they are adjacent in the morphological structure. AspP
  and CaseP are adjacent in the Iridian verbal complex; their
  adjacency triggers fusion, and the resulting portmanteau is the
  Asp.Align suffix. This is the most parsimonious DM account.
\item[\upshape Single head with bundled features.] AspP is a
  single syntactic head whose feature bundle includes both
  $[$\textsc{aspect}: pf/retro/cont/prog/contempl$]$ and
  $[$\textsc{case-pattern}: erg/nom$]$. This is not a violation
  of the Mirror Principle if the feature bundle defines a single
  head, not two. The objection is theoretical: are aspect and
  alignment independent features (motivated independently) or are
  they necessarily co-specified?
\item[\upshape Aspect drives alignment, not vice versa.] The
  alignment pattern is not an independent feature but a
  \emph{consequence} of the aspect value under the Agree mechanism:
  perfective/retrospective Asp features trigger one Agree pattern
  (ergative), nominative Asp features trigger another (nominative).
  There is only one head (AspP); the alignment is derived by the
  case assignment properties of that head. This is the approach
  assumed in §\,\ref{sec:vm-asp-case} and is consistent with
  aspect-driven ergativity in the broader typological literature
  (Coon 2010; Legate 2008).
\end{description}

\noindent The third approach is preferred here. The desideratum
is to formalize exactly which feature of AspP triggers which
case-assigning behaviour, without reducing the analysis to mere
labelling of the phenomenon.

\subsection{Mood--Aspect competition: why does Mood suppress Aspect?}
\label{sec:vm-tension-mood-asp}

The complete suppression of Asp.Align when a Mood morpheme is
present is unusual. In many languages, mood and aspect are
independently realized (e.g., Turkish, Finnish, Georgian). The
Iridian pattern --- non-indicative mood bleeds aspect entirely ---
requires a formal account.

The DM analysis of \term{impoverishment} (Bonet 1991;
Embick 2010) provides a mechanism: a context-sensitive
rule removes the Asp feature bundle from the morphological structure
when Mood features are present. The Asp node, stripped of its
features, is then realized by a null morpheme (or is deleted). The
consequence is that in non-indicative contexts, the Asp.Align
case-assigning machinery does not operate, and alignment defaults
to nominative.

This account remains to be verified: (a) Is nominative alignment
in mood forms a default (elsewhere case), or is it actively
specified by the Mood head? (b) Can aspect distinctions be
recovered in non-indicative contexts through peripheral particles
or converb structures?

\subsection{The Voice head \ird{-n-} and the causative \ird{-na-}: shared form, distinct functions}
\label{sec:vm-tension-n}

The morpheme \ird{-n-} (the transitive/passive Voice head, appearing
before aspect suffixes on transitive stems) and the causative
\ird{-na-} share a common phonological shape. The analysis adopted
in §\,\ref{sec:vm-n} treats them as formally distinct morphemes
occupying distinct functional positions: \ird{-n-} is the Vocabulary
Item of the VoiceP head, while \ird{-na-} is the head of the higher
CausP projection (§\,\ref{sec:vm-caus}). They are not allomorphs of
a single morpheme, and their distribution is mutually reinforcing
rather than in competition: \ird{-na-} may co-occur with \ird{-n-}
on the same stem (the causative takes a transitive complement), while
\ird{-n-} and \ird{-aš-} cannot (they compete for the same VoiceP
slot).

The outstanding question is whether the phonological resemblance
between \ird{-n-} and \ird{-na-} is synchronically active or merely
a residue of a historical relationship. Two positions are possible.
On the \term{historical coincidence} view, the two morphemes derive
from a common ancestor that bifurcated into CausP and VoiceP
exponents in the history of Iridian, and their phonological overlap
is no longer productive. On the \term{morphological family} view, the
shared \ird{n} is a synchronically transparent marker of
argument-introducing functional heads, realized in full (\ird{-na-})
when an additional argument (the causee) is introduced and in reduced
form (\ird{-n-}) when no new argument is added. The grammar takes no
definitive position on this historical question, but regularizes
\ird{-na-} for the causative and \ird{-n-} for the VoiceP head
throughout.

\subsection{The absolutive--topic mapping: Agree or movement?}
\label{sec:vm-tension-abs-top}

The grammar's analysis treats the absolutive argument as the clause
topic, occupying Spec,TopP. Two formal options:

\begin{description}[leftmargin=2em, labelwidth=!, labelindent=0em]
\item[\upshape A-movement to Spec,TopP.] The absolutive DP moves
  from its base position (complement or specifier of VP/vP) to
  Spec,TopP in the narrow syntax. This movement is driven by a
  [\textsc{top}] feature on the Top head, which probes for a
  matching DP. Case is assigned in Spec,TopP (absolutive as
  a structural case of the topic position). This parallels subject
  raising to Spec,TP in English.
\item[\upshape Long-distance Agree without movement.] The absolutive
  DP remains in its base position; the Top head Agrees with it
  long-distance and values its case as absolutive. The DP's
  surface-initial position is then a consequence of base-ordering
  (head-finality places the specifier of the highest projection
  first) rather than movement.
\end{description}

\noindent The movement account is more standard and predicts
reconstruction effects and binding asymmetries that could in
principle be tested. The long-distance Agree account is more
parsimonious but makes weaker predictions. Which analysis is
correct has consequences for how topic drop is modeled: movement
predicts that topic drop is pro-drop in Spec,TopP; long-distance
Agree predicts that topic drop is simply the absence of a lexical
DP in the base position of the absolutive argument.

\subsection{Ergative case: dependent, structural, or inherent?}
\label{sec:vm-tension-erg}

The transitive/ergative case on the agent in ergative-aligned
aspects can be assigned by several mechanisms:

\begin{description}[leftmargin=2em, labelwidth=!, labelindent=0em]
\item[\upshape Dependent case (Marantz 1991).] The
  ergative is assigned to the agent because another DP (the patient)
  has already received absolutive in the same domain. The ergative
  is thus not assigned by any head but is computed as a dependent
  case: marked case on the DP in the presence of another marked-for-
  absolutive DP.
\item[\upshape Structural case assigned by $v^*$.] The transitive
  light verb $v^*$ assigns ergative to the external argument (agent)
  in its specifier. This is the standard Kratzer/Legate analysis for
  ergative languages.
\item[\upshape Asp-driven ergative.] The Asp head, as argued here,
  is responsible for the case pattern. In the ergative aspects, the
  Asp head's case-assigning operation assigns absolutive to the
  patient and ergative to the agent simultaneously. This requires
  the Asp head to case-assign two DPs, which is unusual in standard
  Agree but not unprecedented (cf.\ Baker 2015 two-probe model).
\end{description}

\noindent The Asp-driven account is consistent with the analysis
developed in §\,\ref{sec:vm-asp-case} and is preferred here. Its
key prediction is that the ergative case is aspect-sensitive: only
in perfective and retrospective aspects does the agent receive
ergative, and this is because only those Asp heads have the
case-assigning feature bundle that triggers the ergative pattern.


\subsection{The conjunctive linker \ird{-e}: see §\,\ref{sec:conjunctive-linker}}

The formal analysis of the conjunctive linker \ird{-e} --- including its phonological characterization, morphosyntactic distribution, the competing unified CP-edge, PF boundary-marker, and two-homophone analyses, and the outstanding empirical questions --- is developed in full in the conjunctive linker section above (§\,\ref{sec:conjunctive-linker}).

\subsection{The NEG prefix vs.\ NegP position: see §\,\ref{sec:vm-neg}}
\label{sec:vm-tension-neg}

The prefix-paradox tension and its resolution via the DM
morphological-incorporation account are discussed in full in
§\,\ref{sec:vm-neg}; the empirical commitments of that account
are listed at the end of §\,\ref{sec:vm-neg-exponents}.


%% ================================================================
%%  SUMMARY TABLE
%% ================================================================

\section{Summary: the verbal complex as a syntactic structure}
\label{sec:vm-summary}

Table~\ref{tab:vm-full-structure} gives the full proposed
hierarchical structure of the Iridian finite verbal clause, from
the highest projection (SpeechActP) to the lowest (VP/Root), with
the morphological exponents of each relevant layer.

\begin{table}[ht]
  \footnotesize\sffamily
  \caption{The Iridian verbal clause: functional hierarchy and
  morphological exponents.}\label{tab:vm-full-structure}
  \medskip
  \begin{tabularx}{\textwidth}{l l l L}
    \toprule
    {\sc level} & {\sc projection} & {\sc morphological exponent}
    & {\sc notes} \\
    \midrule
    Highest  & \textsc{SpeechActP} & Zone I particles (peripheral)
    & Post-topic; speech act framing \\
    & \textsc{EvidP}$_\text{outer}$ & \ird{oče, hlavdy, ktady, či}
    & Clause-level evidential/epistemic \\
    & \textsc{TopP} & [absolutive DP]
    & Topic = absolutive; A-movement target \\
    & \textsc{PolP} & \ird{by}
    & Clause typing: polar interrogative \\
    & \textsc{FocP} & [wh-phrase]
    & Focus/wh movement target \\
    \midrule
    & \textsc{EvidP}$_\text{inner}$ & \ird{-ejte} (?)
    & Indirect evidential; inside verbal complex \\
    & \textsc{MoodP} & \ird{-ím, -é, -ká, -ku, -ěv, -ir}, $\emptyset$
    & Non-indicative moods; suppresses Asp when overt \\
    & \textsc{NegP} & \ird{zá-} (prefix)
    & Clausal negation; PF-cliticized to verbal complex \\
    & \textsc{AspP} & \ird{-ek/-ik, -ac, -a, -ime, -ach} (fused w/ align)
    & Sole obligatory head; aspect + case assignment \\
    \midrule
    & \textsc{CausP} & \ird{-na-}
    & Adds external causer; dominates Voice/ApplP \\
    & \textsc{VoiceP / ApplP} & \ird{-aš-} (antipassive) / \ird{-éb-} (applicative)
    & Valency operations; complementary distribution \\
    & \textsc{VP} & Root stem
    & Lexical predicate; complement of V = preverbal \\
    Lowest & \textsc{VP complement} & NP / CP (\ird{po}) / VoiceP$_{\text{-it}}$
    & Object, po-clause, or purposive converb; all preverbal \\
    \bottomrule
  \end{tabularx}
\end{table}

\noindent The preverbal position of all verbal complements --- NP
objects, CP complements (\ird{po}-clauses), and purposive converb
complements --- follows uniformly from the head-final parameter
applied at the VP level. No movement or extraposition is required:
the complement is base-generated to the left of the verb and
surfaces there in the linear string. Higher projections (AspP,
NegP, MoodP, and the inner evidential) are realized as suffixes on the verbal
complex, which surfaces at the right edge of the clause as required by
pervasive head-finality.

%% ----------------------------------------------------------------
%% END: verb-minim.tex
%% ----------------------------------------------------------------


%% Chapter 4: The Noun Phrase
\chapter{Nouns and the Noun Phrase}\label{ch:nouns}

\section{Introduction}

Nominal morphology in Iridian is relatively straightforward. Nouns are inflected for case; there is no grammatical gender, and number marking is optional. The case system interacts closely with the verbal aspect system: the morphological role of a noun phrase in a clause depends in part on the aspect of the verb, a consequence of Iridian's split-ergative alignment (discussed in detail in Chapter~\ref{ch:verbs}).

\section{The case system}

Iridian has three productive cases: the \term{direct}, the \term{transitive}, and the \term{indirect}. A fourth case, the \term{vocative}, exists for a restricted set of nouns used in direct address. The direct case is the default, uninflected form from which all other cases are derived.

The functions of each case are closely tied to the aspect of the verb in the clause:

\begin{itemize}
	\item The \term{direct case} (\Dir{}) is the unmarked or citation form of a noun. It functions as the unmarked or core argument of the clause: the subject of an intransitive verb, the subject of a transitive verb in the nominative-alignment aspects (continuous, progressive, contemplative), and the object of a transitive verb in the transitive-alignment aspects (perfective and retrospective). It is also used as the predicative complement in equative constructions (zero copula) and with the overt copular verb \ird{neh-}; and as the entity whose existence is asserted in existential constructions with \ird{jec-} and \ird{zad-}.

	\item The \term{transitive case} (\Trs{}) is the ergative case. It marks the agent of a transitive verb in the perfective and retrospective aspects. Outside of those contexts, it does not appear as a core argument.

	\item The \term{indirect case} (\Ind{}) is the oblique or non-core case. It marks the patient of a transitive verb in the nominative aspects (continuous, progressive, contemplative); it also marks instruments, beneficiaries, locations, and other peripheral arguments regardless of aspect. Its function is thus both grammatical and adverbial.

	\item The \term{vocative} (\Voc{}) is used for direct address. It exists only for personal names and a small set of common nouns, primarily kinship terms.
\end{itemize}

\pex
\a
\begingl
\gla Marek zmacime.//
\glb Marek.\Dir{} run-\Prog{}//
\glft \trsl{Marek is running.} \lingcontext{intransitive; agent = \Dir{}}//
\endgl
\a
\begingl
\gla Travat Marku piaštnik.//
\glb bread.\Dir{}.\Def{} Marek.\Trs{} eat-\Pf{}//
\glft \trsl{Marek ate the bread.} \lingcontext{transitive perfective; P = \Dir{}, A = \Trs{}}//
\endgl
\a
\begingl
\gla Marek travě piastime.//
\glb Marek.\Dir{} bread.\Ind{} eat-\Prog{}//
\glft \trsl{Marek is eating bread.} \lingcontext{transitive progressive; A = \Dir{}, P = \Ind{}}//
\endgl
\xe

\section{Inflectional classes}

Nouns in Iridian are best described through two intersecting systems: \term{declension class} and \term{stem grade}. The declension class determines which case endings a noun selects; the stem grade determines which shape of the stem appears before that ending. This allows the grammar to treat forms such as \ird{Marek} $\rightarrow$ \ird{Marku} / \ird{Marcí}, \ird{obel} $\rightarrow$ \ird{oblu} / \ird{oblí}, and \ird{lobek} $\rightarrow$ \ird{lubku} not as isolated curiosities, but as regular consequences of the general morphophonemic system described in Chapter~\ref{ch:phon}.

The surface case inventory remains simple --- direct, transitive, indirect, and vocative --- but the stems that feed those cases may occur in several recurrent grades:

\begin{itemize}
	\item the \term{basic stem}, used in the direct/citation form;
	\item the \term{zero-grade stem}, where weak \ird{e} is deleted before an oblique ending;
	\item the \term{front-raised stem}, where zero-grade formation further yields \ird{e/ě/é $\rightarrow$ i/í};
	\item the \term{back-raised stem}, where zero-grade formation yields \ird{o $\rightarrow$ u};
	\item the \term{soft stem}, where a front-vocalic ending palatalizes the final consonant;
	\item the \term{leveled stem}, where an inherited or high-frequency form overrides the productive pattern.
\end{itemize}

Not every noun shows all these grades, but every morphophonemically complex noun can be described by reference to them. Dictionaries should therefore identify nouns not merely by a citation form, but by a compact inflectional profile: \term{Direct form}, \term{declension class}, \term{weak-\ird{e} status}, \term{raising behavior}, and any \term{leveled indirect or vocative form}. In ordinary prose discussion this grammar cites the direct form alone unless some further principal part is needed.

Native direct-case nouns may end in a consonant, in \ird{-a}, in \ird{-o}, in \ird{-e/-ě}, or in \ird{-ou}. Final \ird{-i} and \ird{-u} are avoided in inherited absolutive nouns and occur chiefly in recent loans, abbreviations, or marginal expressive vocabulary. The declension system must therefore account not only for consonant-final stems, but also for several distinct vowel-final root shapes.

Modern standard Iridian has eighteen declension classes. These are lexically assigned and must be learned noun by noun, but they are not arbitrary: they fall into six broader macrofamilies, each with three internal subclasses. Within a given class the endings and stem alternations are regular; what is lexical is chiefly which class a noun belongs to, whether it tolerates deletion and raising, and whether it preserves an inherited short indirect ending or one of the longer areally remodeled endings.

\begin{table}[ht]
	\scriptsize\sffamily
	\caption{The eighteen noun declension classes.}\label{tab:declension}
	\medskip
	\begin{tabularx}{\textwidth}{l L l l l L}
		\toprule
		{\sc class} & {\sc type} & {\sc \Trs{}} & {\sc \Ind{}} & {\sc \Voc{}} & {\sc stem behavior} \\
		\midrule
		1 & hard consonant plain & \ird{-u} & \ird{-í} & \ird{-o} & basic stem; no obligatory weak-\ird{e} deletion \\
		2 & hard consonant deleting & \ird{-u} & \ird{-ení} & \ird{-o} & productive weak-\ird{e} deletion plus expanded oblique ending \\
		3 & hard consonant front-raised & \ird{-u} & \ird{-í} & \ird{-o} & deletion plus \ird{e/ě/é $\rightarrow$ i/í} \\
		4 & hard consonant back-raised & \ird{-u} & \ird{-ení} & \ird{-e} & deletion plus \ird{o $\rightarrow$ u}; long front-vocalic indirect \\
		5 & hard animate arealized & \ird{-a} & \ird{-ovi} & \ird{-e} & Slavic-like animate oblique endings beside inherited stem grades \\
		6 & hard consonant leveled & \ird{-u} & lexical \ird{-í / -ení} & zero/\ird{-o} & productive stem behavior, but ending choice partly analogical \\
		7 & velar plain & \ird{-u} & \ird{-e} & \ird{-e} & regular velar affrication in \Ind{} and \Voc{} \\
		8 & velar deleting & \ird{-u} & \ird{-ení} & \ird{-e} & weak-\ird{e} deletion plus velar softening under an expanded indirect ending \\
		9 & velar arealized & \ird{-a} & \ird{-evi} & \ird{-e} & animate-like oblique remodelled under areal pressure \\
		10 & soft consonant plain & \ird{-y} & \ird{-ě} & \ird{-e} & already soft stem; little further alternation \\
		11 & soft consonant contracted & \ird{-y} & \ird{-i} & \ird{-e} & contracted front oblique with little further alternation \\
		12 & sibilant / affricate expanded & \ird{-y} & \ird{-ení} & \ird{-e} & class-specific spellout after sibilants and affricates \\
		13 & \ird{-a} stem & \ird{-y} & \ird{-ě} & --- & low alternation; default native vocalic class \\
		14 & \ird{-o} stem (\ird{o + m > um}) & \ird{-a} & \ird{-um} & --- & final \ird{-o} replaced in \Trs{}; coalesces with nasal indirect marker \\
		15 & front-vowel stem in \ird{-e/-ě} & \ird{-e} & \ird{-ení} & --- & levelled front-vowel stems with analogically extended indirect \\
		16 & consonant \ird{-a/-am} & \ird{-a} & \ird{-am} & \ird{-e} & conservative oblique family with little softening \\
		17 & weak-\ird{e} \ird{-a/-am} & \ird{-a} & \ird{-am} & \ird{-e} & productive deletion before oblique endings, often with syllabic liquids \\
		18 & contracting \ird{-ou} & \ird{-ovy} & \ird{-ově / -ovení} & --- & inherited contraction plus occasional areal extension of the indirect \\
		\bottomrule
	\end{tabularx}
\end{table}

The classes are numerous because Iridian noun morphology is intentionally noun-heavy. The same four cases recur everywhere, but the path from citation form to oblique form is richly differentiated by class. Some classes preserve short inherited endings, others generalize longer front-vocalic endings such as \ird{-ení}, and a few animate classes have borrowed or convergently renewed endings such as \ird{-ovi} and \ird{-evi}. The language thus preserves an agglutinative syntax and transparent case inventory while allowing a much denser Slavic-looking nominal surface.

\subsection{Macrofamily I: hard consonant stems}

Classes~1--6 comprise the hard consonant family. Here the microfamilies show genuine ending divergence as well as stem divergence. Class~1 preserves the short inherited indirect \ird{-í}. Class~2 extends that oblique to \ird{-ení}, giving forms of the type \ird{louž} $\rightarrow$ \ird{loužení}. Classes~3 and~4 add front-vowel and back-vowel raising respectively, but differ again in whether the indirect is the short \ird{-í} or the expanded \ird{-ení}. Class~5 contains animate nouns remodeled under areal pressure, with \ird{-a} in the transitive and \ird{-ovi} in the indirect. Class~6 keeps a partly leveled oblique paradigm in high-frequency nouns and personal names.

Representative outputs are:

\pex
\a \ird{Marek} $\rightarrow$ \ird{Marku} / \ird{Marcí} \lingcontext{short inherited indirect}
\a \ird{louž} $\rightarrow$ \ird{loužu} / \ird{loužení} \lingcontext{expanded indirect in Class~2}
\a \ird{lobek} $\rightarrow$ \ird{lubku} / \ird{lubcení} \lingcontext{back-raised stem plus long indirect}
\a \ird{Marek} in the arealized animate class $\rightarrow$ \ird{Marka} / \ird{Markovi}
\xe

\subsection{Macrofamily II: velar stems}

Classes~7--9 are a specialized subset of consonant-final nouns whose final consonant is velar. They preserve the visible velar alternations that most strongly index the language's Slavic surface profile. In this family the key analysis is velar affrication: after weak-\ird{e} deletion, /k/ surfaces as \ird{c} before front-vocalic oblique endings. Class~7 takes a short indirect in \ird{-e}; Class~8 uses the longer \ird{-ení}; Class~9 has animate-like areal endings in \ird{-evi}. The distinction among the three classes thus lies not only in deletion and raising, but also in which oblique ending the affricated velar stem selects.

\pex
\a \ird{perek} $\rightarrow$ \ird{prku} / \ird{prce}
\a \ird{človek} $\rightarrow$ \ird{človku} / \ird{človcení}
\a \ird{rok} $\rightarrow$ \ird{roka} / \ird{rocevi}
\xe

\subsection{Macrofamily III: soft consonant stems}

Classes~10--12 contain stems that already end in a soft consonant, a sibilant, or an affricate. Class~10 preserves the inherited \ird{-ě} indirect. Class~11 contracts that ending to \ird{-i}. Class~12 generalizes an expanded \ird{-ení}, often under areal pressure and especially after sibilants. This is the chief source of contrasts such as short \ird{Marcí} beside longer \ird{loužení}.

Because these nouns are already soft or sibilant-final, they usually show less dramatic alternation than hard consonant nouns. Their role in the system is not to showcase weak-\ird{e} deletion, but to preserve a dense inherited declensional texture.

\subsection{Macrofamily IV: vocalic stems}

Classes~13--15 are the main vowel-final families. Class~13 contains the default \ird{-a} stems such as \ird{trava}, with transitive \ird{-y} and indirect \ird{-ě}. Class~14 contains back-vowel \ird{-o} stems, whose transitive replaces the final vowel with \ird{-a} and whose indirect shows the regular coalescence \ird{o + m > um}; thus \ird{bro} $\rightarrow$ \ird{bra} / \ird{brum}. Class~15 contains front-vowel stems in \ird{-e/-ě}; these are often historically complex and are more heavily leveled than the other classes, and many of them now show an analogically extended indirect in \ird{-ení}. Together with the separate \ird{-ou} class below, these vocalic families account for the full set of native vowel-final absolutive shapes.

\pex
\a \ird{trava} $\rightarrow$ \ird{travy} / \ird{travě}
\a \ird{bro} $\rightarrow$ \ird{bra} / \ird{brum}
\a \ird{země} $\rightarrow$ \ird{zeme} / \ird{zemení}
\xe

\subsection{Macrofamily V: the \ird{-a/-am} oblique family}

Classes~16 and~17 preserve the old consonant-final family whose transitive is \ird{-a} and indirect \ird{-am}. Class~16 is conservative and comparatively stable. Class~17, by contrast, subjects the stem to productive weak-\ird{e} deletion before those oblique endings and often resolves the resulting cluster through a syllabic liquid rather than by restoring the vowel. The especially productive subtype is the \term{CeLeC} pattern, where the direct stem contains weak \ird{e} before a liquid and a final consonant, and the oblique stem surfaces as \ird{CLC-} plus the case vowel licensed by its class.

\pex
\a \ird{byl} $\rightarrow$ \ird{byla} / \ird{bylam}
\a \ird{jelen} $\rightarrow$ \ird{jlna} / \ird{jlnam}
\a \ird{berel} in the \ird{CeLeC} subtype $\rightarrow$ \ird{brla} / \ird{brlam}
\a \ird{berel} in the most reduced pronunciation $\rightarrow$ [br̩lam]
\xe

\subsection{Macrofamily VI: the contracting \ird{-ou} class}

Class~18 contains inherited \ird{-ou} nouns, above all deverbal nominals. These follow a distinct but internally regular pattern: transitive \ird{-ovy}, indirect \ird{-ově}, and contraction in the direct definite form. Some members, especially literary or formal nouns, also admit an analogically extended indirect in \ird{-ovení}. This class is highly visible in nominalizations and therefore serves as the chief bridge between noun-heavy morphophonology and the rest of the grammar.

\pex
\a \ird{piaščikou} $\rightarrow$ \ird{piaščikovy} / \ird{piaščikově}
\a \ird{hladou} $\rightarrow$ \ird{hladovy} / \ird{hladovení}
\a \ird{piaščikou + -ot} $\rightarrow$ \ird{piaščikovat}
\xe

\subsection{Case endings and stem selection}\label{sec:noun-stem-grades}

The relation between case and stem grade is summarized below.

\begin{table}[ht]
	\footnotesize\sffamily
	\caption{Case endings and their morphophonemic effects.}\label{tab:noun-stem-selection}
	\medskip
	\begin{tabularx}{0.95\textwidth}{L L L}
		\toprule
		{\sc case slot} & {\sc typical endings} & {\sc stem selected} \\
		\midrule
		Direct & zero & basic stem \\
		Transitive & \ird{-u}, \ird{-y}, \ird{-e}, \ird{-a}, \ird{-ovy} & basic, zero-grade, raised, or vowel-replacing stem according to declension class \\
		Indirect & \ird{-í}, \ird{-e}, \ird{-ení}, \ird{-ovi}, \ird{-evi}, \ird{-ě}, \ird{-i}, \ird{-um}, \ird{-am}, \ird{-ově}, \ird{-ovení} & soft stem for front-vocalic endings; zero-grade or raised stem where the class licenses deletion; some classes replace final root vowels or add areally expanded oblique material \\
		Vocative & \ird{-o}, \ird{-e}, zero & basic or soft stem; often levelled in names and kin terms \\
		\bottomrule
	\end{tabularx}
\end{table}

The transitive and indirect cases are therefore not merely suffixes added to an invariant base. Rather, the ending first selects a stem grade, and only then surfaces in its overt form. The rule order is the same as in Chapter~\ref{ch:phon}: suffix class selection first, then weak-\ird{e} deletion or retention, then any raising, then soft-stem selection if the ending is front-vocalic, and only after that the cluster and voicing repairs that yield the final phonological form.

Irregularity in the noun system is best understood as \emph{class assignment plus stem-grade irregularity}, not as unconstrained suppletion. Some nouns fail to raise where the general pattern would suggest it, some preserve weak \ird{e} in only one oblique, and some classes have generalized longer endings such as \ird{-ení} or animate-like endings such as \ird{-ovi} under areal pressure. Once these lexical facts are known, however, the resulting paradigms are still learnable by class rather than item by item.

\pex
\a \ird{Marek} $\rightarrow$ \ird{Marku} / \ird{Marcí}
\a \ird{perek} $\rightarrow$ \ird{prku} / \ird{prce}
\a \ird{bro} $\rightarrow$ \ird{bra} / \ird{brum}
\a \ird{louž} $\rightarrow$ \ird{loužu} / \ird{loužení}
\a \ird{obel} $\rightarrow$ \ird{oblu} / \ird{oblí}
\a \ird{lobek} $\rightarrow$ \ird{lubku}
\a \ird{piaščikou} $\rightarrow$ \ird{piaščikovy} / \ird{piaščikově}
\xe

\subsection{Irregular nouns}

A small number of nouns have irregular inflectional patterns that do not conform neatly even to the eighteen-class system. These are best treated as explicitly levelled lexemes. Dictionaries accordingly list any exceptional indirect or vocative form when it cannot be recovered from the noun's class and grade behavior alone.

\section{The vocative}

The vocative is used for direct address. It is productive for personal names and for a closed set of common nouns, mostly kinship terms. Outside this set, nouns are addressed with the direct case form.

Nouns with vocative forms include the following kinship terms:

\begin{table}[ht]
	\footnotesize\sffamily
	\caption{Kinship terms with vocative forms.}\label{tab:vocative}
	\medskip
	\begin{tabularx}{0.7\textwidth}{L l l}
		\toprule
		{\sc direct} & {\sc vocative} & {\sc meaning} \\
		\midrule
		\ird{papka}  & \ird{papko}  & father \\
		\ird{mamka}  & \ird{mamko}  & mother \\
		\ird{mlazka} & \ird{mlazko} & brother \\
		\ird{voska}  & \ird{vosko}  & sister \\
		\bottomrule
	\end{tabularx}
\end{table}

For personal names the vocative is formed in a manner predictable from the inflectional class of the name (see Table~\ref{tab:declension} for the hard consonant plain type exemplified by \ird{Marek}, vocative \ird{Marko}). The vocative list is not exhaustive; other common nouns may have vocatives, but their use has become increasingly restricted to formal and literary registers.

\section{Definiteness}\label{sec:definiteness}

Iridian possesses a postposed determiner, which appears as \ird{-ot} after a consonant and as \ird{-t} after a vowel. Thus \ird{bylot} signifies \trsl{the child}, and \ird{travat} \trsl{the bread}. This suffix is not to be understood as answering in every particular to the English definite article. Its office is rather to mark a noun as referring to something already present to the mind of the speaker and hearer, whether by previous mention, by the circumstances of the utterance, or by some other sort of immediate intelligibility. It is therefore more exact to say that the suffix marks an \textit{anchored} or \textit{determinate} reference than that it merely marks definiteness in the narrower sense.

In form the determiner is attached to the fully inflected noun. Since the direct case is without overt ending, the suffix there appears immediately upon the nominal stem; where, however, the noun bears some case termination, the determiner follows that termination. One may accordingly set side by side such forms as \ird{byl} \trsl{child}, \ird{bylot} \trsl{the child}, \ird{bylam} \trsl{child} in the indirect case, and \ird{bylamot} \trsl{the child} in the same case.

\ex
\begingl
\gla byl \qquad bylot \qquad bylam \qquad bylamot//
\glb child.\Dir{} \qquad child.\Dir{}.\Def{} \qquad child.\Ind{} \qquad child.\Ind{}.\Def{}//
\glft \trsl{child, the child, child (oblique), the child (oblique)}//
\endgl
\xe

The use of the determiner is optional. The bare noun is the ordinary and least marked form, and may be employed where the reference is indefinite, generic, kind-denoting, weak, or otherwise not yet fixed to the common ground of discourse. The suffixed form, on the other hand, presents the referent as one which is in some manner already recoverable. Very often it is anaphoric; not seldom it is situationally unique; at other times it is contrastively selected or resumed as a discourse-prominent participant. The opposition is therefore not simply that between \trsl{definite} and \trsl{indefinite}, but rather that between a noun which is left unanchored and one which is brought into a more determinate relation with the discourse.

This may be seen from a simple contrast. In the sentence below the object is construed as a particular loaf or portion of bread already identifiable in the context:

\ex
\begingl
\gla Travat Marku piaštnik.//
\glb bread.\Dir{}.\Def{} Marek.\Trs{} eat-\Pf{}//
\glft \trsl{Marek ate the bread.}//
\endgl
\xe

If, however, the noun appears without the determiner, the reference is left less fixed and may be understood in a weaker or less individuated sense:

\ex
\begingl
\gla Trava Marku piaštnik.//
\glb bread.\Dir{} Marek.\Trs{} eat-\Pf{}//
\glft \trsl{Marek ate bread; Marek ate some bread.}//
\endgl
\xe

It must be observed that this suffix does not serve to indicate grammatical relation as such. It does not mark subjecthood, nor is it imposed by the mere fact that a noun stands as a core argument. In transitive clauses of the perfective or retrospective, it is indeed especially frequent with direct-case patient nouns; but this is because such nouns are often conceived as individuated and discourse-recoverable, not because the syntax requires the suffix. The same determiner may, where the sense calls for it, be found upon nouns in other cases also; yet outside the direct case such use is generally less neutral and more plainly expressive of topical, resumptive, contrastive, or otherwise salient reference.

The vocative excludes the determiner. In forms of direct address the person addressed is already sufficiently identified by the speech-act itself, and no further anchoring is needed. Generic expressions likewise dispense with the suffix as a rule, the bare noun being in such contexts the ordinary vehicle of the class or kind as such.

Beside this determiner the language possesses two preposed particles, \ird{si} and \ird{druh}, which are employed when the speaker wishes to introduce a referent as particularized, yet not as already established in the common ground. These particles thus stand, in point of meaning, between the wholly bare noun and the noun provided with the determiner. They do not make the noun definite; rather, they indicate what may be called specific indefiniteness.

Of these two particles, \ird{si} is the lighter and more usual. It answers approximately to \trsl{a certain} or \trsl{some particular}, and is the common means of introducing a referent which the speaker has in view, though the hearer has not yet been led to identify it.

\ex
\begingl
\gla si byl//
\glb \Spec{} child.\Dir{}//
\glft \trsl{a certain child; some particular child}//
\endgl
\xe

The particle \ird{druh} is weightier in tone and somewhat more selective in force. It often suggests \trsl{a given} or \trsl{one particular}, and may accordingly appear where the speaker means to single out the referent with greater insistence, or where the style is more elevated or deliberate.

\ex
\begingl
\gla druh byl//
\glb \Spec{} child.\Dir{}//
\glft \trsl{a certain particular child; one given child}//
\endgl
\xe

The difference between \ird{si} and \ird{druh} is not absolute, and their spheres overlap. In many passages either may stand without serious change of denotation. Still, the general tendency is plain enough: \ird{si} is slighter, more colourless, and more readily employed in ordinary discourse, whereas \ird{druh} is stronger, more selective, and not seldom somewhat more literary in effect.

These particles do not ordinarily co-occur with the determiner \ird{-ot/-t}. The reason is semantic. A noun marked by \ird{si} or \ird{druh} is presented as particularized indeed, but not yet anchored; a noun marked by \ird{-ot/-t} is presented as already anchored and recoverable. The two constructions are therefore, in the normal language, held apart.

Number does not directly alter this opposition. The noun itself is not obligatorily inflected for singular and plural. Plurality, where it is to be expressly marked, is indicated by the proclitic \ird{ne}, which stands at the left edge of the noun phrase. The plural marker and the determiner are independent of one another in form and function: the former denotes plurality, the latter anchoring.

\ex
\begingl
\gla ne byl \qquad ne bylot//
\glb \Pl{} child.\Dir{} \qquad \Pl{} child.\Def{}//
\glft \trsl{children \qquad the children}//
\endgl
\xe

The particles of specific indefiniteness may likewise accompany plural expressions:

\ex
\begingl
\gla ne si byl \qquad ne druh byl//
\glb \Pl{} \Spec{} child.\Dir{} \qquad \Pl{} \Spec{} child.\Dir{}//
\glft \trsl{some particular children \qquad a certain set of children}//
\endgl
\xe

When \ird{ne} is joined with a numeral or other quantifying word, it no longer denotes mere plurality, but rather an approximate number. In such combinations the determiner, as well as the particles \ird{si} and \ird{druh}, is ordinarily avoided, since approximate quantified expressions are not commonly treated as either anchored or narrowly individualized noun phrases.

A word must also be added concerning demonstratives. In the unmarked construction a demonstrative stands without the postposed determiner. Nevertheless, the language permits, in emphatic or contrastive style, the combination of demonstrative and determiner in the same noun phrase. Such double determination is not neutral, but serves to single out the referent with unusual force, whether resumptive, affective, or contrastive.

\ex
\begingl
\gla ša byl \qquad ša bylot//
\glb \Dem{}.\Prox{} child.\Dir{} \qquad \Dem{}.\Prox{} child.\Def{}//
\glft \trsl{this child \qquad this very child; this particular child here}//
\endgl
\xe

The full semantic range of the determiner, together with its relation to deixis, anaphora, topicality, contrast, and the introduction and subsequent tracking of referents, will be considered more fully in Chapter~\ref{ch:discourse}. The formal paradigm of the demonstratives themselves is given in Section~\ref{sec:demonstratives}.
\section{Number}

Iridian nouns are not formally inflected for number. A noun in the direct case form may be interpreted as singular or plural depending on context. The plural particle \ird{ne} may be used to explicitly mark plurality; it is a proclitic attaching to the leftmost element of the noun phrase.

\ex
\begingl
\gla ne byl//
\glb \Pl{} child//
\glft \trsl{children}//
\endgl
\xe

The use of \ird{ne} is optional. In contexts where number is recoverable from numerals, quantifiers, or shared knowledge, the particle is typically omitted. When used with a numeral or other quantifier, \ird{ne} signals approximation rather than plain plurality: \ird{ne hrona byl} means \trsl{about three children}, not \trsl{three children.}

\ird{Ne} is a proclitic and attaches to the front of the noun phrase, including any preceding demonstratives or modifiers. It cannot appear inside a possessive noun phrase between an oblique possessor and the head noun. \ird{Ne} cannot be combined with a \ird{zad-} (negative existential) clause.

\section{Uses of the indirect case}\label{sec:indirect-uses}

Beyond its role as the patient case of transitive clauses in the nominative-alignment aspects, the indirect case is the general non-core case of the noun phrase. It covers the functions that in older descriptions were distributed across several narrower cases. In modern analysis these uses form a single domain: the indirect marks dependence on another noun, relation to a source or whole, and a wide range of peripheral participant roles.

\subsection{Possession and nominal dependence}

The most basic non-core use of the indirect is to mark a dependent noun within the noun phrase, especially a possessor. The dependent noun precedes the head noun and forms a tight constituent with it.

\pex
\irdp{Marcí dum}{Marek's house}\\
\irdp{mámcí hašek}{my mother's bag}\\
\irdp{ša študencí tóm}{this student's book}
\xe

Demonstratives and other phrasal modifiers precede the entire possessive phrase and do not intervene between the oblique possessor and the head noun. The plural clitic \ird{ne}, by contrast, is placed immediately before the element whose plurality it marks.

\pex
\a \irdp{ša študencí tóm}{this student's book}
\a \irdp{ně študencí tóm}{the students' book}
\a \irdp{študencí ně tóm}{the student's books}
\xe

The same oblique dependency is used more broadly for close nominal relations that English often renders with \trsl{of}: kinship, authorship, association, and other kinds of inherent or identifying linkage.

\subsection{Material and partitive relations}

The indirect also marks the substance from which something is made and the noun that supplies the whole or source domain in a partitive relation. In this use the head noun denotes an object, measure, or portion, while the indirect noun names the material or larger whole from which it is drawn.

\ex
\irdp{kuní prosc}{silver spoon}
\xe

\ex
\begingl
\gla Ona pive štava unarížčem.//
\glb one beer-\Ind{} mug-\Dir{} \Refl{}-order-\mk{av-pv-1s}//
\glft \trsl{I ordered a mug of beer.}//
\endgl
\xe

Where the construction refers to an unbounded substance or to a portion abstracted from a whole, the indirect is the normal choice. By contrast, when the language counts discrete members of a set, Iridian prefers an overt quantificational construction rather than a bare oblique partitive dependent.

\pex
\a
\begingl
\gla Kroumaště na po zma jeca pivo.//
\glb refrigerator.\Ind{} in still few \Exst{}-\Cont{} beer.\Dir{}//
\glft \trsl{There's still some beer left in the refrigerator.}//
\endgl
\a
\begingl
\gla Študencí na hroná.//
\glb student-\Ind{} in three//
\glft \trsl{three of the students}//
\endgl
\xe

\subsection{Source and movement away}

With verbs of departure, removal, and separation, the indirect marks the source from which motion proceeds. This source reading contrasts with the direct case, which in comparable verbal environments may instead foreground the affected place as a more central argument.

\pex
\a
\begingl
\gla Dumí palžek.//
\glb house-\Ind{} leave-\Pf{}//
\glft \trsl{I left the house.}//
\endgl
\a
\begingl
\gla Dum palzinek.//
\glb house.\Dir{} leave-\Pf{}//
\glft \trsl{I left the \emph{house}.}//
\endgl
\xe

The same source-oriented oblique extends naturally to nouns followed by postpositions of separation, origin, exclusion, and replacement.

\section{Pronouns and demonstratives}\label{sec:pronouns}

Iridian distinguishes first and second person by means of a small set of personal pronouns. There are no independent personal pronouns of the third person. Reference to third persons is instead managed by the demonstrative system, together with a reduced set of anaphoric forms derived from it. The personal pronouns are inflected for the three ordinary cases, and each has beside it a possessive clitic, of which more will be said presently.

\begin{table}[ht]
	\footnotesize\sffamily
	\caption{Personal pronouns and possessive clitics.}\label{tab:personal-pronouns}
	\medskip
	\begin{tabularx}{0.75\textwidth}{L l l l l}
		\toprule
		              & {\sc direct} & {\sc transitive} & {\sc indirect} & {\sc clitic} \\
		\midrule
		1\textsc{sg}  & \ird{dá}   & \ird{dam}  & \ird{že}   & \ird{-(e)m}  \\
		2\textsc{sg}  & \ird{já}   & \ird{jí}   & \ird{jena} & \ird{-(e)š}  \\
		1\textsc{pl}  & \ird{mé}   & \ird{mně}  & \ird{mi}   & \ird{-(e)m}  \\
		2\textsc{pl}  & \ird{tová} & \ird{tvy}  & \ird{tě}   & \ird{-(e)st} \\
		\bottomrule
	\end{tabularx}
\end{table}

The indirect case of the personal pronoun serves, among its other functions, to express the possessor, precisely as the indirect case of an ordinary noun may do elsewhere in the language (Section~\ref{sec:indirect-uses}). Thus \irdp{že dum}{my house} is formally no different in structure from the corresponding nominal possessor construction. The possessor stands before the head noun in the usual prenominal position.

Alongside this analytic construction there exists a shorter and in many contexts commoner possessive expression by means of enclitic forms attached to the right edge of the noun. The first-person clitic appears as \ird{-m} after a vowel-final stem and as \ird{-em} after a consonant-final stem; the second-person singular behaves in like manner as \ird{-š} or \ird{-eš}, and the second-person plural as \ird{-st} or \ird{-est}. Hence \irdp{dumem}{my house; our house}. It will be observed that the first singular and the first plural have here fallen together in the same clitic form, \ird{-(e)m}; the distinction, where it matters, is left to the context or else restored by the use of the free pronoun.

The free indirect form and the clitic are not mutually exclusive. On the contrary, the language not seldom employs both at once where emphasis, contrast, or explicitness is desired, as in \irdp{že dumem}{my house}. Much the same liberty is found when the noun is at the same time accompanied by a demonstrative: \irdp{tuto že tom}{that book near you of mine}, \irdp{tuto tomem}{that book near you, mine}, and \irdp{tuto že tomem}{that book near you of mine} are all possible, though they differ in weight and emphasis.

Third-person reference belongs properly to the demonstrative system. Iridian possesses three demonstrative stems, distinguished according to the familiar threefold deictic opposition of nearness to the speaker, nearness to the addressee, and distance from both. These may be termed, in the usual fashion, proximal, medial, and distal.

\begin{table}[ht]
	\footnotesize\sffamily
	\caption{Demonstrative paradigm.}\label{tab:demonstratives}
	\medskip
	\begin{tabularx}{0.65\textwidth}{L l l l}
		\toprule
		           & {\sc direct} & {\sc transitive} & {\sc indirect} \\
		\midrule
		Proximal   & \ird{tak}  & \ird{tam}   & \ird{ten}  \\
		Medial     & \ird{tuto} & \ird{toho}  & \ird{taly} \\
		Distal     & \ird{je}   & \ird{jaky}  & \ird{jení} \\
		\bottomrule
	\end{tabularx}
\end{table}

In their primary and local use, these forms answer to the ordinary distinctions of deixis. Thus \irdp{je tom}{that book yonder} contains the distal form, whereas \irdp{tuto pect}{that cup near you} contains the medial. The indirect case of the demonstrative may further serve as a possessor in the same manner as the indirect of a noun or personal pronoun. One may therefore say \irdp{ten dum}{his or her house}, \irdp{taly pect}{his or her cup}, or \irdp{jení vrad}{his, her, or their town}. In such expressions the deictic opposition is preserved, but it is readily extended into discourse. The choice of demonstrative may indicate not only where the possessor stands in relation to speaker and hearer, but also how the referent is situated in the flow of the discourse.

This latter point is of some importance. When used pronominally, the three demonstratives do not merely point in space; they also arrange third-person referents according to accessibility and prominence. The proximal is naturally employed of the current topic or most salient participant; the medial may denote a secondary but still active referent, often one associated with the addressee’s sphere; and the distal tends toward what may conveniently be called obviative use, namely a referent less central, more displaced, backgrounded, or set over against another.

\begin{table}[ht]
	\footnotesize\sffamily
	\caption{Discourse function of the demonstrative series.}\label{tab:dem-discourse}
	\medskip
	\begin{tabularx}{0.9\textwidth}{L L L}
		\toprule
		{\sc form}   & {\sc spatial value}    & {\sc discourse value} \\
		\midrule
		Proximal & near speaker       & current topic, most accessible referent \\
		Medial   & near addressee     & addressee-linked or secondary active referent \\
		Distal   & away from both     & obviative, backgrounded, displaced, or contrastive referent \\
		\bottomrule
	\end{tabularx}
\end{table}

By this means Iridian is able to keep apart several third-person participants without the aid of verbal agreement. The forms are not mere substitutes for a single abstract \trsl{he} or \trsl{she}, but carry with them a definite discourse colouring.

\pex
\begingl
\gla Tak vrbnik, tuto spasnik, je těsnik.//
\glb \Prox{}.\Dir{} stand.up-\Pf{}, \Med{}.\Dir{} sit.down-\Pf{}, \Dist{}.\Dir{} leave-\Pf{}//
\glft \trsl{He stood up, he sat down, that one left.}//
\endgl
\xe

Alongside the full demonstratives there exists a reduced anaphoric set, phonologically derived from the distal series. These forms are used only pronominally, and only where the referent is already highly active and unmistakable in the discourse. They do not properly point; they merely resume.

\begin{table}[ht]
	\footnotesize\sffamily
	\caption{Reduced anaphoric forms (from distal demonstrative).}\label{tab:anaphoric-pronouns}
	\medskip
	\begin{tabularx}{0.6\textwidth}{L l l}
		\toprule
		{\sc case}   & {\sc full distal} & {\sc reduced} \\
		\midrule
		Direct     & \ird{je}   & \ird{i}  \\
		Transitive & \ird{jaky} & \ird{ji} \\
		Indirect   & \ird{jení} & \ird{ni} \\
		\bottomrule
	\end{tabularx}
\end{table}

The distinction is chiefly one of pragmatic force. The full demonstrative retains its deictic or contrastive value, and may be chosen whenever the speaker wishes to indicate distance, opposition, or a stronger gesture of reference. The reduced form, by contrast, is lighter and more nearly a pure anaphor. In broad terms one may say that \ird{tak} points, whereas \ird{i} merely carries forward an already established referent.

In ordinary connected speech, moreover, Iridian shows a marked tendency to omit pronouns wherever they are recoverable from the situation or the preceding discourse. It is, in this sense, a strongly pro-drop language. Overt personal pronouns are introduced chiefly for contrast, insistence, disambiguation, or some comparable rhetorical purpose. This tendency is especially notable in the second person. Direct pronominal reference to the addressee is relatively marked, and may suggest impatience, familiarity, or an undesirable directness. Courteous usage therefore prefers, wherever an overt expression is required, some noun of address instead: a personal name, a title, a kin term, the designation of an office, or the like. The most unmarked course is commonly silence; next in neutrality comes the respectful nominal expression; the overt second-person pronoun is the strongest and least neutral of the available choices.

There remains to be noticed the reflexive. Iridian possesses a distinct reflexive anaphor, whose forms show some remodelling under Slavic influence. This element is not a free pronoun of reference like the demonstratives, but a bound anaphor whose office is solely to mark local coreference.

\begin{table}[ht]
	\footnotesize\sffamily
	\caption{Reflexive anaphor forms.}\label{tab:reflexive}
	\medskip
	\begin{tabularx}{0.5\textwidth}{L l}
		\toprule
		{\sc case}   & {\sc form} \\
		\midrule
		Direct     & \ird{svem} \\
		Transitive & \ird{svě}  \\
		Indirect   & \ird{svy}  \\
		\bottomrule
	\end{tabularx}
\end{table}

The reflexive does not introduce a new discourse participant, nor does it rank one referent against another. It merely indicates that a given argument is coreferential with the local controller. For that reason it does not compete with the demonstratives in the ordinary tracking of discourse referents. The contrast is plain in the following examples:

\pex
\a
\begingl
\gla Svem Marku vednik.//
\glb \Refl{}.\Dir{} Marek.\Trs{} see-\Pf{}//
\glft \trsl{Marek saw himself.}//
\endgl
\a
\begingl
\gla I Marku vednik.//
\glb 3.\Dist{}.\Dir{} Marek.\Trs{} see-\Pf{}//
\glft \trsl{Marek saw him.} \lingcontext{non-reflexive; referent distinct from Marek}//
\endgl
\xe

The same opposition appears in possessive expressions. A demonstrative possessor remains non-reflexive; it refers to some other salient discourse participant, not back to the controller of the clause.

\pex
\a
\begingl
\gla Jení dum Marku vednik.//
\glb \Dist{}.\Ind{} house.\Dir{} Marek.\Trs{} see-\Pf{}//
\glft \trsl{Marek saw his house.} \lingcontext{someone else's house, not Marek's}//
\endgl
\xe

This distinction proves particularly useful where more than one third-person referent is simultaneously in play. The demonstrative system can separate discourse roles and degrees of accessibility; the reflexive, for its part, cuts through ambiguity by marking local identity directly.

\pex
\a
\begingl
\gla Je tam vednik.//
\glb \Dist{}.\Dir{} \Dist{}.\Trs{} see-\Pf{}//
\glft \trsl{He\textsubscript{1} saw him\textsubscript{2}.} \lingcontext{two distinct referents}//
\endgl
\a
\begingl
\gla Svem tam vednik.//
\glb \Refl{}.\Dir{} \Dist{}.\Trs{} see-\Pf{}//
\glft \trsl{He saw himself.} \lingcontext{reflexive; agent and patient coreferential}//
\endgl
\xe

\subsection{Interrogatives}

% [Port from 2020 archive — jede, ježe, jehát, jehu, jemí, jena, jak, zajehu, jiká, jišká, jeně, jení]

\subsection{Indefinite and negative pronouns}

% [Port from 2020 archive — nět/zide, niže/niho, etc.; note niho is superseded by zad- in the current analysis]

\section{Postpositions and case governance}\label{sec:prepositions}

Iridian is strictly head-final in its adpositional phrases. The governed noun phrase precedes the relational marker, and the marker itself is therefore best analyzed as a \term{postposed clitic adposition}. These items are phonologically weak, lean prosodically on the preceding noun phrase, and attach to the full governed phrase rather than to the noun stem alone. Orthographically they are here written as separate words for ease of reading, but morphologically they behave like enclitics.

All such clitic adpositions govern the \term{indirect case}. Older descriptions sometimes distributed these relations across several case labels or treated some of the markers as prepositions, but the synchronic generalization is simpler: the dependent noun phrase is inflected in the indirect, and the following clitic specifies the relation. The differences among the forms are therefore semantic, not differences in case government.

\subsection*{Morphological behavior}

The clitic attaches to the right edge of the governed noun phrase as a whole. A noun with modifiers, determiners, or possessors is therefore inflected first for the indirect case, after which the adposition follows as a phonological dependent. The syntax remains head-final throughout: the noun phrase comes first, and the clitic adposition closes the phrase.

\subsection*{Inventory}

The common clitic adpositions are \ird{na} \trsl{in, inside, on, at}, \ird{de} \trsl{into, to}, \ird{my} \trsl{for}, \ird{u} \trsl{near, by}, \ird{še} \trsl{with}, \ird{z} \trsl{from, out of}, \ird{nam} \trsl{without}, \ird{pale} \trsl{instead of}, and \ird{ahte} \trsl{except for}. Location, goal, proximity, benefit, accompaniment, source, exclusion, and replacement thus all belong to the same formal pattern: an indirect noun phrase followed by a phonologically weak clitic adposition.

\subsection*{Spatial and goal relations}

\pex
\a \irdp{dumě na}{in the house; at home}
\a \irdp{hradě de}{into the town; to the town}
\a \irdp{hradě u}{near the town}
\xe

The same pattern extends naturally to more specific spatial and path relations. In addition to the broad locative and goal clitics, the language also admits more specialized forms such as \ird{pod} \trsl{under, beneath}, \ird{nad} \trsl{over, above}, \ird{měz} \trsl{between, among}, \ird{do} \trsl{up to, as far as}, and \ird{mně} \trsl{through, via}. These too follow an indirect host noun phrase and behave as phonologically weak enclitics.

\pex
\a \irdp{stolě pod}{under the table}
\a \irdp{mostě nad}{over the bridge}
\a \irdp{domě měz}{between the houses; among the houses}
\a \irdp{hradě do}{up to the town}
\a \irdp{lesě mně}{through the forest; via the forest}
\xe

\subsection*{Participant and dependency relations}

\pex
\a \irdp{bylam my}{for the child}
\a \irdp{Jancí še}{with Janek}
\a \irdp{Marcí z}{from Marek}
\a \irdp{záhárí nam}{without sugar}
\a \irdp{Jancí pale}{instead of Janek}
\a \irdp{Jancí ahte}{except for Janek}
\xe

\subsection*{Clausal examples}

In full clauses the same structure is preserved: the host noun phrase is indirect, and the following adposition is glossed by its nearest English equivalent.

\pex
\a
\begingl
\gla Kroumaště na po zma jeca pivo.//
\glb refrigerator.\Ind{} in still few \Exst{}-\Cont{} beer.\Dir{}//
\glft \trsl{There's still some beer left in the refrigerator.}//
\endgl
\a
\begingl
\gla Bylam my Jancí še stóžách.//
\glb child-\Ind{} for Janek-\Ind{} with go-\Ctp{}//
\glft \trsl{(I am) going for the child with Janek.}//
\endgl
\a
\begingl
\gla Marcí z houba.//
\glb Marek-\Ind{} from gift.\Dir{}//
\glft \trsl{a gift from Marek}//
\endgl
\xe

The governing relation is thus constant across the system: the noun phrase takes the indirect, while the following clitic adposition contributes the relational meaning. This analysis captures both the head-final syntax of the phrase and the reduced phonological status of the adposition itself.

\section{Noun phrase structure}\label{sec:np-structure}

% ============================================================
% STATUS: STUB — to be developed
% The NP in Iridian is head-final: modifiers precede the head noun.
% Order within the NP:
%   ne (plural) > demonstrative > oblique possessor > stative-verb modifier > head
% Stative verbs as attributives: this is the primary adjectival strategy.
% The attributive suffix -ní can nominalize/adjectivalize a stative verb
% but the stative converb is also used.
% ============================================================

The canonical Iridian noun phrase is head-final. Modifiers of all kinds precede the head noun. The general ordering within the noun phrase is:

\medskip
\ird{ne} (plural) $>$ demonstrative $>$ oblique possessor $>$ property-root describer / stative-verb modifier $>$ head noun
\medskip

Possessors therefore occupy the same prenominal dependent slot described in Section~\ref{sec:indirect-uses}: they stand directly before the head noun and remain inside the noun phrase, while demonstratives and the plural clitic attach outside the possessive grouping.

Iridian has a very small closed set of lexical property roots used as simple
adnominal describers. These roots denote especially salient notions such as
\trsl{good}, \trsl{beautiful}, and \trsl{bad}, and they occur directly before
the head noun. Their use is lexically restricted and non-productive: speakers
do not freely coin new property-root describers for ordinary descriptive
content.

Outside this small residue, stative verbs (see Section~\ref{sec:stative} in
Chapter~\ref{ch:verbs}) serve as the primary attributive modification
strategy. Where English uses an adjective before a noun, Iridian normally uses
either a stative verb in an attributive/converbial form or a nominalized
stative verb. In many cases a property root and a stative verb form a lexical
pair: the root functions as an adnominal describer, while the corresponding
stative verb serves as the ordinary predicate.

% [Examples to be added]

\section{Event and participant nominals}\label{sec:event-nominals}

Iridian can form nouns directly from verbs. The resulting forms — called
\term{event nominals} and \term{participant nominals} — behave like
ordinary nouns in every grammatical respect, taking case endings, the
definiteness marker, and the plural clitic on exactly the same terms as
any other noun. What distinguishes them is their meaning: an event nominal
names the action or event itself (roughly what English does with
\trsl{the eating} or \trsl{the running}), while a participant nominal names
a person or thing defined by the role they play in that event (as English
does with \trsl{the eater} or \trsl{the food}).

Both types are formed by adding the suffix \ird{-ou} to a verb stem.
The details of which stem form to use in each case are treated in
Chapter~\ref{ch:verbs}; here the concern is with the resulting nouns
as noun phrases: how they inflect, how they interact with definiteness, and
how they encode the participants of the original event.

\subsection{Inflection and definiteness}

Nouns ending in \ird{-ou} follow the inflectional pattern of Class~II
nouns (see the section on inflectional classes above). In the direct case
the form ends in \ird{-ou}; the transitive ending is \ird{-ovy}; the
indirect is \ird{-ově}. The noun \ird{piaščikou} \trsl{the one who ate}
illustrates the full set:

\begin{table}[ht]
	\footnotesize\sffamily
	\caption{Inflection and definiteness for the participant nominal
	\ird{piaščikou}.}\label{tab:ou-inflection}
	\medskip
	\begin{tabularx}{0.72\textwidth}{L l l}
		\toprule
		{\sc case} & {\sc plain} & {\sc definite} \\
		\midrule
		Direct     & \ird{piaščikou}  & \ird{piaščikovat}  \\
		Transitive & \ird{piaščikovy} & \ird{piaščikovyt}  \\
		Indirect   & \ird{piaščikově} & \ird{piaščikovět}  \\
		\bottomrule
	\end{tabularx}
\end{table}

The definiteness marker (see Section~\ref{sec:definiteness}) attaches
after whichever case ending the noun already bears. The only irregularity
is in the direct case: the sequence \ird{-ou\,+\,-ot} does not surface as
\ird{*-out} but contracts to \ird{-ovat}. In the other cases the marker is
added without contraction: \ird{-ovy} yields \ird{-ovyt}, and \ird{-ově}
yields \ird{-ovět}.

\subsection{Two kinds of verb-derived noun}

Among nouns built on \ird{-ou}, it is useful to draw a line between two
kinds.

The first kind is the \term{lexical} or \term{resultant nominal}: a
noun tied to a specific verb that names the typical result, product, or
correlate of that verb's action. The verb for \trsl{to eat} can produce
\ird{piaštou} with the conventionalised meaning \trsl{food} — not any
particular act of eating, but the thing characteristically eaten or
produced by eating. Such forms must be learned verb by verb; many verbs
lack them entirely, relying instead on unrelated or inherited nouns.

The second kind is the \term{gerund}: the fully productive way of
turning any verb into a noun that names the event itself. Gerunds are
formed freely from any verb and always refer to the action as an abstract
event rather than to a fixed result. They correspond most closely to the
English nominalisation in \trsl{-ing} when used as a noun (\trsl{the
eating}, \trsl{the running}, \trsl{the remembering}).

Gerunds admit a further distinction based on what they contain. A
\term{simple} gerund names the event without specifying who was involved:
it means just \trsl{the eating}, nothing more. A \term{complex} gerund
names the event and also expresses the participants of that event inside the
noun phrase: \trsl{Marek's eating of the bread}, \trsl{Janek's
remembering}. No separate morphology marks the difference; the same
\ird{-ou} suffix produces both. What makes a given gerund simple or complex
is simply whether case-marked nouns appear inside it.

\subsection{Complex gerunds and their participants}

When a gerund expresses its participants internally, both the doer
(\textit{agent}) and the thing acted upon (\textit{theme}) appear in the
indirect case. This is the same case used throughout the noun phrase for
possession and dependent relations (see Section~\ref{sec:indirect-uses}).
The pattern is consistent with the rest of the grammar: just as
\ird{Marcí dum} means \trsl{Marek's house}, with the possessor Marek
in the indirect case, the doer of a gerundival event also appears in the
indirect case when expressed inside the noun phrase.

One consequence is that a complex gerund containing both a doer and a
theme will have two indirect-case nouns with no formal difference between
them. Consider:

\ex
\ird{Jancí podohletou} --- \trsl{the act of remembering Janek}
\textit{or} \trsl{Janek's act of remembering}
\xe

This is a genuine ambiguity. The indirect-case noun \ird{Jancí} (Janek)
could be either the person doing the remembering or the person being
remembered. Iridian does not resolve this ambiguity by grammatical rule.
A reader or listener is expected to use context, and this is by design:
complex event nominals regularly carry a degree of referential looseness
that a full finite clause would not. The grammar treats this as a natural
property of the construction rather than a defect.

When the speaker wishes to be explicit, two strategies are available.
The first is \term{word order}: the doer-indirect precedes the
theme-indirect. In a construction of the form \textit{A B gerund}, A is
the doer and B is the theme. The second is the \term{definiteness
marker}: a theme noun bearing the \ird{-ot/-t} suffix
(Section~\ref{sec:definiteness}) is read unambiguously as the theme and
not as the doer. Since the \ird{-ou} direct ending contracts with \ird{-ot}
to give \ird{-ovat}, this marker is visually distinctive when the theme
itself is a nominalized form.

\subsection{Participant nominals}

Participant nominals use the same \ird{-ou} suffix, but instead of naming
an event they name a \textit{participant} in that event, identified by the
role they occupy. The role a given noun names is determined by which verb
stem the suffix is attached to.

Iridian inflects verbs to foreground different participants as their
central argument (the system is described in Chapter~\ref{ch:verbs}).
Attaching \ird{-ou} to a given stem form yields a noun naming whichever
participant that form places at the centre. Three stem types are relevant
here:

\begin{itemize}

	\item The \term{antipassive stem} is the form of a transitive verb
	in which the doer of the action is treated as the central argument and
	the thing acted upon is expressed as a peripheral, indirect-case
	dependent. Adding \ird{-ou} to this stem produces an
	\term{agent-oriented} participant nominal meaning \trsl{the one who
	does X}. All five indicative aspect stems are available here. Two of
	them show special surface effects already noted in Chapter~\ref{ch:verbs},
	and two more contract regularly with \ird{-ou}:
	the perfective antipassive \ird{-ek} yields \ird{-ikou} (so
	\ird{piaščikou} means \trsl{the one who ate}), and the continuous
	antipassive \ird{-a} yields \ird{-ova} (so \ird{piaščova} means
	\trsl{the one who eats} or \trsl{the one who is eating}, whether
	habitually or at the present moment). The retrospective \ird{-aní}
	yields \ird{-anou}, and the progressive \ird{-ime} yields \ird{-imou};
	hence \ird{piaščanou} \trsl{the one who has eaten},
	\ird{piaščimou} \trsl{the one who is in the act of eating}, and
	\ird{piaščachou} \trsl{the one who is about to eat}.

	\item The \term{plain transitive stem} — the ordinary, unmarked form
	of a transitive verb — produces a \term{patient-oriented} participant
	nominal meaning \trsl{the thing that X is done to}. In Iridian's default
	treatment of transitive verbs (discussed in Chapter~\ref{ch:verbs}),
	it is the thing acted upon that occupies the central argument position
	of the construction. Attaching \ird{-ou} to this stem therefore names
	the recipient of the action. Here too all five aspects are available:
	\ird{piaštnikou} means \trsl{that which was eaten}, \ird{piaštnanou}
	\trsl{that which has been eaten}, \ird{piaštnova} \trsl{that which is
	habitually eaten}, \ird{piaštnimou} \trsl{that which is being eaten},
	and \ird{piaštnachou} \trsl{that which is about to be eaten}. More
	concretely, such forms may be interpreted as \trsl{food} understood as
	the participant affected by eating, with the chosen aspect supplying the
	expected temporal contour.

	\item The \term{applicative stem} is the form of a verb that promotes
	an indirect participant — typically a beneficiary, recipient, or other
	non-patient role — into the central argument position. Attaching \ird{-ou}
	to this stem produces a noun referring to that promoted participant.
	Thus \ird{piaštébikou}, built on the applicative stem of the verb for
	\trsl{to eat}, means \trsl{the one who was fed} or \trsl{the one for
	whom eating was done}.

\end{itemize}

The logic throughout is the same: whichever participant a given voice- and
aspect-marked verb form treats as its central argument, attaching \ird{-ou}
to that form yields a noun that names that participant. Antipassive-based
nominals name the doer; plain-transitive-based nominals name the thing
acted upon; applicative-based nominals name the participant promoted by the
applicative. Aspect contributes the expected interpretation of the event
profile: completed (\ird{piaščikou}), experientially or result-relevant
past (\ird{piaščanou}), ongoing or habitual (\ird{piaščova}),
currently in progress (\ird{piaščimou}), or intended / about to occur
(\ird{piaščachou}).

Participant nominals are particularly useful when the speaker wants to
identify someone by what they do or what is done to them, rather than by a
separate descriptive noun. They therefore tend to appear in constructions
that introduce or identify discourse participants, a topic taken up in
Chapter~\ref{ch:discourse}.

\subsection{The attributive form in \ird{-ovní}}

Participant nominals are Iridian's primary means of expressing what English
conveys with the relative pronouns \trsl{who}, \trsl{which}, and
\trsl{that}. The \ird{-ou} nominal stands alone as a freestanding noun; its
counterpart, the \term{attributive form} in \ird{-ovní}, is built from the
same verb stem and serves as a modifier directly preceding a head noun. The
construction that results is Iridian's standard equivalent of a restrictive
relative clause.

The two forms share a single organising principle: a nominalized verb phrase
denotes whichever participant is structurally \textit{privileged} within that
clause package. The suffix \ird{-ou} or \ird{-ovní} identifies the privileged
participant; every other core argument within the same nominalized domain
appears in the indirect case. The nominalized clause is its own enclosed
domain, so this rule operates within it independently of any surrounding
structure. When the noun phrase headed or modified by \ird{-ovní} is
embedded in a higher clause, it takes whatever case that outer clause
requires.

The paradigm below uses \ird{vrava} (Class~III, \trsl{letter}), \ird{pres}
(Class~II, \trsl{man}), and the verb \ird{obroh-} \trsl{to read}:

\pex
\a \ird{vravat presy obrohnik.} --- \trsl{The man read the letter.}
\a \ird{vravět obrohnikou} --- \trsl{the one who read the letter}
\a \ird{vravět obrohnikovní pres} --- \trsl{the man who read the letter}
\xe

In (\lastx b), the participant nominal picks out the reader as its referent.
The letter appears in the indirect case — \ird{vravě}, with the definiteness
suffix giving \ird{vravět} — as a non-privileged argument inside the
nominalized domain. In (\lastx c), the same stem takes the attributive suffix
\ird{-ovní} and immediately precedes the head noun \ird{pres}. The head noun
takes whatever case the surrounding clause requires; the \ird{-ovní} form
does not change.

\subsubsection*{Definiteness}

The attributive form \ird{-ovní} qualifies the head noun as something defined
by a verbal notion but does not itself encode definiteness. Definiteness,
where intended, is shown by the ordinary suffix \ird{-t}/\ird{-ot} on the
head noun, following case-inflection in the regular way; the \ird{-ovní}
modifier is unaffected.

\subsubsection*{Nested attributives}

Because each nominalized clause is its own domain, attributives may be
stacked recursively. The inner nominalized phrase becomes an argument within
the outer nominalized clause, appearing there in the indirect case as
required. The example below, introducing additionally \ird{zasd-} \trsl{to
write} and \ird{naz} \trsl{woman}, embeds one attributive inside another:

\pex
\a \ird{vravat naza zasdnik.} --- \trsl{The woman wrote the letter.}
\a \ird{nazam zasdnikovní vrava} --- \trsl{the letter that the woman wrote}
\a \ird{[nazam zasdnikovní vravět] obrohnikovní pres} ---
   \trsl{the man who read the letter that the woman wrote}
\xe

In (\lastx b), considered in isolation, the head noun \ird{vrava} appears in
the direct case: the letter is the privileged participant of the inner domain
(the thing written), and the woman appears as \ird{nazam} in the indirect.
In (\lastx c), this same inner phrase is embedded as a non-privileged
argument of the outer nominalized clause. Accordingly the inner head noun
shifts to the indirect: \ird{vrava} surfaces as \ird{vravět} (indirect with
definiteness suffix). The head noun of an \ird{-ovní} phrase always takes the
case required by the clause it is embedded in, not the case it would bear if
freestanding.

\subsubsection*{Practical limits on embedding}

The system is formally recursive, but in practice speakers rarely go beyond
two or three levels of embedding, and a single embedded attributive is by far
the most common. Deeper nesting produces noun phrases that are unwieldy to
produce and difficult to parse in speech. Speakers typically rephrase heavily
embedded content as a sequence of simpler clauses or a converb chain rather
than extending the nominalization further.


%% Chapter 5: The Simple Sentence
\chapter{The Simple Sentence}\label{ch:simple}

% ============================================================
% STATUS: SCAFFOLD — largely new content; some material available in
% archive/2020/chapters/07-clause-structure.tex.
% Key theoretical framing: topic-comment structure is the primary
% organizational principle of the Iridian clause. The split-ergative
% alignment can be recast as: the absolutive is typically the topic;
% the ergative demotes the agent away from topic status.
% ============================================================

\section{Topic-comment structure and information structure}\label{sec:topic-comment}

% ============================================================
% The basic Iridian clause is organized around a topic-comment structure.
% The topic (what the sentence is "about") is unmarked (absolutive) and
% typically appears sentence-initially. The comment (what is said about
% the topic) follows. The alignment system is closely tied to this:
%   - Perfective/retrospective: P is the topic (ABS), A is demoted (ERG)
%   - Nominative aspects: A/S is the topic (ABS), P is non-topic (OBL)
% This means aspect choice is partly pragmatic — the speaker chooses the
% aspect that places the intended topic in the unmarked (absolutive) position.
% Antipassive demotes P and restores A to topic/ABS position.
% ============================================================

The basic Iridian clause organizes information around a topic-comment structure. The \textbf{topic} --- the entity the sentence is about --- is expressed in the absolutive case and typically occupies the initial position in the clause. The \textbf{comment} --- what is predicated of the topic --- follows. This organization is reflected in and reinforced by the split-ergative alignment: in the perfective and retrospective aspects, the patient is promoted to the absolutive (becoming the topic), while in the nominative aspects, the agent occupies the absolutive (topic) position.

Aspect choice is therefore partly a pragmatic decision: the speaker selects the aspect that places the discourse-salient argument in the absolutive position.

% [Examples to be added]

\section{Word order}\label{sec:word-order}

The canonical word order in Iridian is SOV (subject-object-verb), where the finite verb is clause-final. This order is relatively rigid for the finite verb itself, which must appear last. The ordering of nominal arguments before the verb is more flexible and is governed by information structure rather than strict grammatical requirements. Topic-initial and focus-final tendencies apply within the pre-verbal domain.

% [Examples: SOV default, pragmatic reordering, focus movement]

\section{The copula and existential constructions}\label{sec:copula}

Iridian distributes predication across three structurally and semantically distinct constructions. Where many European languages rely on a single copular verb for all such functions, Iridian treats equative identity, marked stative or inchoative predication, and existential predication separately. All overt predicative verbs are clause-final and inflect for the full aspect system; each paradigm is strongly suppletive and morphologically irregular.

\subsection{Equative and class-membership clauses: the zero copula}\label{sec:zero-copula}

The default pattern for equative clauses --- those asserting identity, role, or class membership --- is zero copula in the indicative. The predicative complement appears in the direct case in clause-final position; no overt verb is present. This is the stylistically neutral, unmarked construction.

\pex
\a
\begingl
\gla Marek doktor.//
\glb Marek.\Dir{} doctor.\Dir{}//
\glft \trsl{Marek is a doctor.}//
\endgl
\a
\begingl
\gla Iridian malno.//
\glb Iridian.\Dir{} language.\Dir{}//
\glft \trsl{Iridian is a language.}//
\endgl
\xe

Inclusive and class-membership predications --- of the type \textit{dogs are animals} --- follow the same zero-copula pattern and carry no inference beyond simple predication. The predicative complement always appears in the direct case regardless of which case would otherwise be assigned in the same position.

\subsection{The marked copular verb \ird{neh-}: becoming and inchoative predication}\label{sec:neh}

The root \ird{neh-} provides the overt copular verb. Its paradigm is strongly suppletive and morphologically irregular; principal forms are given in Table~\ref{tab:neh-paradigm}.

\begin{table}[ht]
	\footnotesize\sffamily
	\caption{Principal forms of \ird{neh-} (marked copula).}\label{tab:neh-paradigm}
	\medskip
	\begin{tabularx}{0.72\textwidth}{L l}
		\toprule
		{\sc form} & {\sc iridian} \\
		\midrule
		Continuous          & \ird{neha} \\
		Progressive         & \ird{nehime} \\
		Contemplative       & \ird{nehach} \\
		Perfective          & \ird{nek} \\
		Retrospective       & \ird{neaní} \\
		Sequential converb  & \ird{nesu} (colloquial \ird{nu}) \\
		Simultaneous converb & \ird{nešěc} \\
		\bottomrule
	\end{tabularx}
\end{table}

Since neutral equative predication is expressed by zero copula (§\,\ref{sec:zero-copula}), overt \ird{neh-} is semantically marked. It is not used for simple statements of identity or class membership. Instead it signals \textbf{becoming}, \textbf{entering a state}, or \textbf{transition}: the subject is understood as moving from one classification or condition into another.

\pex
\a
\begingl
\gla Marek doktor neha.//
\glb Marek.\Dir{} doctor.\Dir{} \Cop{}-\Cont{}//
\glft \trsl{Marek is becoming a doctor.}//
\endgl
\a
\begingl
\gla Ša byl kapitan nek.//
\glb \Dem{}.\Prox{} child.\Dir{} captain.\Dir{} \Cop{}-\Pf{}//
\glft \trsl{This child became captain.}//
\endgl
\xe

In the continuous aspect, \ird{neh-} may describe a newly entered or recently changed state, with the inchoative flavour modulated by context; a purely stable state assertable without implication of transition calls for zero copula instead. The predicative complement of \ird{neh-} stands in the direct case, as with zero-copula equatives.

The sequential converb \ird{nesu} --- colloquially \ird{nu} --- is a suppletive form established independently of the regular aspect paradigm. Its grammaticalization as a discourse particle is described in Section~\ref{sec:seq-copula} of Chapter~\ref{ch:complex}.

\subsection{The existential verb \ird{jec-}: existence, presence, and having}\label{sec:jec}

The existential verb \ird{jec-} is a true clause-final verb asserting the existence or presence of a referent. It inflects for all aspects; its paradigm is strongly suppletive and morphologically irregular, with the perfective and retrospective built on a distinct suppletive root. Principal forms are given in Table~\ref{tab:jec-paradigm}.

\begin{table}[ht]
	\footnotesize\sffamily
	\caption{Principal forms of \ird{jec-} (existential verb).}\label{tab:jec-paradigm}
	\medskip
	\begin{tabularx}{0.72\textwidth}{L l}
		\toprule
		{\sc form} & {\sc iridian} \\
		\midrule
		Continuous          & \ird{jeca} \\
		Progressive         & \ird{jecime} \\
		Contemplative       & \ird{jecach} \\
		Perfective          & \ird{žek} \\
		Retrospective       & \ird{žaní} \\
		Sequential converb  & \ird{žu} \\
		\bottomrule
	\end{tabularx}
\end{table}

The division of labour between \ird{jec-} and the zero-copula construction is loosely parallel to the Japanese distinction between the existential verbs \textit{iru} and \textit{aru} on the one hand and the equative \textit{da}/\textit{desu} on the other: \ird{jec-} covers existence and presence; zero copula covers identity and class membership.

\paragraph{Existence and presence.} The primary use of \ird{jec-} is to assert that a referent exists or is present. The entity in question stands in the direct case.

\pex
\a
\begingl
\gla Byl jeca.//
\glb child.\Dir{} \Exst{}-\Cont{}//
\glft \trsl{There is a child. / A child is present.}//
\endgl
\a
\begingl
\gla Kroumaště na po zma jeca pivo.//
\glb refrigerator.\Ind{} in still few \Exst{}-\Cont{} beer.\Dir{}//
\glft \trsl{There is still some beer left in the refrigerator.}//
\endgl
\xe

\paragraph{Location.} \ird{Jec-} is also used when locating an entity: the entity appears in the direct case and the location is expressed by an indirect noun phrase followed by a postposed clitic adposition.

\pex
\a
\begingl
\gla Marek dumě na jeca.//
\glb Marek.\Dir{} house.\Ind{} in \Exst{}-\Cont{}//
\glft \trsl{Marek is at home.}//
\endgl
\xe

\paragraph{Having.} Possession is expressed through \ird{jec-} with an oblique possessor: `I have X' is construed as `there is X at me', with the possessor in the indirect case and the possessed item in the direct.

\pex
\a
\begingl
\gla Bylam jeca tóm.//
\glb child.\Ind{} \Exst{}-\Cont{} book.\Dir{}//
\glft \trsl{The child has a book. (Lit.\ At the child there is a book.)}//
\endgl
\xe

\subsection{The negative existential verb \ird{zad-}}\label{sec:zad}

The root \ird{zad-} is the lexically negative counterpart of \ird{jec-}: it asserts the non-existence or absence of a referent. Because \ird{zad-} is inherently negative, it does not combine with the standard negation prefix \ird{zá-}. Its paradigm is given in Table~\ref{tab:zad-paradigm}.

\begin{table}[ht]
	\footnotesize\sffamily
	\caption{Principal forms of \ird{zad-} (negative existential verb).}\label{tab:zad-paradigm}
	\medskip
	\begin{tabularx}{0.72\textwidth}{L l}
		\toprule
		{\sc form} & {\sc iridian} \\
		\midrule
		Continuous          & \ird{zada} \\
		Progressive         & \ird{zadime} \\
		Contemplative       & \ird{zadach} \\
		Perfective          & \ird{žatek} \\
		Retrospective       & \ird{žataní} \\
		Sequential converb  & \ird{zadu} \\
		\bottomrule
	\end{tabularx}
\end{table}

\pex
\a
\begingl
\gla Pivo zada.//
\glb beer.\Dir{} \NegExst{}-\Cont{}//
\glft \trsl{There is no beer.}//
\endgl
\a
\begingl
\gla Marek zada.//
\glb Marek.\Dir{} \NegExst{}-\Cont{}//
\glft \trsl{Marek is not here. / Marek is absent.}//
\endgl
\xe

The plural particle \ird{ne} cannot be combined with a \ird{zad-} clause; see Section~\ref{sec:definiteness} of Chapter~\ref{ch:nouns}.

\section{Questions}\label{sec:questions}

\subsection{Yes-no questions}

% [To be written: polar question formation — intonation, question particle]

\subsection{Content questions}

Content questions (\textit{wh}-questions) require the interrogative to occupy the topic (initial) position of the clause. \textit{Wh}-fronting is obligatory in Iridian; in-situ interrogatives are not grammatical in standard Iridian.

% [Port from 2020 archive §\,sec:content-questions; update case labels]

\subsection{Tag, echo, and indirect questions}

% [To be written]

\section{Negation in clausal context}\label{sec:clausal-negation}

At the clausal level, negation is expressed by the verbal prefix \ird{zá-} (reduced to \ird{za-} in unstressed environments). The prohibitive mood suffix \ird{-éma} handles negation of commands. The negative existential verb \ird{zad-} handles negation of existential constructions; because it is lexically negative, it does not take the prefix \ird{zá-} (see §\,\ref{sec:zad}). Interaction with the perfective aspect (negated perfective = `should have done but didn't') is discussed in Chapter~\ref{ch:verbs}.

% [Full treatment to be added]

\section{Comparative and superlative constructions}\label{sec:comparison}

% ============================================================
% The 2020 archive described a comparative construction using the agentive
% case (now: transitive case) for the standard of comparison, and a comparative
% particle. This needs to be updated to the new case terminology.
% ============================================================

% [Port and revise from 2020 archive; update "agentive of comparison" to
%  "transitive case in comparison"]


%% Chapter 6: Complex Sentences and Clause Chaining
\chapter{Complex Sentences and Clause Chaining}\label{ch:complex}

\section{Overview: clause-combining strategies}

Iridian employs several strategies for combining clauses. The dominant strategy is \textbf{converb chaining}, in which all non-final clauses are expressed as non-finite converb forms linked to a single main (finite) clause. Additional strategies include coordination with conjunctions, quotative constructions for reported speech, and \ird{še}-clauses for a limited set of complement and adjunct types.

The centrality of the converb to Iridian clause-combining means that what would be expressed as multiple finite clauses in a Slavic language typically appears in Iridian as a chain of converbs followed by a single finite verb. This has consequences for information structure (only the final, finite verb carries aspect and alignment), for switch-reference tracking, and for discourse organization (see Chapter~\ref{ch:discourse}).

\section{Converb chains: structure and alignment}\label{sec:converb-chains}

\subsection{The finite-verb-final rule}

In a converb chain, all verbs except the last are non-finite (converbs). Only the final verb is finite and carries aspect, mood, and alignment. This rule is absolute in standard written Iridian; in colloquial and dialectal speech, finite verbs may occasionally appear in non-final positions in certain coordinated structures.

\subsection{Accusative alignment within converb clauses}

Case marking within converb clauses always follows nominative-accusative alignment: the absolutive argument is the agent or sole argument, and the patient appears in the indirect (oblique) case. The ergative case is never found inside a converb clause, because converbs carry no aspect and the ergative alignment is triggered only by the perfective and retrospective aspects.

\subsection{Switch-reference and topic continuity}\label{sec:switch-reference}

By default, the subject of a converb clause is coreferent with the subject of the following clause (same-subject, SS reading). When the subjects differ (different-subject, DS), a pronominal clitic is attached to the converb to mark the switch. If both clauses have full nominal subjects, the full NPs disambiguate and no clitic is required.

\section{Sequential constructions}\label{sec:sequential}

\subsection{Forms: \ird{-u} and \ird{-ení}}

The sequential converb has two forms with subtly different temporal force. The \ird{-u} form implies a temporally bounded, completed sub-event preceding the main clause event. The \ird{-ení} form is temporally underspecified and is compatible with both sequential and simultaneous readings.

\pex
\a
\begingl
\gla Marek zdravení Janek znovime.//
\glb Marek sleep-\Seq{} Janek study-\Prog{}//
\glft \trsl{While Marek sleeps, Janek is studying.} \lingcontext{simultaneous reading with \ird{-ení}}//
\endgl
\xe

\subsection{The copula converb and the discourse particle \ird{nu}}\label{sec:seq-copula}

The marked copular verb \ird{neh-} (see §\,\ref{sec:neh} of Chapter~\ref{ch:simple}) has the sequential converb form \ird{nesu}, a suppletive form established independently of the regular aspect paradigm. In colloquial speech, \ird{nesu} is contracted to \ird{nu}. The form \ird{nu} has been grammaticalized in some dialects as a discourse particle inserted between clauses in the manner of English `and then', without strong sequential implication and without retaining the inchoative/copular force of the underlying verb.

\subsection{Interaction with applicative and causative}

\pex
\a \ird{piašt-u} / \ird{piašt-ení} \trsl{having eaten / and [then] ate\ldots}
\a \ird{piašt-éb-u} / \ird{piašt-éb-ní} \trsl{having fed someone}
\a \ird{piašt-nau} / \ird{piašt-na-ní} \trsl{having made someone eat}
\xe

\section{Temporal constructions}\label{sec:temporal}

\subsection{The simultaneous / manner converb (\ird{-ěc})}

The simultaneous converb (\ird{-ěc}) describes an action that is ongoing or co-occurring with the main clause event, and also functions as a manner converb. The suffix belongs to the palatalizing domain of front-vocalic morphology (see Section~\ref{sec:front-palatalization}), not to the special sibilant-replacement system of the antipassive. When the converb and main clause have different subjects, the sequential \ird{-ení} is preferred for purely simultaneous meaning.

\subsection{The six temporal converbs}

\begin{table}[ht]
	\footnotesize\sffamily
	\caption{Temporal converb suffixes.}\label{tab:tempconverbs}
	\medskip
	\begin{tabularx}{0.9\textwidth}{L l L}
		\toprule
		{\sc suffix} & {\sc meaning} & {\sc function} \\
		\midrule
		\ird{-ezak}    & until              & main clause holds until converb-clause event \\
		\ird{-eškady}  & when, around when  & loose temporal co-occurrence \\
		\ird{-eeby}    & since              & from the time of the converb-clause event \\
		\ird{-ešhoume} & as soon as         & immediately upon the converb-clause event \\
		\ird{-eebymy}  & after              & after the converb-clause event \\
		\ird{-esny}    & before             & before the converb-clause event \\
		\bottomrule
	\end{tabularx}
\end{table}

\section{Causal constructions}\label{sec:causal}

% ============================================================
% STATUS: STUB — causal converb suffix not yet documented.
% The 2026 notes list converb types 1, 2, 4, 5 and skip 3 (causal).
% This section is a placeholder until the suffix is confirmed.
% ============================================================

A causal converb introduces the reason or cause for the main clause event (`because X, Y'). The suffix is not yet confirmed in the current revision and will be documented here once established.

\section{Concessive constructions}\label{sec:concessive}

The concessive converb (\ird{-ouc}) introduces a conceded or contrasted clause (`although X, Y'). It licenses the irrealis mood in the main clause when the conceded content is contrary to expectations.

\section{Conditional constructions}\label{sec:conditional}

\subsection{The four conditional converbs}

\begin{table}[ht]
	\footnotesize\sffamily
	\caption{Conditional converb suffixes.}\label{tab:conditional}
	\medskip
	\begin{tabularx}{0.8\textwidth}{L l L}
		\toprule
		{\sc suffix} & {\sc meaning} & {\sc notes} \\
		\midrule
		\ird{-ic}   & if (possibility)       & does not trigger palatalization \\
		\ird{-emy}  & if not (possibility)   & \\
		\ird{-ezmy} & if/when (expectation)  & \\
		\ird{-ena}  & if/when not (expected) & \\
		\bottomrule
	\end{tabularx}
\end{table}

\subsection{No palatalization with \ird{-ic}}\label{sec:cond-nopala}

Despite beginning with the front vowel \ird{-i-}, the conditional converb \ird{-ic} does not trigger palatalization of the thematic consonant. This is a historical accident: the suffix was grammaticalized during a period when this \ird{-i-} lacked palatalizing force. The conditional is therefore distinguishable from the irrealis (\ird{-ěv}), the simultaneous converb (\ird{-ěc}), and the conjunctive linker \ird{-e}, all of which belong to the productive palatalizing domain.

\section{Purposive constructions}\label{sec:purposive}

The purposive converb (\ird{-it}) expresses purpose (`in order to X') and serves as the complement of modal and desiderative verbs. Palatalization is irregular, applying mainly to Class~I stems.

\subsection{The intentional future reanalysis}

In colloquial speech, the purposive converb has been reanalyzed as a future/intentional marker with first-person clitics when no governing verb is present:

\ex
\begingl
\gla Dá trave piašcitem.//
\glb \textsc{1sg} bread.\Ind{} eat-\Purp{}-\textsc{1sg}//
\glft \trsl{I'm gonna eat bread.} \lingcontext{colloquial intentional future}//
\endgl
\xe

\section{Contrastive constructions}\label{sec:contrastive}

The contrastive converb (\ird{-ema}) introduces a contrasting clause (`X, but/whereas Y'), covering simple contrast, adversative contrast, and parallel contrast.

\section{Coordination}\label{sec:coordination}

% ============================================================
% STATUS: STUB — systematic treatment needed.
% Key conjunctions: a (additive), še (additive / complementizer),
% má / ozná (adversative), disjunctives.
% ============================================================

Coordination is expressed by \ird{a} (additive), \ird{še} (additive, also complementizer), \ird{má} / \ird{ozná} (adversative), and a set of disjunctives.

\section{Quotative constructions and reported speech}\label{sec:quotative}

% ============================================================
% STATUS: STUB — quotative complementizer to-že, direct quotatives,
% evidential/quotative verbal suffix. Port from 2020 archive.
% ============================================================

Reported speech uses the quotative complementizer \ird{to-že} for indirect quotation and allows fronted direct quotations before a \textit{verbum dicendi}. Clause-linking and certain quotative constructions further employ the verbal linker \ird{-e}, whose morphophonemics are described in Section~\ref{sec:conjunctive-linker} of Chapter~\ref{ch:verbs}. The evidential/quotative morpheme on the verb signals attribution to another speaker.

\section{Epistemic particles and evidential extensions}\label{sec:evidential-particles}

% ============================================================
% STATUS: STUB — see open question in ROADMAP.md §3.
% ============================================================

\section{Fixed expressions with converbs}\label{sec:fixed-converbs}

% ============================================================
% Thanks, apologies, congratulations, greetings that are historically
% grammaticalized converb forms — to be documented.
% ============================================================


%% Chapter 7: Discourse
\chapter{Discourse}\label{ch:discourse}

% ============================================================
% STATUS: SCAFFOLD — largely new; no direct predecessor in 2020 archive.
% Some material in archive/2020/chapters/08-semantics-and-usage.tex may be relevant.
% This chapter is theoretically the most ambitious of the rewrite:
% it reframes the alignment system as a discourse-level phenomenon and
% treats the converb chain as the basic unit of narrative structure.
% ============================================================

\section{Sentence-chains as discourse units}\label{sec:discourse-chains}

% ============================================================
% Iridian discourse is organized in "sentence chains": sequences of
% converb clauses capped by a finite verb. Each sentence chain functions
% as a single discourse unit. The chain has:
%   - Converb clauses = background, supporting, or subordinate events
%   - Finite verb = the foregrounded, topically relevant event
% This parallels the foreground/background distinction found in many
% verb-final or SOV languages with converb systems.
% ============================================================

\section{Topic continuity and topic shift}\label{sec:topic-continuity}

% The absolutive participant tends to persist as topic across sentence
% chains. When the topic shifts, this is typically signaled by:
%   - Lexical re-introduction of the new topic
%   - Switch-reference clitic on the converb
%   - A demonstrative or proximate/obviative distinction
% The obviative/proximate system of demonstratives (ša vs. dní) is the
% primary cross-clause topic-tracking device.

\section{Definiteness in discourse context}\label{sec:discourse-definiteness}

% The definiteness marker -ot/-t on a noun signals that the referent is
% identifiable to the hearer. In discourse, first mention is typically
% bare (no definiteness marker); subsequent mention may use the marker,
% the demonstrative, or drop the noun entirely (pro-drop).
% The interaction of definiteness, topic, and absolutive alignment
% is the key thread here.

\section{Aspect choice as narrative strategy}\label{sec:aspect-narrative}

% The split-ergative alignment means that aspect choice simultaneously
% encodes:
%   (a) the temporal/aspectual shape of the event
%   (b) which participant is the discourse topic
% In narrative, speakers use the perfective/retrospective to bring the
% patient into topic/absolutive position (foregrounding the event outcome)
% and use nominative aspects to keep the agent as topic (foregrounding
% the agent's activity).
% This is the core pragmatic principle behind the aspect system.

In extended discourse, aspect choice is inseparable from participant tracking. The perfective and retrospective place the patient in the absolutive and are therefore preferred when the discourse is about what happened \emph{to} a participant: they are the natural vehicle for narrative foreground, event culmination, and topic shift to an affected referent. The nominative aspects, by contrast, preserve the agent or sole participant as absolutive and are accordingly favoured where the continuity of an actor matters more than the bounded result of the event.

The same principle explains the discourse value of the antipassive. Speakers choose the antipassive when a transitive event is relevant chiefly as part of an agent's ongoing activity and not because the patient itself needs to be established, identified, or tracked. The construction keeps the agent in absolutive position, suppresses any patient from the grammatical core, and thereby functions as a topic-maintenance device. It is especially natural when the patient is generic, incorporated into the event type, inferable from context, or simply of no further narrative consequence; for the morphological details see Chapter~\ref{ch:verbs}.

The causative serves a different discourse purpose. By adding a causer, it allows the speaker to recast an event as one of inducement, delegated action, permission, or coercion. This is a useful viewpoint-management device in narrative: a simple predicate presents an event as directly performed, whereas a causative can foreground authority, distribute responsibility across two participants, or distance the discourse topic from the event's direct execution. In perfective contexts this often yields a strongly interventionist reading; in retrospective contexts it more readily supports permissive, indirect, or weak-control interpretations. The formal argument structure of the causative is described in Chapter~\ref{ch:verbs}; what matters here is that the construction lets the speaker decide whether the discourse should centre on the doing, the causing, or the managed consequences of the action.

\section{Tail-head linkage}\label{sec:tail-head}

% Tail-head linkage is a discourse-structuring device where the final
% clause of one sentence chain is repeated (in whole or in abbreviated
% form) at the start of the next chain. This is attested in the sample
% texts and is a common device in narrative Iridian.

\section{Discourse particles}\label{sec:discourse-particles}

% ============================================================
% STATUS: STUB — to be documented
% Categories:
%   - Topic-management particles (topic shift, topic continuation)
%   - Stance particles (mirative, dubitative, concessive)
%   - Connectives (now, then, well, anyway)
% Some of these were described as "evidential" in the 2020 archive
% (under the QUOT/Infer/Rep labels). They should be reconciled with
% the evidential slot in the morphological template.
% ============================================================

\section{Register variation}\label{sec:register-discourse}

% Literary vs. colloquial clause-linking preferences:
%   - Literary: tends toward full converb chains, explicit nominalization
%   - Colloquial: nu as discourse particle, purposive-future, shorter chains
%   - The increasing use of finite coordinated clauses under Slavic influence


%% Chapter 8: Derivational Morphology
\chapter{Derivational Morphology}\label{ch:derivational}

% ============================================================
% STATUS: SCAFFOLD — port and revise from
% archive/2020/chapters/06-derivational-morphology.tex.
% Key review points:
%   - Nominal derivation: update case terminology in examples
%   - Agent nominals: relation to the new antipassive stem (-kou)
%   - Abstract nominals: interaction with the new case system
%   - Verbal derivation: how does it interact with the antipassive /
%     applicative / causative slot? Are there derivational prefixes that
%     interact with valency?
%   - Borrowing: update to include recent borrowings and assimilation
%     patterns under the new stem-class system (the German -irná verbs)
% ============================================================

\section{Nominal derivation}

\subsection{Diminutives and augmentatives}

% [Port from 2020 archive]

\subsection{Agent nouns: \ird{-(h)ár} and \ird{-ist}}

Agent nouns denoting the habitual or professional performer of an action are formed with the suffixes \ird{-(h)ár} and \ird{-ist} (the latter a Slavic borrowing, now productive):

% [Examples: fisherman, teacher, activist, etc.]
% Note: the agent nominal -kou (from the antipassive stem) is a verbal
% derivation documented in §\,\ref{sec:nominalized}; the -ár/-ist forms
% are purely nominal derivation.

\subsection{Place nouns: \ird{-(i)mašt}}

% [Port from 2020 archive]

\subsection{Abstract nouns: \ird{-(i)žnám}}

% [Port from 2020 archive]

\section{Verbal derivation}

% ============================================================
% STATUS: STUB — the interaction of valency operations (antipassive,
% applicative, causative) with derivational morphology needs to be
% clarified. The question is whether some "derivational" operations
% in the 2020 archive are now better analyzed as instances of the
% applicative or causative, or whether they are genuinely separate.
% ============================================================

\section{Compounding}

% [Port from 2020 archive; update examples]

\section{Borrowing and assimilation patterns}

% The 2020 archive described:
%   - Phonological assimilation of loanwords (stress, consonant substitution)
%   - German loanwords in -irná (treated as a separate verbal class:
%     thematic consonant = -r, active root in -irž-)
%   - Recent borrowings (English, international) showing less assimilation
% These patterns remain valid in the 2026 framework; update terminology only.



%% Chapter 9: Semantics, Pragmatics, and Usage
\chapter{Semantics, Pragmatics, and Usage}\label{ch:semantics}

% ============================================================
% STATUS: SCAFFOLD — port and revise from
% archive/2020/chapters/08-semantics-and-usage.tex.
% ============================================================

\section{Register and formality}\label{sec:register}

% [Port from 2020 archive; update case terminology in examples]

\section{Politeness and the T-V distinction}\label{sec:politeness}

% ============================================================
% The 2020 archive had a detailed treatment of the T-V distinction:
%   - 2pl used as polite singular
%   - 1pl used for speaker humility in formal contexts
%   - Pro-drop as a politeness strategy (avoiding pronouns entirely)
%   - Use of names and titles as quasi-honorifics
% The plural marker ne/ně as an honorific with proper names and
% kinship terms (ne Stám Kovárž, ne mlazka) should be discussed here.
% ============================================================

\section{Honorifics}\label{sec:honorifics}

% [Develop from 2020 archive notes on ně/ne as honorific particle]

\section{Phatic expressions}\label{sec:phatic}

% [Greetings, leave-takings, thanks, apologies — cross-reference with
%  §\,\ref{sec:fixed-converbs} in Chapter~\ref{ch:complex}]

\section{Idiomatic expressions}\label{sec:idioms}

% [Selected idioms from the 2020 archive, updated glosses]



\end{document}
